\section{Notation}\label{app:notation}

Table~\ref{tab:notation} summarizes the notation used throughout the paper.

\begin{table}[h]
\centering
\small
\begin{tabular}{cl}
\toprule
Symbol & Meaning \\
\midrule
$\mathcal{M}$ & Language model \\
$x$ & Input instance \\
$y_i$ & $i$-th stochastic sample \\
$N$ & Number of samples \\
$T$ & Sampling temperature \\
$c(\cdot)$ & Confidence signal \\
SJ & Soft Jaccard similarity \\
FK & Fleiss' $\kappa$ \\
EM & Exact-match voting rate \\
VC ($c_{\mathrm{vc}}$) & Entity voting confidence \\
LP ($c_{\mathrm{lp}}$) & Log-probability (mean token log-likelihood) \\
\bottomrule
\end{tabular}
\caption{Notation used throughout the paper.}
\label{tab:notation}
\end{table}

\section{Robustness Analysis}\label{app:robustness}

We exclude instances with empty gold annotations from all correlation analyses (CoNLL: 697/3{,}453 instances, 20.2\%; SciERC: 22/551, 4.0\%).
For these instances, F1 is trivially zero regardless of prediction quality, and consistency signals reflect agreement on the empty set rather than extraction quality. Including them inflates or deflates Spearman $\rho$ depending on signal behavior at the boundary.

\begin{table}[t]
\centering
\small
\caption{Signal robustness across configurations (Spearman $\rho$). Top: temperature on Qwen SciERC. Middle: sample size on Qwen SciERC. Bottom: transfer across model families (both 8B) (LLaMA CoNLL $N{=}16$, 3-seed mean $\pm\sigma$). $^\dagger$Mean over 3 seeds ($\sigma{=}0.013$; individual: 0.398, 0.423, 0.410). $\sigma$ computed as sample standard deviation.}
\label{tab:robustness}
\begin{tabular}{llcc}
\toprule
& & NER $\rho$ & RE $\rho$ \\
\midrule
\multicolumn{4}{l}{\textit{Temperature} (Qwen SciERC, $N{=}8$)} \\
\quad SJ & $T{=}0.7$ & 0.364 & 0.234 \\
\quad SJ & $T{=}1.0$ & 0.360 & 0.250 \\
\quad FK & $T{=}0.7$ & 0.289 & 0.184 \\
\quad FK & $T{=}1.0$ & 0.267 & 0.275 \\
\quad LP & $T{=}0.7$ & 0.249 & --- \\
\quad LP & $T{=}1.0$ & 0.205 & 0.266 \\
\midrule
\multicolumn{4}{l}{\textit{Sample size} (Qwen SciERC, $T{=}1.0$)} \\
\quad SJ & $N{=}8$  & 0.360 & 0.250 \\
\quad SJ & $N{=}16$ & 0.411$^\dagger$ & 0.281 \\
\quad FK & $N{=}8$  & 0.267 & 0.275 \\
\quad FK & $N{=}16$ & 0.283 & 0.279 \\
\midrule
\multicolumn{4}{l}{\textit{Across model families} (LLaMA CoNLL, $N{=}16$, 3 seeds)} \\
\quad SJ &  & $0.474{\pm}0.009$ & --- \\
\quad FK &  & $0.467{\pm}0.009$ & --- \\
\quad EM &  & $0.473{\pm}0.010$ & --- \\
\quad VC &  & $0.470{\pm}0.006$ & --- \\
\quad LP &  & $0.311{\pm}0.003$ & --- \\
\bottomrule
\end{tabular}
\end{table}

\begin{figure}[t]
\centering
\includegraphics[width=\columnwidth]{fig_temperature_robustness}
\caption{Temperature robustness on SciERC ($N{=}8$). Top: NER; bottom: RE. SJ maintains stable $\rho$ across $T{=}0.7$ and $T{=}1.0$, while FK shows mixed sensitivity. LP shows the largest NER temperature effect ($\Delta\rho{=}{+}0.044$); RE $T{=}0.7$ LP data unavailable.}
\label{fig:temperature_robustness}
\end{figure}

\paragraph{Temperature sensitivity.}
SJ NER $\rho$ shifts from 0.364 ($T{=}0.7$) to 0.360 ($T{=}1.0$), a difference of 0.004.
FK NER $\rho$ drops by 0.022, while RE $\rho$ increases by 0.091, likely reflecting reduced diversity at lower temperature interacting differently across tasks.
LP exhibits the largest temperature effect on NER ($\rho{=}0.249$ at $T{=}0.7$ vs.\ $0.205$ at $T{=}1.0$, $\Delta{=}{+}0.044$), consistent with lower temperature concentrating probability mass and improving LP's discriminative power for quality estimation; RE $T{=}0.7$ LP data is not available.
Across all comparisons, the signal hierarchy is preserved at both temperatures on NER (SJ${>}$FK${>}$LP at $T{=}0.7$: $0.364{>}0.289{>}0.249$; at $T{=}1.0$: $0.360{>}0.267{>}0.205$).

\paragraph{Sample size scaling.}
Increasing from $N{=}8$ to $N{=}16$ improves SJ on Qwen SciERC NER: $\rho$ rises from $0.356{\pm}0.014$ (5 seeds) to $0.411{\pm}0.010$ (three seeds: 0.398, 0.423, 0.410), a gain of $+0.055$.
The magnitude of $N$-scaling gains is setting-dependent (LLaMA SciERC NER yields $+0.016$ and Qwen CoNLL NER $+0.014$), though the direction is consistently positive.
LLaMA CoNLL exhibits the most uniform scaling: all four consistency signals gain $\Delta\rho{\approx}{+}0.042$ (range: $+0.039$ to $+0.044$) while LP gains only $+0.004$ (Table~\ref{tab:quality_estimation}).
FK $N$-scaling is seed-dependent: Qwen SciERC gains range from $+0.005$ to $+0.023$ across seeds (mean $+0.016$), while LLaMA CoNLL yields $+0.039$ (3-seed mean); individual seed FK shifts vary substantially under the unified pipeline, so this should not be interpreted as a robust trend.
VC $N$-scaling is consistently positive ($+0.005$ on Qwen CoNLL, $+0.018$ on LLaMA SciERC, $+0.030$ on Qwen SciERC), while LP is approximately $N$-invariant ($\Delta{\leq}0.034$).
SJ, EM, and VC scale consistently across all tested settings; only FK shows seed-dependent variation that does not constitute a reliable trend, structural and binary-agreement signals are the most robust to $N$.
At $N{=}16$, SJ converges to comparable NER performance with VC ($\Delta{=}0.002$) while exhibiting lower cross-seed variance ($\sigma{=}0.010$ vs.\ $0.026$), a more stable alternative at higher sample sizes.

\paragraph{Transfer across model families.}
LLaMA-3.1-8B-Instruct shows a similar SJ--FK ordering on SciERC NER (SJ $\rho{=}0.325{\pm}0.021$, FK $\rho{=}0.256{\pm}0.024$, $\Delta\rho{=}0.069$; permutation $p{=}0.06$, borderline).
On CoNLL, the gap across model families is minimal (Table~\ref{tab:quality_estimation}): SJ achieves $\rho{=}0.431$ (LLaMA) vs.\ 0.436 (Qwen), $\Delta{=}0.005$.
At $N{=}16$, the structural-over-LP advantage holds across all four tested NER settings: Qwen SciERC (VC $\rho{=}0.409$ vs.\ LP 0.239, $\Delta{=}{+}0.170$), Qwen CoNLL (EM $\rho{=}0.463$ vs.\ LP 0.328, $\Delta{=}{+}0.135$),\footnote{Qwen CoNLL $N{=}16$ was not recomputed under the unified pipeline; LP $\rho{=}0.328$ may differ slightly from the registry value (0.315) due to instance filtering.} LLaMA SciERC (VC $\rho{=}0.452$ vs.\ LP 0.277, $\Delta{=}{+}0.175$), and LLaMA CoNLL 3-seed (SJ $\rho{=}0.474{\pm}0.009$ vs.\ LP $0.311{\pm}0.003$, $\Delta{=}{+}0.163$).

\paragraph{FK independence: instance-level permutation test.}\label{app:fk_independence}

To verify that Fleiss' $\kappa$ reflects genuine within-instance agreement rather than marginal-rate artifacts, we conduct instance-level permutation tests across three configurations.
For each instance, we permute sample labels 10{,}000 times to generate a null distribution of FK under independence, then compare the observed FK against this null via a one-sample $t$-test and permutation $p$-value.

\begin{table}[t]
\centering
\small
\caption{Instance-level FK independence test. Permutation tests confirm FK reflects true within-instance agreement (all $p{<}0.001$). Split-half reliability $r$ computed by randomly partitioning $N$ samples into two halves and correlating their FK scores.}
\label{tab:fk_independence}
\resizebox{\columnwidth}{!}{
\begin{tabular}{lrcccc}
\toprule
Config & $n$ & FK mean & 95\% CI & $t$ & Split-half $r$ \\
\midrule
SciERC (Qwen) & 529 & 0.402 & [0.376, 0.428] & 30.8 & 0.395 \\
CoNLL (LLaMA) & 2{,}756 & 0.927 & [0.919, 0.935] & 225.9 & 0.578 \\
CoNLL (Qwen) & 2{,}756 & 0.720 & [0.706, 0.734] & 102.1 & 0.520 \\
\bottomrule
\end{tabular}
}
\end{table}

Table~\ref{tab:fk_independence} reports the results.
All configurations reject the independence null at $p{<}0.001$ ($n{=}529$--$2{,}756$), FK captures genuine inter-sample agreement rather than chance label co-occurrence.
Split-half reliability ranges from $r{=}0.395$ (SciERC) to $r{=}0.578$ (CoNLL LLaMA), moderate stability: FK estimates from half the samples correlate positively with those from the other half, though the SciERC value suggests noisier estimation on harder tasks with fewer samples.
The higher FK means on CoNLL ($0.720$--$0.927$) vs.\ SciERC ($0.402$) mirror the difficulty gradient observed in the main quality estimation results.

\section{Calibration Analysis}\label{app:calibration}

\begin{figure*}[t]
\centering
\includegraphics[width=\textwidth]{fig_reliability}
\caption{Reliability diagrams on SciERC ($N{=}8$, $T{=}1.0$). Top: mean F1 vs.\ binned confidence; bottom: instance counts per bin. Left: NER; right: RE. NER signals are reasonably well-calibrated (calibration ECE 0.069--0.096), while RE shows greater miscalibration (ECE 0.109--0.339), RE calibration remains challenging for compound extraction tasks.}
\label{fig:reliability}
\end{figure*}

Reliability diagrams use 10 equal-width bins.
NER signals are reasonably well-calibrated (calibration ECE 0.069--0.096), while RE exhibits greater miscalibration (ECE 0.109--0.339), translating consistency scores into calibrated confidence estimates is feasible for NER but remains challenging for compound extraction tasks.

\section{Additional Analyses}\label{app:additional}

\begin{figure*}[t]
\centering
\includegraphics[width=\textwidth]{fig_signal_heatmap}
\caption{Quality estimation heatmap (Spearman $\rho$) across five signals and four configurations. Left: full evaluation set; right: conditional set (non-zero F1 instances). On the full set, structural signals dominate on SciERC NER while all signals converge on CoNLL. On the conditional set, EM emerges as the strongest signal on both NER ($\rho{=}0.360$) and RE ($\rho{=}0.417$), while LP RE collapses to near-zero ($\rho{=}0.011$).}
\label{fig:signal_heatmap}
\end{figure*}

\paragraph{Per-relation diagnostic.}
RE signal degradation is relation-specific.
\textsc{Feature-Of} has near-zero correlation ($\rho{=}0.004$), while \textsc{Hyponym-Of} achieves $\rho{=}0.646$.
The primary bottleneck is entity span errors, which account for 55.5\% of RE false negatives. This is a prerequisite failure that consistency signals cannot detect, as they measure agreement among equally incorrect extractions.
We note that RE analysis relies on a single dataset (SciERC); generalization to other RE benchmarks remains to be validated.

\paragraph{Capability stratification.}
For NER, SJ effectiveness scales with structural complexity: on SciERC instances with 6+ entities, $\rho{=}0.547$, compared to $\rho{=}0.415$ for 1--2 entities.
The SJ advantage over LP remains stable across sentence-length quartiles ($\Delta\rho{\approx}0.08$--$0.20$).
For RE, SJ $\rho$ increases monotonically with sentence length (Q1: 0.110 $\to$ Q4: 0.473; Fisher $z$-test $p{=}0.004$): longer, more relation-dense sentences provide richer structure for consistency-based estimation.
Both proxies are weakly correlated with F1 ($|r|{<}0.15$), they capture input complexity rather than directly predicting output quality; however, entity count and sentence length are themselves moderately correlated ($\rho{=}0.650$), so they reflect overlapping rather than independent complexity dimensions.
For RE, entity count and relation count are strongly correlated ($\rho{=}0.769$) and thus reflect a single structural complexity dimension rather than independent axes.
On the high-complexity RE subset ($n{=}81$, high relation count), signals improve substantially: VC reaches $\rho{=}0.630$, LP $\rho{=}0.554$, and SJ $\rho{=}0.502$, all far above the full-set levels.
Richer relational structure provides more signal for consistency-based estimation, complementing the task-level boundary established by WNUT-17: structural signals are most effective within a ``sweet spot'' of sufficient model capability and sufficient output complexity.

\paragraph{Length vs.\ log-probability.}
Figure~\ref{fig:length_logprob_scatter} shows that mean log-probability is weakly positively correlated with output length (NER: $\rho{=}0.386$, $n{=}4{,}232$; RE: $\rho{=}0.326$, $n{=}3{,}176$), ruling out length-penalty confounding as an explanation for LP's selection failure: longer outputs receive \textit{higher}, not lower, token-averaged log-probability under constrained decoding.

\begin{figure}[t]
\centering
\includegraphics[width=\columnwidth]{fig_length_logprob_scatter}
\caption{Output length vs.\ mean log-probability on SciERC ($N{=}8$, $T{=}1.0$). Positive correlations (NER: $\rho{=}0.386$; RE: $\rho{=}0.326$) rule out length-penalty confounding.}
\label{fig:length_logprob_scatter}
\end{figure}

\paragraph{Evaluation circularity and LOO-SJ equivalence.}
Consensus-proximity selection chooses $\hat{y} = \arg\max_{y_i} \sum_{j \neq i} \text{SJ}(y_i, y_j)$.
Leave-one-out SJ (LOO-SJ) for sample $y_i$ is $\text{LOO-SJ}_i = \frac{1}{N-1}\sum_{j \neq i} \text{SJ}(y_i, y_j)$.
Since the $\frac{1}{N-1}$ constant is shared across all candidates within an instance, $\arg\max_{y_i} \text{LOO-SJ}_i = \arg\max_{y_i} \sum_{j \neq i} \text{SJ}(y_i, y_j)$; the two selection rules are equivalent.
Standard (non-LOO) SJ adds the self-comparison term $\text{SJ}(y_i, y_i) = 1$ to each candidate, which is constant and thus does not affect ranking.
Empirically, LOO-SJ yields $\Delta\rho{=}{-}0.014$ relative to standard SJ (a correlation-level difference), with a maximum per-instance score difference of $3.1{\times}10^{-4}$ (negligible at the individual prediction level).
Evaluation circularity is negligible.

\paragraph{SJ normalizer robustness.}
Equation~\ref{eq:sj} uses $\max(|g_i|, |g_j|)$ as the denominator; an alternative is the union size $|g_i \cup g_j|$.
Table~\ref{tab:sj_normalizer} compares the two normalizers across all six model--dataset combinations ($N{=}8$, $T{=}1.0$, seed${=}42$): $|\Delta\rho|{<}0.01$ in every case, the choice between max and union normalization does not affect conclusions.

\begin{table}[h]
\centering
\small
\caption{SJ normalizer comparison (Spearman $\rho$; $N{=}8$, $T{=}1.0$, seed${=}42$). Max: $\max(|g_i|,|g_j|)$; Union: $|g_i \cup g_j|$.}
\label{tab:sj_normalizer}
\begin{tabular}{@{}ll ccc@{}}
\toprule
Dataset & Model & $\rho$(max) & $\rho$(union) & $\Delta\rho$ \\
\midrule
CoNLL & Qwen3-8B & .436 & .438 & $-$.002 \\
CoNLL & Qwen3-4B & .422 & .421 & $+$.000 \\
CoNLL & LLaMA-8B & .431 & .432 & $-$.000 \\
SciERC & Qwen3-8B & .360 & .365 & $-$.005 \\
SciERC & Qwen3-4B & .422 & .423 & $-$.001 \\
SciERC & LLaMA-8B & .363 & .372 & $-$.009 \\
\bottomrule
\end{tabular}
\end{table}

\paragraph{Conditional filtering details.}
On LLaMA conditional RE (non-zero F1 instances), SJ recovers to $\rho{=}0.405$ while FK drops to $\rho{=}0.259$, consistent with the NER structural advantage; on full-set RE, however, the SJ--FK difference is non-significant ($\Delta\rho{=}{-}0.005$, $p{=}0.53$; preliminary).
At $N{=}16$, the divergence persists: SJ maintains $\rho{=}0.308$ on conditional RE while LP drops to $\rho{=}{-}0.026$.
EM ranks first across all four tested conditional configurations: SciERC NER ($\rho{=}0.360$), CoNLL Qwen $N{=}16$ ($\rho{=}0.466$), CoNLL LLaMA 3-seed ($\rho{=}0.462$), and SciERC RE ($\rho{=}0.417$). This holds despite full-set rankings ranging from \#1 (CoNLL Qwen) to \#5 (SciERC RE).
On the full set, EM is diluted by trivial agreement: instances where all $N$ samples produce identical (possibly incorrect) outputs receive EM${=}1$ regardless of quality, compressing the signal's dynamic range.
On the conditional set, where the model demonstrably produces varied outputs, exact agreement becomes the strongest quality signal.
The margin over SJ is largest on RE ($\Delta\rho{=}{+}0.171$), where compound extraction decisions amplify output variability, and smallest on CoNLL LLaMA ($\Delta\rho{=}{+}0.002$), where all signals converge near $\rho{\approx}0.46$.

\paragraph{WNUT-17 full signal analysis.}
Table~\ref{tab:wnut} reports all five signals on WNUT-17 with full-set and conditional $\rho$.

\begin{table}[t]
\centering
\small
\caption{Quality estimation on WNUT-17 (Qwen, $N{=}16$, $T{=}1.0$, $n{=}689$). Cond: non-zero F1 instances ($n{=}334$). Full-set $\rho$ reflects binary success prediction; conditional $\rho$ measures quality gradation. LP's conditional $\rho$ is highest despite negative full-set $\rho$.}
\label{tab:wnut}
\begin{tabular}{lcccc}
\toprule
Signal & $\rho$ & AUROC & Cond.\ $\rho$ & Cond.\ $p$ \\
\midrule
VC      & .724 & .913 & .038 & .487 \\
FK      & .125 & .555 & .169 & .002 \\
SJ      & .089 & .517 & .288 & ${<}$1e-7 \\
EM      & .050 & .487 & .333 & ${<}$1e-9 \\
LP & $-$.410 & .219 & .350 & ${<}$1e-10 \\
\bottomrule
\end{tabular}
\end{table}
The full-set $p$-values are: VC ${<}$1e-112, FK $p{=}0.001$, SJ $p{=}0.020$, EM $p{=}0.194$, LP ${<}$1e-28.
VC's high full-set $\rho{=}0.724$ is entirely attributable to the zero/non-zero boundary (conditional $\rho{=}0.038$, $p{=}0.49$), while LP's conditional $\rho{=}0.350$ is the highest among all signals on the non-zero subset, full-set correlations on high-failure-rate datasets conflate binary success prediction with quality gradation.

\paragraph{Cross-configuration selection analysis.}
Table~\ref{tab:cross_selection} extends the selection analysis to five additional configurations.
On LLaMA SciERC, where output diversity remains substantial (18\% identical outputs, 39\% unique), LP achieves the highest selection F1 (0.652) at $N{=}8$ but still falls within bootstrap noise of single-sample (0.646), replicating the gap observed on Qwen SciERC (Table~\ref{tab:selection}).
On LLaMA Few-NERD ($n{=}32{,}565$, 3-seed mean), LP is the only signal with positive selection ($+0.36$~pp), consistent with LP being the most stable positive signal across model families (Qwen Few-NERD: $+1.39$~pp, Table~\ref{tab:fewnerd}). All structural signals degrade selection (VC $-2.39$~pp, SJ $-0.52$~pp), consistent with the Qwen pattern. LLaMA Few-NERD degeneracy ($23.4{\pm}4.7$\%, 3 seeds) is higher than Qwen ($13.0$\%), illustrating the capability-diversity tension: LLaMA achieves stronger greedy F1 ($0.784$ vs.\ $0.748$, $+3.6$~pp) but at the cost of higher degeneracy ($+10.4$~pp) and reduced oracle headroom ($10.0$ vs.\ $12.4$~pp).
On CoNLL, 48--87\% of instances produce identical outputs across all $N$ samples; this reflects a ceiling effect (single-sample F1 ${\approx}$0.91--0.92) rather than signal failure, and selection analysis on these configurations is uninformative because there is little per-sample variance for any signal to exploit.

\begin{table*}[t]
\centering
\small
\caption{Best-of-$N$ selection F1 (macro, instance-averaged) across configurations ($T{=}1.0$, 3-seed mean). Best non-oracle signal per row in \textbf{bold}. LP is the only signal with positive selection on both Few-NERD configurations (Qwen $+1.39$~pp, Table~\ref{tab:fewnerd}; LLaMA $+0.36$~pp). At $N{=}16$ on LLaMA CoNLL, VC achieves small but significant gains ($+0.28$~pp, $p{<}0.03$; see text below). CoNLL shows especially tight clustering where greedy already exceeds 0.90.}
\label{tab:cross_selection}
\begin{tabular}{llccccccc}
\toprule
Config & $n$ & Greedy & Oracle & SJ & FK & EM & LP & VC \\
\midrule
LLaMA/SciERC $N{=}8$ & 529 & .646 & .778 & .643 & .639 & .634 & \textbf{.652} & .633 \\
LLaMA/SciERC $N{=}16$ & 529 & .645 & .806 & \textbf{.649} & \textbf{.649} & .640 & .646 & .640 \\
\addlinespace
LLaMA/Few-NERD $N{=}8$ & 32{,}565 & .784 & .884 & .779 & .779 & .774 & \textbf{.788} & .760 \\
\addlinespace
Qwen/CoNLL $N{=}8$ & 2{,}756 & .908 & .948 & \textbf{.899} & .893 & .898 & .897 & .895 \\
LLaMA/CoNLL $N{=}8$ & 2{,}756 & .924 & .943 & .924 & \textbf{.925} & .924 & .924 & .923 \\
\bottomrule
\end{tabular}
\end{table*}

\paragraph{Diminishing returns of $N$-scaling.}
Our $N{=}8$ and $N{=}16$ comparisons (Table~\ref{tab:robustness}) characterize the marginal value of additional samples for test-time scaling \cite{brown2024monkeys, snell2025scaling}.
For quality estimation, structural signals benefit consistently but sublinearly: SJ $\rho$ improves by $+0.047$ (Qwen SciERC, 3-seed mean) and $+0.032$ (LLaMA CoNLL, 3-seed mean) when doubling from $N{=}8$ to $N{=}16$, while log-probability is approximately $N$-invariant ($\Delta\rho{\leq}0.021$).
Structural signals use $\binom{N}{2}$ pairwise comparisons (28 pairs at $N{=}8$, 120 at $N{=}16$), but the $\rho$ improvement is far less than proportional, each additional comparison yields diminishing information.
This asymmetry gives structural signals a relative advantage in the compute-constrained regime: at $N{=}8$, SJ already achieves $\rho{=}0.431$ on LLaMA CoNLL (vs.\ LP $\rho{=}0.307$, $\Delta{=}{+}0.124$), and the gap \textit{widens} at $N{=}16$ ($\Delta{=}{+}0.163$) because structural signals continue to benefit from additional pairwise evidence while LP does not.
For sample selection, additional samples increase the oracle upper bound (more draws raise the probability of a high-quality candidate), but signal-based selection remains flat because the selector's inability to identify the best candidate is a granularity limitation, not a sample size limitation.
The practical implication is that $N{=}8$ captures most of the achievable quality estimation benefit; further scaling widens rather than closes the correlation-selection gap absent stronger per-sample selectors.

\paragraph{$N{=}16$ CoNLL selection: small but significant.}
While quality estimation benefits sublinearly from $N$-scaling, selection behavior at $N{=}16$ shows a nuanced exception on LLaMA CoNLL.
VC achieves the largest gain over single-sample ($+0.28$~pp, $p{=}0.003$--$0.027$ across 3 seeds, all significant), with SJ, FK, and EM at $+0.22$--$0.23$~pp and LP at $+0.12$~pp (2/3 seeds significant).
These gains are statistically significant but practically minimal: oracle headroom at $N{=}16$ is only 2.5~pp, so VC captures approximately 10\% of the available headroom.
This reinforces the primary finding: even under the most favorable conditions (high baseline performance, more candidates, and aggregated multi-seed evidence), signal-based selection recovers only a fraction of the oracle headroom available.

\section{Bootstrap CI for Signal Comparisons}\label{app:bootstrap_signal}

Table~\ref{tab:bootstrap_signal_ci} reports bootstrap confidence intervals for individual signal $\rho$ and pairwise $\Delta\rho$ on SciERC NER ($N{=}8$, $n{=}529$, $B{=}10{,}000$, seed${=}42$).
The SJ--LP gap ($\Delta\rho{=}{+}0.155$) is clearly significant, the structural advantage reported in \S\ref{sec:qe_results} is robust.
The VC--SJ gap ($\Delta\rho{=}{+}0.019$) is not significant; these two structural signals perform comparably.
The SJ--EM gap ($\Delta\rho{=}{+}0.065$) is borderline significant, with the CI lower bound at $+0.003$ and upper bound at $+0.130$.

\begin{table}[t]
\centering
\small
\caption{Bootstrap 95\% CIs for signal $\rho$ on SciERC NER ($N{=}8$, $n{=}529$, $B{=}10{,}000$, seed${=}42$). Top: individual signals. Bottom: pairwise $\Delta\rho$; ``Sig.'' indicates whether the 95\% CI excludes zero.}
\label{tab:bootstrap_signal_ci}
\begin{tabular}{lcc}
\toprule
Signal & $\rho$ & 95\% CI \\
\midrule
VC & +0.379 & [+0.299, +0.456] \\
SJ & +0.360 & [+0.277, +0.437] \\
EM & +0.295 & [+0.212, +0.371] \\
FK & +0.267 & [+0.176, +0.354] \\
LP & +0.205 & [+0.119, +0.291] \\
\bottomrule
\end{tabular}

\vspace{0.5em}

\begin{tabular}{lccc}
\toprule
Comparison & $\Delta\rho$ & 95\% CI & Sig. \\
\midrule
SJ $-$ LP & +0.155 & [+0.074, +0.234] & Yes \\
VC $-$ SJ & +0.019 & [$-$0.033, +0.073] & No \\
SJ $-$ EM & +0.065 & [+0.003, +0.130] & Marginal$^\dagger$ \\
\bottomrule
\end{tabular}

\smallskip
{\footnotesize $^\dagger$Marginally significant without multiple comparison correction; would not survive Bonferroni correction for 10 pairwise comparisons ($\alpha_{\text{adj}}{=}0.005$).}
\end{table}

At $N{=}16$ ($B{=}10{,}000$, permutation test), the structural advantage extends to FK comparisons.
On SciERC NER ($n{=}529$), SJ--FK ($\Delta\rho{=}0.147$--$0.160$, all $p{\leq}0.002$) and VC--FK ($\Delta\rho{=}0.138$--$0.188$, all $p{\leq}0.003$) are clearly significant across all three seeds; EM--FK is seed-dependent ($\Delta\rho{=}0.054$--$0.093$, $p{=}0.048$--$0.15$).
On CoNLL ($n{=}2{,}756$), all signals significantly outperform FK: EM $\Delta\rho{=}0.071$ ($p{<}0.001$), SJ $\Delta\rho{=}0.056$ ($p{<}0.001$), VC $\Delta\rho{=}0.022$ ($p{=}0.040$).

\paragraph{Seed-replication analysis.}
Under a separate seed-replication analysis (5 SciERC seeds, Fisher $z$ + Holm correction), SJ~${>}$~LP reaches $p_{\text{adj}}{=}0.013$ ($\rho_{\text{SJ}}{=}0.356{\pm}0.014$, $\Delta\rho{=}0.151{\pm}0.006$), providing additional evidence for the robustness of structural signals.

\section{Holm--Bonferroni Enumeration}\label{app:holm}

Table~\ref{tab:holm_family} enumerates all $k{=}57$ statistical tests reported in this paper under a single global Holm--Bonferroni correction.
Under this most conservative family, no core NER claim survives correction ($p_{\text{adj}}{=}0.055$ for the most liberal core claim); two peripheral tests also reach significance (CoNLL few-shot entity entropy, RE LP selection reversal under free-form decoding).
We report per-family corrections in the main text (\S\ref{sec:prescriptive}); the global enumeration is provided here for transparency.

\begin{table*}[t]
\centering
\small
\caption{Complete enumeration of $k{=}57$ statistical tests subject to Holm--Bonferroni correction. Raw $p$: uncorrected $p$-value as reported. $p_{\text{adj}}$: Holm--Bonferroni adjusted $p$-value. $^\dagger$Survives correction at $\alpha{=}0.05$.}
\label{tab:holm_family}
\begin{tabular}{clllcc}
\toprule
\# & Family & Test Description & Raw $p$ & $p_{\text{adj}}$ & Sig. \\
\midrule
\multicolumn{6}{l}{\textit{Selection}} \\
1 & & SciERC NER LP selection ($N{=}8$, Qwen, Table~3) & $0.290$ & $1.0$ &  \\
2 & & SciERC NER SJ selection ($N{=}8$, Qwen, Table~3) & $0.670$ & $1.0$ &  \\
3 & & SciERC NER FK selection ($N{=}8$, Qwen, Table~3) & $0.790$ & $1.0$ &  \\
4 & & WNUT-17 LP selection ($N{=}8$, Qwen) & $0.720$ & $1.0$ &  \\
5 & & CoNLL LP selection ($N{=}8$, Qwen) & $1.0$ & $1.0$ &  \\
6 & & Few-NERD LP selection ($N{=}8$, Qwen, 3-seed) & ${<}.001$ & $0.055$ & \\
7 & & SciERC NER SJ selection ($N{=}8$, LLaMA-3B) & $0.175$ & $1.0$ &  \\
8 & & SciERC NER LP selection ($N{=}32$, Qwen) & $0.091$ & $1.0$ &  \\
9 & & SciERC NER FK selection ($N{=}32$, Qwen) & $0.101$ & $1.0$ &  \\
10 & & SciERC NER SJ selection ($N{=}32$, Qwen) & $0.244$ & $1.0$ &  \\
11 & & SciERC NER EM selection ($N{=}32$, Qwen) & $0.355$ & $1.0$ &  \\
12 & & SciERC NER VC selection ($N{=}32$, Qwen) & $1.0$ & $1.0$ &  \\
13 & & SciERC RE LP selection ($N{=}16$, 4B, Fisher) & $0.058$ & $1.0$ &  \\
14 & & SciERC RE VC selection ($N{=}16$, 4B, Fisher) & $0.973$ & $1.0$ &  \\
15 & & SciERC RE FK selection ($N{=}16$, 4B, Fisher) & $0.985$ & $1.0$ &  \\
16 & & SciERC RE SJ selection ($N{=}16$, 4B, Fisher) & $0.993$ & $1.0$ &  \\
17 & & SciERC RE EM selection ($N{=}16$, 4B, Fisher) & $0.999$ & $1.0$ &  \\
18 & & SciERC LP selection, non-tied subset ($\varepsilon{=}0.05$) & $0.004$ & $0.180$ &  \\
19 & & Few-NERD LP selection ($N{=}8$, LLaMA, 3-seed) & ${<}.001$ & $0.055$ & \\
20 & & SciERC LP favorable subset ($+$2.47 pp) & $0.015$ & $0.585$ &  \\
21 & & CoNLL LP favorable subset ($-$1.78 pp) & ${<}.001$ & $0.055$ & \\
\addlinespace
\multicolumn{6}{l}{\textit{Construction}} \\
22 & & SciERC LP-weighted construction (Qwen, 3-seed) & $0.004$ & $0.180$ &  \\
23 & & SciERC uniform construction (Qwen, 3-seed) & $0.023$ & $0.874$ &  \\
24 & & Few-NERD LP-weighted construction (Qwen) & ${<}.001$ & $0.055$ & \\
25 & & Few-NERD uniform construction (Qwen) & $0.003$ & $0.138$ &  \\
26 & & CoNLL uniform construction (Qwen) & $0.500$ & $1.0$ &  \\
27 & & SciERC ZS construction (72B) & $0.008$ & $0.320$ &  \\
28 & & SciERC FS construction (72B) & $0.544$ & $1.0$ &  \\
29 & & CoNLL FS construction (72B) & $0.417$ & $1.0$ &  \\
30 & & Few-NERD FS construction (72B) & $0.401$ & $1.0$ &  \\
31 & & SciERC cross-model LP-weighted construction & ${<}.001$ & $0.055$ & \\
\addlinespace
\multicolumn{6}{l}{\textit{Calibration}} \\
32 & & SciERC MF-Platt vs Raw LP & $0.720$ & $1.0$ &  \\
33 & & CoNLL MF-Platt vs Raw LP & $0.310$ & $1.0$ &  \\
34 & & Few-NERD MF-Platt vs Raw LP & $0.830$ & $1.0$ &  \\
\addlinespace
\multicolumn{6}{l}{\textit{Signal $\Delta\rho$}} \\
35 & & SciERC $\Delta\rho$ SJ$-$LP ($N{=}8$, seed 42) & ${<}.001$ & $0.055$ & \\
36 & & SciERC $\Delta\rho$ VC$-$SJ ($N{=}8$, seed 42) & $0.300$ & $1.0$ &  \\
37 & & SciERC $\Delta\rho$ SJ$-$EM ($N{=}8$, seed 42) & $0.045$ & $1.0$ &  \\
38 & & CoNLL $\Delta\rho$ EM$-$FK ($N{=}16$) & ${<}.001$ & $0.055$ & \\
39 & & CoNLL $\Delta\rho$ SJ$-$FK ($N{=}16$) & ${<}.001$ & $0.055$ & \\
40 & & CoNLL $\Delta\rho$ VC$-$FK ($N{=}16$) & $0.040$ & $1.0$ &  \\
41 & & SciERC $\Delta\rho$ SJ$-$FK (LLaMA, permutation) & $0.060$ & $1.0$ &  \\
42 & & SciERC RE $\Delta\rho$ SJ$-$FK (full-set) & $0.530$ & $1.0$ &  \\
43 & & RE SJ $\rho$ Q1 vs Q4 sentence length (Fisher $z$) & $0.004$ & $0.180$ &  \\
\addlinespace
\multicolumn{6}{l}{\textit{Correlation}} \\
44 & & Cross-dataset degeneracy--LP gain ($n{=}23$) & $0.810$ & $1.0$ &  \\
45 & & Per-type degeneracy--headroom (Few-NERD, Spearman) & $0.004$ & $0.180$ &  \\
46 & & Per-type degeneracy--headroom (exact permutation) & $0.005$ & $0.205$ &  \\
47 & & Per-type degeneracy--headroom ($\text{df}{=}4$ adjusted) & $0.033$ & $1.0$ &  \\
48 & & SciERC per-type degeneracy--LP ($\rho{=}0.14$) & $0.790$ & $1.0$ &  \\
49 & & Kruskal--Wallis per-type LP heterogeneity (SciERC) & $0.796$ & $1.0$ &  \\
\addlinespace
\multicolumn{6}{l}{\textit{FS error}} \\
50 & & FS vs ZS error overlap (SciERC, permutation) & $0.840$ & $1.0$ &  \\
51 & & FS vs ZS error overlap (CoNLL, permutation) & $0.640$ & $1.0$ &  \\
52 & & FS vs ZS entity entropy (SciERC) & $0.810$ & $1.0$ &  \\
53 & & FS vs ZS entity entropy (CoNLL, Wilcoxon) & ${<}10^{-4}$ & ${<}.001$ & $\dagger$ \\
\addlinespace
\multicolumn{6}{l}{\textit{Other}} \\
54 & & LP distribution seed 42 vs 456 (Mann--Whitney) & $0.550$ & $1.0$ &  \\
55 & & RE LP selection reversal (free-form, 4-seed) & ${<}.001$ & $0.006$ & $\dagger$ \\
56 & & WNUT-17 conditional $\rho$ VC ($p$-value) & $0.487$ & $1.0$ &  \\
57 & & WNUT-17 conditional $\rho$ FK ($p$-value) & $0.002$ & $0.094$ &  \\
\bottomrule
\end{tabular}
\end{table*}

\section{Diversity and Selection Analysis}\label{app:diversity}

% fig_selection_gap promoted to main text (\S\ref{sec:selection_gap})

\paragraph{Output diversity quantification.}
Table~\ref{tab:diversity} summarizes output diversity across configurations.
On SciERC, output diversity is substantial: Qwen produces 43\% unique outputs per instance (mean 3.5 distinct out of $N{=}8$), with per-sample F1 standard deviation of 0.144 and mean pairwise SJ of 0.703.
LLaMA shows similar diversity (39\% unique, F1 std 0.11).
On CoNLL, diversity collapses: 59\% (Qwen) to 89\% (LLaMA) of instances produce identical outputs across all $N$ samples; this reflects a ceiling effect where single-sample F1 already exceeds 0.90.
WNUT-17 shows 55\% identical outputs, attributable to the 53.7\% zero-F1 rate.

\begin{table}[t]
\centering
\small
\caption{Output diversity across configurations ($N{=}8$, $T{=}1.0$). ``Identical'' = fraction of instances with all $N$ outputs identical. ``Unique'' = mean fraction of distinct outputs per instance.}
\label{tab:diversity}
\begin{tabular}{lcc}
\toprule
Config & Identical & Unique \\
\midrule
Qwen/SciERC & 11\% & 43\% \\
LLaMA/SciERC & 18\% & 39\% \\
Qwen/CoNLL & 57\% & --- \\
LLaMA/CoNLL & 88\% & --- \\
WNUT-17 & 55\% & --- \\
\bottomrule
\end{tabular}
\end{table}

\paragraph{Oracle-improvable fraction.}
Across SciERC configurations, 39--42\% of instances contain at least one sample with F1 exceeding single-sample\footnote{Oracle-improvable and diversity-stratified analyses use $n{=}529$ (excluding instances with empty gold annotations; see \S\ref{sec:evaluation}); the diversity quantification in Table~\ref{tab:diversity} uses the full test set.}: Qwen NER 42.0\% (222/529, mean headroom $+$31.9~pp among improvable instances), Qwen RE 41.3\% (164/397, $+$45.3~pp), LLaMA NER 39.3\% (208/529, $+$33.7~pp).
On CoNLL, only 4--9\% of instances are improvable, consistent with the ceiling effect.

\paragraph{Diversity-stratified selection gap.}
On the 44\% of Qwen SciERC instances with high diversity (${\geq}4$ unique outputs among 8 samples, $n{=}232$), oracle headroom reaches $+$16.3~pp, yet all consistency-based signals (SJ, FK, EM, VC) select at or below single-sample, and LP, the only per-sample signal, achieves merely $+$0.8~pp.
LLaMA replicates this pattern: on its high-diversity subset ($n{=}186$), oracle headroom reaches $+$20.1~pp while consistency signals remain at or below single-sample and LP captures only $+$2.0~pp.
This rules out insufficient diversity as an explanation for the correlation-selection gap.

\paragraph{Instance degeneracy statistics.}
Table~\ref{tab:degeneracy} quantifies instance degeneracy across configurations. ``Constant F1'' denotes instances where all $N$ samples receive identical instance-level F1; ``Constant utility'' denotes instances where MBR (SJ) utility is identical across candidates, preventing ranking.

\begin{table}[t]
\centering
\small
\caption{Instance degeneracy statistics ($T{=}1.0$). ``Constant F1'': all $N$ samples receive identical instance-level F1. ``Constant utility'': MBR (SJ) utility identical across candidates. 4B data confirms degeneracy scales with model capability.}
\label{tab:degeneracy}
\resizebox{\columnwidth}{!}{
\begin{tabular}{lccc}
\toprule
& \shortstack{SciERC 8B\\($N{=}8$)} & \shortstack{SciERC 4B\\($N{=}8$)} & \shortstack{CoNLL 8B\\($N{=}8$)} \\
\midrule
Constant F1 & 12.2\% & 10.4\% & 61.7\% \\
Constant utility & 13\% & 7.0\% & 64\% \\
Oracle gap closed (non-deg.) & 23\% & 25.2\% & 36\% \\
Within-inst.\ SJ $\rho$ (median) & 0.524 & 0.550 & 0.990 \\
\bottomrule
\end{tabular}
}
\end{table}

\paragraph{TOST equivalence testing.}\label{app:tost}
The primary evidence against selection gains comes from effect sizes and confidence intervals: Cohen's $d{=}0.038$ for the best signal (NER LP) and bootstrap 95\% CI [$-$3.0, $+$4.2]~pp spanning zero.
Table~\ref{tab:tost} supplements this with two one-sided tests (TOST), testing whether each signal's selection F1 is equivalent to single-sample within a bound $\delta$.
The $\delta$ values are contextualized against oracle headroom (12--18~pp): $\delta{=}5$~pp corresponds to forgoing 28--38\% of the available headroom, $\delta{=}3$~pp forgoing 17--23\%, and $\delta{=}2$~pp less than 15\%.
These proportions use the same test data that defines the headroom, introducing circularity: $\delta$ is anchored to the very quantity it contextualizes.
We address this below with externally anchored equivalence margins.
The breakpoint $\delta^*$, the minimum equivalence margin at which each signal passes TOST, ranges from approximately 1.3~pp (SJ, ${\sim}$10\% of oracle headroom) to 3.4~pp (FK, ${\sim}$26\%) on NER, and from approximately 2.9~pp (LP, ${\sim}$16\%) to 6.5~pp (VC, ${\sim}$36\%) on RE.
On NER ($n{=}529$), at $\delta{=}{\pm}2$~pp, SJ and LP are confirmed equivalent to single-sample (TOST $p{<}0.05$); at $\delta{=}{\pm}3$~pp, four of five signals pass; at $\delta{=}{\pm}5$~pp, all five pass.
On RE ($n{=}397$, preliminary), EM and VC fail equivalence even at $\delta{=}{\pm}5$~pp due to systematic underperformance ($\Delta{=}{-}3.2$~pp and $-4.5$~pp respectively).
The effect sizes are negligible ($d{=}0.038$) and confidence intervals span zero, with TOST confirming equivalence at practically meaningful margins.

\begin{table}[t]
\centering
\small
\caption{TOST equivalence testing: signals vs.\ single-sample on SciERC ($N{=}8$, $T{=}1.0$). \checkmark: equivalent at $p{<}0.05$; \texttimes: not equivalent. $\Delta$: mean F1 difference (signal $-$ single-sample).}
\label{tab:tost}
\resizebox{\columnwidth}{!}{
\begin{tabular}{l cccc cccc}
\toprule
& \multicolumn{4}{c}{\textbf{NER}} & \multicolumn{4}{c}{\textbf{RE}} \\
\cmidrule(lr){2-5} \cmidrule(lr){6-9}
Signal & $\Delta$ & $\pm$2 & $\pm$3 & $\pm$5 & $\Delta$ & $\pm$2 & $\pm$3 & $\pm$5 \\
\midrule
LP & +0.6 & \checkmark & \checkmark & \checkmark & $-$0.9 & \texttimes & \checkmark & \checkmark \\
SJ & +0.1 & \checkmark & \checkmark & \checkmark & $-$1.8 & \texttimes & \texttimes & \checkmark \\
FK & $-$0.5 & \texttimes & \checkmark & \checkmark & $-$1.6 & \texttimes & \texttimes & \checkmark \\
EM & $-$0.8 & \texttimes & \checkmark & \checkmark & $-$3.2 & \texttimes & \texttimes & \texttimes \\
VC & $-$1.0 & \texttimes & \texttimes & \checkmark & $-$4.5 & \texttimes & \texttimes & \texttimes \\
\bottomrule
\end{tabular}
}
\end{table}

\paragraph{External equivalence margins.}
The oracle-headroom anchors above introduce circularity: $\delta$ is calibrated against the very quantity it contextualizes.
To resolve this, we supplement with three externally grounded margins:
(1) $\delta_{\mathrm{SEM}}$: the bootstrapped standard error of single-sample F1 (0.43--1.33~pp depending on configuration), which represents measurement noise;
(2) $\delta_{\mathrm{practical}}{=}0.5$~pp, the threshold below which NER differences are considered negligible \cite{ratinov2009design};
(3) $\delta_{\mathrm{literature}}{=}1.0$~pp, the approximate significant-difference threshold established in NER evaluation methodology \cite{bergkirkpatrick2012empirical}.

Results under these external anchors diverge sharply by configuration.
On LLaMA CoNLL ($n{=}2{,}756$), all signals pass equivalence at every $\delta$ ($p{<}0.001$), consistent with the near-zero selection deltas ($|\Delta|{\leq}0.1$~pp) and tight CIs at this sample size.
On Qwen CoNLL ($n{=}2{,}756$), no signal passes at any external $\delta$: selection F1 degrades by 0.9--1.5~pp relative to single-sample, which exceeds even $\delta_{\mathrm{literature}}$.
On SciERC NER ($n{=}529$), only SJ passes at $\delta_{\mathrm{SEM}}{=}1.33$~pp (the widest SEM bound); all other signals and tighter margins fail.
These external anchors confirm that the equivalence conclusion is configuration-dependent: LLaMA CoNLL selection is robustly equivalent to single-sample under any reasonable margin, while SciERC and Qwen CoNLL selection shows detectable degradation at externally calibrated thresholds.

\paragraph{Multiple comparison considerations.}\label{app:multiple_comparisons}
Our primary selection analysis involves ${\sim}$20 comparisons (5 signals $\times$ 3--4 configurations).
Since all primary selection results are non-significant ($p{>}0.05$), applying Benjamini-Hochberg correction can only make these results more conservative, strengthening the ``no signal beats greedy'' conclusion.
The sole positive finding, the non-tied LP subset ($+2.9$~pp, $p{=}0.004$; \S\ref{sec:discussion}), survives Bonferroni correction over the 5 signal comparisons ($p_{\mathrm{adj}}{=}0.020 < 0.05$). However, the $\varepsilon{=}0.05$ threshold is selected post-hoc from a 7-point sweep (Appendix~\ref{app:epsilon_robustness}); joint correction across thresholds and signals yields $p_{\mathrm{adj}}{=}0.14$, so this finding is exploratory.
For quality estimation signal rankings, differences between the top signals (e.g., VC vs.\ SJ on SciERC) should be interpreted as descriptive rather than as statistically confirmed orderings, given the limited power of 3-seed comparisons.

\paragraph{When would consistency-based selection work?}\label{app:when_selection_works}
Our negative results are specific to fine-tuned models with constrained decoding on structured IE. Several conditions outside our experimental scope could restore selection headroom.
(a)~Models with poor calibration or in-context learning setups (without fine-tuning) produce more diverse outputs, widening the gap between candidates that consistency signals can exploit.
(b)~Closed-source APIs where log-probabilities are unavailable make consistency signals the only viable selection mechanism; our structural signals (SJ, VC, FK, EM) provide quality estimation ($\rho$ up to 0.45) even without model-internal access.
(c)~Ensemble or mixture-of-experts settings generate structurally diverse candidates by design, reducing instance degeneracy below the levels observed in our single-model sampling.
(d)~Tasks where the model exhibits systematic overconfidence (e.g., hallucination-prone generation) may benefit from consensus filtering, as disagreement among samples would carry more discriminative signal than in our well-calibrated fine-tuned regime.

\section{Supervised and Model-Internal Signal Analysis}\label{app:supervised_combination}

\subsection{Supervised Signal Combination}

Supervised classifiers (logistic regression, MLP, random forest) trained on all five signals (5-fold CV, 3-seed mean, SciERC, $N{=}16$, $+17.9$~pp headroom) all perform at or below greedy (F1${=}0.644$; LR: $-0.4$~pp; MLP: $-1.6$~pp; RF: $-3.1$~pp).
Adding pairwise interaction features does not help (AUROC 0.634 vs.\ 0.638).
To test whether richer representations could close the gap, we train a DeBERTa-v3-base verifier that encodes input-output pairs and predicts per-sample F1.
DeBERTa achieves $\rho{=}0.329{\pm}0.087$ (5-fold CV, 3 seeds). Despite moderate quality prediction, selection utility is poor: on SciERC, DeBERTa closes $-26.5{\pm}26.5$\% of the oracle gap (13 of 15 runs negative; vs.\ LP's $+4.2{\pm}10.7$\%); on Few-NERD, DeBERTa recovers only $+3.5$\% (single 80/20 split; vs.\ LP's $+16.3$\%).
Cross-dataset transfer fails entirely ($-27$\% to $-72$\% gap closure).
Grid-search ensemble optimization over signal weights confirms LP dominance: across configurations, LP receives mean weight $0.78$, while FK and SJ receive zero weight; VC and EM serve only as secondary regularizers.
Neither handcrafted signal combinations nor learned text representations close the selection gap; only model-internal log-probability provides meaningful selection signal, consistent with an information asymmetry between the generating model and external observers.

\subsection{Entity-Level Verifier Transfer}\label{app:verifier_transfer}

A LightGBM entity-level verifier (5 features: agreement count, VC, LP, SJ, FK) achieves AUC${=}0.923{\pm}0.002$ under out-of-fold evaluation, confirming that entity correctness is learnable from the available signals.
In-domain OOF construction yields $+0.30$\,pp overall, with per-dataset gains of $+1.28$\,pp (SciERC), $+0.55$\,pp (CoNLL), and $+0.02$\,pp (Few-NERD).
The near-zero Few-NERD gain despite high AUC shows that frequency-based threshold construction already captures most of the verifier's signal.
Cross-dataset transfer to Few-NERD is negative: $-0.76$\,pp (B2-v2), $-1.02$\,pp (B2-v3), $-1.27$\,pp (B2-v4 Replace, AUC${=}0.822$), and $-0.85$\,pp (B2-v4 Augment, AUC${=}0.855$), all below greedy.
The learned feature--correctness mapping is dataset-specific: entity agreement patterns that predict correctness on scientific text do not transfer to fine-grained entity types.

\subsection{Hidden State Probes}\label{app:hidden_probe}

Entity-level probes confirm that quality information is partially encoded in hidden representations (AUROC${=}0.63$\textendash$0.88$, dataset-dependent; \S\ref{sec:mechanistic}); we test whether this translates to sample-level selection gain by training probes on the generating model's hidden states.

\paragraph{Setup.}
Unlike the entity-level probe in \S\ref{sec:mechanistic} (which classifies individual entities on Few-NERD/CoNLL/SciERC), this sample-level analysis uses all 551 SciERC instances (including 22 with empty gold annotations) to predict per-sample F1.
For each of $N{=}8$ samples per instance, we extract the mean-pooled final-layer hidden state ($d{=}4096$) from the generating Qwen3-8B model and train a linear probe and a two-layer MLP (hidden dim${=}256$) to predict per-sample F1 (5-fold CV, 3 seeds).
Entity-level probes use dataset-specific generation temperatures ($T{=}0.5$ for Few-NERD, $T{=}0.8$ for CoNLL, $T{=}1.0$ for SciERC).

\begin{table}[t]
\centering
\small
\caption{Hidden state probes vs.\ LP for sample selection (5-fold CV ${\times}$ 3 seeds). $\rho$: Spearman correlation with per-sample F1. Sel.\ F1: selection F1. Gap\%: oracle gap closure. SciERC greedy${=}0.644$ (seed 42), oracle${=}0.775$ (seed 42).}
\label{tab:hidden_probe}
\resizebox{\columnwidth}{!}{%
\begin{tabular}{@{}l ccc cc@{}}
\toprule
& \multicolumn{3}{c}{\textbf{SciERC}} & \multicolumn{2}{c}{\textbf{Few-NERD}} \\
\cmidrule(lr){2-4} \cmidrule(lr){5-6}
Method & $\rho$ & Sel.\ F1 & Gap\% & $\rho$ & Gap\% \\
\midrule
LP & .283 & $\mathbf{.650}$ & $+$4.8 & .606 & 12.3 \\
Hidden Linear & $.179{\pm}.019$ & $.613{\pm}.002$ & $-$23.4 & \multicolumn{2}{c}{---} \\
Hidden MLP & $\mathbf{.388{\pm}.013}$ & $.637{\pm}.003$ & $-$5.1 & $\mathbf{.629{\pm}.007}$ & 11.7 \\
\bottomrule
\end{tabular}}
\end{table}

\paragraph{Cross-dataset asymmetry.}
On SciERC, the MLP probe achieves higher quality correlation than LP ($\rho{=}0.388$ vs.\ $0.283$, $+37$\%); hidden states encode information beyond what LP captures.
Yet higher correlation does not yield better selection: MLP selection F1 ($0.637$) falls below LP's ($0.650$), with gap closure turning negative ($-5.1$\% vs.\ $+4.8$\%; Table~\ref{tab:hidden_probe}).
The linear probe performs worse still ($\rho{=}0.179$, gap closure $-23.4$\%).
On Few-NERD, the MLP provides no increment over LP (gap closure $11.7$\% vs.\ $12.3$\%).
A learned token-level LP aggregator (1D CNN over per-token log-probabilities) matches mean LP on Few-NERD (gap closure $56.4$\% vs.\ $55.8$\%); LP's bottleneck is not token aggregation.
Despite partial entity quality encoding in hidden representations (entity-level probe AUROC${=}0.63$\textendash$0.88$; \S\ref{sec:mechanistic}), the correlation-selection gap persists at the sample level: in the high-capability regime (greedy F1${=}0.644$), inter-sample F1 differences are small enough that improved ranking provides no marginal selection value, consistent with the capability-diversity tension (\S\ref{sec:discussion}).

\subsection{Method Coverage Matrix}\label{app:method_coverage}

Table~\ref{tab:method_coverage} summarizes the experimental coverage of each entity-level method evaluated in this work. Individual methods were tested on subsets of the full dataset $\times$ model matrix; the collective coverage spans three NER datasets and two model families.

\begin{table}[h]
\centering
\small
\caption{Experimental coverage of entity-level methods. \checkmark = tested; --- = not tested. All methods use $N{=}8$, $T{=}1.0$ unless noted. Section references indicate where results appear.}
\label{tab:method_coverage}
\begin{tabular}{lccccl}
\toprule
Method & SciERC & Few-NERD & CoNLL & Qwen & LLaMA \\
\midrule
ESJ (\S\ref{sec:mechanistic}) & --- & \checkmark & --- & \checkmark & --- \\
MRSC (\S\ref{sec:mechanistic}) & --- & \checkmark & --- & \checkmark & --- \\
LSC (\S\ref{sec:mechanistic}) & --- & \checkmark & --- & \checkmark & --- \\
Het-T (App.~\ref{app:robustness}) & \checkmark & --- & --- & \checkmark & --- \\
ASB (App.~\ref{app:dals}) & \checkmark & \checkmark & \checkmark & \checkmark & --- \\
SCD (App.~\ref{app:supervised_combination}) & \checkmark & \checkmark & \checkmark & \checkmark & --- \\
Verifier (\S\ref{sec:verifier}) & \checkmark & \checkmark & \checkmark & \checkmark & --- \\
\midrule
QE signals (\S\ref{sec:qe_results}) & \checkmark & \checkmark & \checkmark & \checkmark & \checkmark \\
\bottomrule
\end{tabular}
\end{table}

\section{How Self-Verification Fails}\label{sec:mechanistic}

The diagnostic framework (\S\ref{sec:method}) identifies \textit{where} self-consistency fails; this section probes \textit{why} by testing whether models can self-verify at the entity level.
Probe experiments use Qwen2.5-7B (not the fine-tuned Qwen3-8B used in main experiments); conclusions are architecture-specific.
We evaluate self-verification and diagnostic probes on Qwen2.5-7B\footnote{ESJ, MRSC, and LSC use Qwen2.5-7B-Instruct (zero-shot self-verification capability) rather than the fine-tuned Qwen3-8B used in the main experiments. Conclusion direction is consistent with the main experiments' negative results across all fine-tuned selection strategies. The pattern---selection failure across all signals---is independently confirmed in the Qwen3-8B main experiments, making the introspection boundary a model-independent finding within our tested scale. Probe conclusions are limited to the Qwen2.5 architecture; linear probes may underestimate representational capacity (nonlinear alternatives are left to future work).} fine-tuned for NER on Few-NERD ($N{=}8$, seed 42; full results in Table~\ref{tab:self_verification}).

\paragraph{Entity Self-Judgment (ESJ).}
For each extracted entity, we prompt the generating model: ``Is [\textit{entity}] a [\textit{type}] entity in this sentence?''
Entity-level AUROC is $0.512$---indistinguishable from chance---with P(yes$|$correct)${=}0.500$ and P(yes$|$incorrect)${=}0.500$.
ESJ-weighted construction degrades F1 to $0.711$ (vs.\ greedy $0.794$).

\paragraph{Multi-Round Self-Consistency (MRSC).}
Following iterative self-refinement \cite{madaan2023selfrefine}, we re-prompt on disputed entities (agreement ${<}0.6$) with focused verification; P(yes) separation between gold and non-gold entities is $0.0006$ ($0.5003$ vs.\ $0.4997$).
MRSC best F1 ($0.763$) falls below greedy.

\paragraph{Latent Self-Consistency (LSC).}
Mean-pooled final-layer hidden states ($d{=}3584$) provide per-sample latent representations.
Centroid selection ($0.759$) underperforms greedy ($0.794$); latent entity confidence ($0.839$) outperforms majority vote ($0.792$) but does not close the oracle gap ($0.882$).

\paragraph{Hidden-state probe.}
A logistic regression on mean-pooled final-layer hidden states at entity spans achieves AUROC $0.636$ (Few-NERD), $0.882$ (CoNLL), $0.630$ (SciERC), evaluated with GroupKFold CV (5-fold ${\times}$ 3 seeds, groups${=}$source instance to prevent entity-level leakage; \S\ref{app:hidden_probe}).
All three exceed ESJ ($0.51$); only CoNLL approaches frequency agreement ($0.868$).
Cross-dataset transfer is weak: leave-one-dataset-out AUROC drops to $0.559$ (Few-NERD${\to}$CoNLL) and $0.671$ (CoNLL${\to}$Few-NERD), indicating that entity error patterns are dataset-specific rather than reflecting a universal introspection boundary, limiting probe-based approaches without per-dataset calibration.
At the sample level, an MLP probe closes only 13.9\% of the oracle gap on SciERC (\S\ref{app:hidden_probe}).

\paragraph{The introspection boundary.}
Entity quality is partially encoded in hidden representations: a supervised probe (AUROC $0.63$\textendash$0.88$) consistently exceeds self-judgment (ESJ $0.51$) but only marginally exceeds multi-sample agreement ($0.87$) on CoNLL, suggesting consensus captures information beyond single-pass representations while the prompt interface fails to access even partial hidden-state signal.
LODO transfer drops to $0.56$\textendash$0.67$, so even direct hidden-state access would not yield a universal quality signal---consistent with the negative cross-dataset verifier transfer (\S\ref{sec:verifier}).

\section{Entity-Level Self-Verification Probes}\label{app:self_verification}

Three independent probes test whether the generating model can self-verify at the entity level, using Qwen2.5-7B fine-tuned on Few-NERD ($N{=}8$, seed 42).

\begin{table}[t]
\centering
\small
\caption{Entity-level self-verification probes (Qwen2.5-7B, Few-NERD, $N{=}8$). ESJ: entity self-judgment via yes/no prompts. MRSC: multi-round re-prompting on disputed entities. LSC: latent-space centroid selection. MV: majority vote baseline.}
\label{tab:self_verification}
\begin{tabular}{@{}lcc@{}}
\toprule
Method & Key Metric & F1 \\
\midrule
Greedy ($T{=}0$) & --- & .794 \\
MV ($\theta{=}0.375$) & --- & .792 \\
Oracle & --- & .882 \\
\midrule
ESJ (best $\theta{=}0.02$) & AUROC${=}0.512$ & .711 \\
MRSC (best $\theta{=}0.5$) & $\Delta P(\text{yes}){=}0.001$ & .763 \\
LSC (centroid) & $\rho_{\text{LP-latent}}{=}0.52$ & .759 \\
LSC (entity conf.) & --- & .839 \\
\midrule
Freq.\ agreement & AUROC${=}0.868$ & --- \\
\bottomrule
\end{tabular}
\end{table}

\paragraph{Entity Self-Judgment (ESJ).}
For each extracted entity, we prompt the model: ``Is [entity] a [type] entity in this sentence?''
The entity-level AUROC is $0.512$---indistinguishable from chance---with P(yes$|$correct)${=}0.500$ and P(yes$|$incorrect)${=}0.500$ (separation $0.001$).
ESJ-weighted construction degrades F1 to $0.711$ (vs.\ greedy $0.794$), confirming the model lacks entity-level self-judgment.

\paragraph{Multi-Round Self-Consistency (MRSC).}
We re-prompt the model on disputed entities (agreement ${<}0.6$) with a focused verification query and aggregate with Round~2 weight${=}2$.
P(yes) separation between gold and non-gold entities is $0.0006$ ($0.5003$ vs.\ $0.4997$), and MRSC best F1 ($0.763$) falls below both greedy and MV.

\paragraph{Latent Self-Consistency (LSC).}
Mean-pooled final-layer hidden states ($d{=}3584$) provide latent representations per sample.
Centroid selection ($0.759$) underperforms greedy ($0.794$); latent entity confidence ($0.839$) outperforms MV ($0.792$) but does not close the oracle gap ($0.882$).

\paragraph{Frequency agreement as the viable signal.}
Cross-sample entity frequency agreement achieves AUROC${=}0.868$ for entity correctness prediction; a supervised hidden-state probe (AUROC${=}0.63$\textendash$0.88$; \S\ref{sec:mechanistic}) exceeds this only on the simplest dataset (CoNLL, $0.882$), while falling below on complex NER tasks.
The probe consistently exceeds prompt-based self-evaluation (ESJ AUROC${=}0.512$), confirming that hidden states encode partial entity quality signal inaccessible through the prompt interface, consistent with the complexity-dependent introspection boundary diagnosed in the main text (\S\ref{sec:mechanistic}).

\section{Cross-Dataset Entailment Self-Judgment}\label{app:esj}

Table~\ref{tab:esj_cross} compares entailment-based self-judgment (ESJ) with agreement-based signals across three datasets.
ESJ achieves near-random AUROC on all datasets, confirming that the self-verification failure observed on Few-NERD (\S\ref{sec:mechanistic}) generalizes.

\begin{table}[h]
\small
\centering
\caption{Entity-level self-verification: ESJ vs.\ agreement-based signals across datasets (Qwen3-8B, $N{=}8$). ESJ AUROC is indistinguishable from chance on all datasets.}
\label{tab:esj_cross}
\begin{tabular}{lcccc}
\toprule
Dataset & ESJ AUROC & ESJ $\rho$ & Agree.\ AUROC & Agree.\ $\rho$ \\
\midrule
SciERC  & 0.512 & 0.020 & 0.792 & 0.515 \\
CoNLL   & 0.553 & 0.087 & 0.941 & 0.777 \\
Few-NERD & 0.495 & $-$0.008 & 0.881 & 0.673 \\
\bottomrule
\end{tabular}
\end{table}

\section{Multi-Seed Stability}\label{app:multiseed}

Table~\ref{tab:multiseed} reports quality estimation metrics across multiple seeds (Qwen: 5; LLaMA: 3) for both model families on SciERC ($N{=}8$, $T{=}1.0$).
On NER, the signal ranking VC${>}$SJ${>}$EM${>}$FK${>}$LP is identical across all model-seed combinations, with $\sigma{\leq}0.025$ for all signals (rankings are stable across all seeds).
On RE (Qwen only, preliminary), VC remains the top signal ($\rho{=}0.309{\pm}0.024$), but LP rises to \#2 ($\rho{=}0.297{\pm}0.022$), above FK, SJ, and EM. This reversal from the NER ranking reflects RE's higher sensitivity to per-sample scoring.
LLaMA absolute values are systematically lower than Qwen on NER, consistent with the directionally-consistent-but-magnitude-dependent characterization in the main text.

\begin{table}[t]
\centering
\small
\caption{Multi-seed quality estimation on SciERC ($N{=}8$, $T{=}1.0$; Qwen: 5 seeds (42, 123, 456, 2024, 3141); LLaMA: 3 seeds (42, 123, 456)). NER: both models; RE: Qwen only (preliminary). Signal ranking is stable across seeds ($\sigma{\leq}0.025$). $\sigma$ is population standard deviation over seeds.}
\label{tab:multiseed}
\resizebox{\columnwidth}{!}{
\begin{tabular}{lccccc}
\toprule
& LP & FK & SJ & VC & EM \\
\midrule
\multicolumn{6}{l}{\textit{NER Spearman $\rho$ (mean $\pm\sigma$)}} \\
Qwen & .205${\pm}$.015 & .254${\pm}$.016 & .356${\pm}$.014 & .386${\pm}$.019 & .307${\pm}$.008 \\
LLaMA & .162${\pm}$.025 & .256${\pm}$.024 & .325${\pm}$.021 & .384${\pm}$.014 & .279${\pm}$.011 \\
\midrule
\multicolumn{6}{l}{\textit{RE Spearman $\rho$ (mean $\pm\sigma$, Qwen)}} \\
Qwen & .297${\pm}$.022 & .290${\pm}$.012 & .249${\pm}$.002 & .309${\pm}$.024 & .135${\pm}$.016 \\
\midrule
\multicolumn{6}{l}{\textit{NER AUROC (mean $\pm\sigma$)}} \\
Qwen & .579${\pm}$.008 & .600${\pm}$.016 & .659${\pm}$.016 & .677${\pm}$.011 & .631${\pm}$.013 \\
LLaMA & .565${\pm}$.008 & .593${\pm}$.006 & .632${\pm}$.012 & .673${\pm}$.009 & .606${\pm}$.005 \\
\bottomrule
\end{tabular}
}
\end{table}

\section{Training Configuration}\label{app:training}

Both models are fine-tuned with identical hyperparameters (Table~\ref{tab:training}) to ensure comparisons across model families reflect model differences rather than training confounds.

\begin{table}[t]
\centering
\small
\caption{Training configuration. Both Qwen3-8B and LLaMA-3.1-8B-Instruct use identical settings unless noted.}
\label{tab:training}
\resizebox{\columnwidth}{!}{
\begin{tabular}{ll}
\toprule
Parameter & Value \\
\midrule
LoRA rank / $\alpha$ & 32 / 64 \\
LoRA target & All linear modules \\
LoRA dropout & 0 \\
Optimizer & AdamW \\
Learning rate & 2e-4, cosine schedule \\
Warmup ratio & 0.1 \\
Effective batch size & 16 (2 $\times$ grad.\ accum.\ 8) \\
Epochs & 5 \\
Max sequence length & 2048 \\
Precision & bf16 \\
Gradient checkpointing & Enabled \\
Max gradient norm & 1.0 \\
Training seed & 42 \\
Framework & LLaMA Factory \\
\midrule
\multicolumn{2}{l}{\textit{Dataset-specific}} \\
SciERC split & 1861 / 275 / 551 \\
CoNLL / WNUT-17 & +\,load\_best\_model\_at\_end, \\
& eval\_steps${=}$200, save\_limit${=}$5 \\
\bottomrule
\end{tabular}
}
\end{table}

\paragraph{Prompt template.}
For NER, we use a structured prompt:
\begin{quote}
\small
\texttt{System:} You are an expert in Named Entity Recognition.\\
\texttt{User:} Extract all named entities from the following text. Output as JSON.\\
Text: \texttt{\{input\_text\}}\\
Entity types: \texttt{\{schema\_types\}}
\end{quote}
RE follows the same structure with ``Relation Extraction'' and relation type lists substituted accordingly.

\paragraph{Constrained decoding.}
Inference uses vLLM with XGrammar-based JSON schema enforcement \cite{dong2024xgrammar, willard2023guided}, which constrains outputs to valid entity/relation structures matching the target schema.
This eliminates format errors but may compress per-sample diversity by restricting the output space.\footnote{SciERC training used \texttt{enable\_thinking=True} (LLaMA Factory 0.9.4 default), injecting empty \texttt{<think></think>} blocks (4--5 tokens) into each training response with loss computation; CoNLL and WNUT explicitly set \texttt{enable\_thinking=False}. Constrained decoding at inference time enforces JSON-only output, suppressing thinking tokens. See \S\ref{sec:setup} and Limitations for discussion of this confound.}

\section{Thinking Token Impact Analysis}\label{app:thinking_token}

Qwen3-8B's SciERC checkpoint was trained with empty \texttt{<think></think>} blocks prepended to each response (\S\ref{sec:setup}).
We quantify the potential impact on LP scores through three complementary arguments.

\paragraph{LP shift magnitude.}
Empty thinking tokens (mean 4.5 per sample) constitute 2.1\% of tokens on SciERC.
The estimated LP shift is 0.001 nats per token (2.68\% of the median within-instance LP range of 0.041 nats).
Because all samples share identical thinking-token LP offsets, cumulative LP rankings are invariant; mean LP changes selection for ${<}1.5\%$ of instances.

\paragraph{Constant prefix cancellation.}
All $N$ samples within an instance share the same prompt and, under constrained decoding, the same absence of thinking tokens at inference time.
Any LP bias from the training-time thinking prefix is therefore a near-constant additive shift across all candidates within an instance, leaving within-instance LP \textit{rankings} unchanged.

\paragraph{CoNLL control group.}
CoNLL and WNUT-17 checkpoints were trained without thinking tokens (\texttt{enable\_thinking=False}).
If thinking tokens caused LP compression on SciERC, CoNLL should exhibit lower tied fractions.
Instead, CoNLL shows \textit{higher} tied fractions at both temperatures (86.8\% at $T{=}1.0$, 95.8\% at $T{=}0.7$) than SciERC (60.3\% at $T{=}1.0$), directly contradicting the thinking-token hypothesis.

\section{Log-Probability Score Compression}\label{app:logprob_ties}

Under constrained decoding, log-probability scores are severely compressed within instances, limiting per-sample discrimination (Table~\ref{tab:logprob_ties}, Figure~\ref{fig:lp_range_cdf}).
The median within-instance LP range is only 0.041 nats; we use $\varepsilon{=}0.05$ as an illustrative threshold to characterize this compression.
At this cutpoint, 60.3\% of instances ($n{=}529$, excluding empty gold annotations) have a within-sample range below $\varepsilon$, and 75\% of pairwise comparisons are effective ties ($|\Delta|{<}0.026$).
The CDF rises smoothly with no cliff or discontinuity (Figure~\ref{fig:lp_range_cdf}); compression is a pervasive regime rather than a bimodal split.
This compression is stable across seeds: seed${=}456$ yields median LP range${=}0.040$ nats and tied fraction${=}62.2\%$ (vs.\ seed${=}42$: $0.041$ nats, $60.3\%$).
Seed~456 uses a different subtask flag (joint vs.\ ner for seeds 42 and 123); the LP distributions are statistically indistinguishable (median LP: 8.37 vs.\ 8.20 nats; Mann--Whitney $p{=}0.55$); the subtask change does not affect compression.
For comparison, SJ has only 12.5\% tied instances at the same threshold; SJ's advantage over LP for selection is one of score resolution, not information content.

\begin{figure}[t]
\centering
\includegraphics[width=\columnwidth]{fig_lp_range_cdf}
\caption{CDF of within-instance LP range on SciERC NER ($N{=}8$, $n{=}529$). The median LP range is 0.041 nats (dashed line), below $\varepsilon{=}0.05$. The smooth CDF confirms compression is pervasive, not bimodal.}
\label{fig:lp_range_cdf}
\end{figure}

\begin{table}[t]
\centering
\small
\caption{Log-probability score compression on SciERC NER ($N{=}8$, $n{=}529$). ``Tied'' = fraction of instances where within-instance LP range ${<}\varepsilon$.}
\label{tab:logprob_ties}
\begin{tabular}{lc}
\toprule
$\varepsilon$ threshold & Tied instances (\%) \\
\midrule
0.01 & 1.1\% \\
0.05 & 60.3\% \\
0.1  & ${\sim}$84\% \\
0.5  & ${\sim}$99\% \\
\bottomrule
\end{tabular}
\end{table}

\paragraph{Non-tied subset.}
Among instances with LP range ${\geq}0.05$ ($n{=}210$, 40\% of total), log-probability selection \textit{significantly} outperforms single-sample ($+2.9$~pp F1, $p{=}0.004$), and within-instance $\rho(\text{LP}, \text{F1})$ rises from 0.08 (tied subset) to 0.39 (non-tied subset).
As a continuous complement to the binary tied/non-tied analysis, pooled within-instance $\rho(\text{LP}, \text{F1}){=}0.248$ across all instances confirms weak but positive per-sample signal; however, only 21.6\% of instances achieve $p{<}0.05$ significance, and $\rho(\text{LP range}, \text{F1 range}){=}0.328$ shows that instances with more LP spread also exhibit more F1 variability. LP carries information, but compression limits its exploitation.
A free-form NER decoding ablation (Appendix~\ref{app:freeform_ablation}) confirms this compression persists without constrained decoding.

\paragraph{$\varepsilon$ sensitivity sweep (post-hoc exploratory).}
Table~\ref{tab:epsilon_sweep} sweeps the tied/non-tied threshold from $\varepsilon{=}0.01$ to $0.20$ as a post-hoc exploratory analysis characterizing how selection gain varies with compression severity.
The tied fraction increases smoothly with no cliff or discontinuity, and the selection gain $\Delta$ grows monotonically as the non-tied subset narrows to progressively higher-spread instances.
$\varepsilon{=}0.05$ yields the strongest statistical separation ($p{=}0.004$); at $\varepsilon{\geq}0.15$, the non-tied subset becomes too small for reliable inference ($n{\leq}36$).

\begin{table}[t]
\centering
\small
\caption{$\varepsilon$ sensitivity sweep for LP selection on SciERC NER ($N{=}8$, seed${=}42$, $n{=}529$). ``Tied\%'' = fraction with within-instance LP range ${<}\varepsilon$. $\Delta$: selection F1 gain over single-sample on the non-tied subset. The tied fraction increases smoothly; $\varepsilon{=}0.05$ yields the strongest statistical separation.}
\label{tab:epsilon_sweep}
\begin{tabular}{ccccc}
\toprule
$\varepsilon$ & Tied\% & Non-tied $n$ & $\Delta$ (pp) & $p$ \\
\midrule
0.01 & 1.1\% & 523 & $+$0.58 & .21 \\
0.02 & 13.6\% & 457 & $+$0.83 & .15 \\
0.03 & 31.8\% & 361 & $+$1.31 & .079 \\
\textbf{0.05} & \textbf{60.3\%} & \textbf{210} & \textbf{$+$2.86} & \textbf{.004} \\
0.10 & 83.9\% & 85 & $+$3.69 & .034 \\
0.15 & 93.2\% & 36 & $+$4.15 & .11 \\
0.20 & 96.4\% & 19 & $+$4.59 & .13 \\
\bottomrule
\end{tabular}
\end{table}

\paragraph{$\varepsilon$ robustness of selection F1.}\label{app:epsilon_robustness}
The preceding analysis stratifies LP by a tied/non-tied threshold; consistency-based signals (SJ, FK, VC, EM) also use $\varepsilon$ to define MBR-style proximity in selection.
Table~\ref{tab:epsilon_robustness} reports the maximum selection F1 variation across $\varepsilon \in [0, 0.2]$ for all signals and configurations.
At $\varepsilon{=}0.05$, consistency signals change by at most 1.09~pp (mean 0.21~pp), with EM completely $\varepsilon$-invariant (integer scoring); LP is the most sensitive (up to 3.65~pp across the full range), as expected.
Spearman $\rho$, the primary quality estimation metric throughout, is a continuous rank correlation that does not depend on any discretization threshold; the $\varepsilon$ parameter affects only the selection analysis.

\begin{table}[t]
\centering
\small
\caption{Maximum selection F1 variation (pp) across $\varepsilon \in [0, 0.2]$. At $\varepsilon{=}0.05$ (used throughout), consistency signals change by ${\leq}1.09$~pp; EM is invariant (integer scoring). LP is most sensitive as $\varepsilon$ corresponds to its tied/non-tied partition. SciERC: $n{=}529$; CoNLL: $n{=}2{,}756$.}
\label{tab:epsilon_robustness}
\begin{tabular}{lcccc}
\toprule
& \multicolumn{2}{c}{\textbf{SciERC}} & \multicolumn{2}{c}{\textbf{CoNLL}} \\
\cmidrule(lr){2-3} \cmidrule(lr){4-5}
Signal & Qwen & LLaMA & Qwen & LLaMA \\
\midrule
SJ & 0.92 & 1.45 & 0.99 & 0.18 \\
FK & 1.03 & 0.93 & 0.60 & 0.15 \\
VC & 2.33 & 0.53 & 0.72 & 0.18 \\
EM & 0.00 & 0.00 & 0.00 & 0.00 \\
LP & 3.65 & 2.58 & 3.00 & 0.31 \\
\bottomrule
\end{tabular}
\end{table}

\paragraph{Non-tied subset analysis.}
Table~\ref{tab:nontied_main} presents the non-tied subset analysis referenced in \S\ref{sec:discussion}.

\begin{table}[t]
\centering
\small
\caption{Non-tied subset analysis (SciERC, Qwen3-8B, $N{=}8$, $T{=}1.0$). Instances split at LP range $\varepsilon{=}0.05$. Within-instance $\rho$ and selection gain both concentrate in the non-tied subset. Full $\varepsilon$ sweep in Table~\ref{tab:epsilon_sweep}.}
\label{tab:nontied_main}
\resizebox{\columnwidth}{!}{
\begin{tabular}{lccc}
\toprule
Subset & $n$ & $\rho$(LP,F1) & LP $\Delta$ vs.\ greedy \\
\midrule
Tied (range ${<}0.05$) & 319 & 0.08 & ${\approx}0$ \\
Non-tied (range ${\geq}0.05$) & 210 & 0.39 & $+$2.9~pp ($p{=}0.004$) \\
\midrule
All instances & 529 & 0.248 & $+$0.2~pp ($p{=}0.29$) \\
\bottomrule
\end{tabular}
}
\end{table}

\paragraph{Within-instance LP-quality alignment.}\label{app:within_instance_rho}

Table~\ref{tab:within_rho} reports median within-instance $\rho(\text{LP}, \text{F1})$ across datasets, computed over non-degenerate instances (those with ${\geq}2$ distinct F1 values across $N{=}8$ samples).
Few-NERD exhibits substantially stronger alignment ($\rho{=}0.667$) than SciERC ($\rho{=}0.262$), consistent with LP selection succeeding only on Few-NERD.
CoNLL alignment is moderate ($\rho{\approx}0.51$), but the valid-instance fraction is low (38\% for Qwen, 11\% for LLaMA) because degeneracy eliminates most headroom before alignment can operate.
LP-quality alignment is the primary predictor of selection direction; degeneracy constrains headroom but does not independently predict whether LP selection gains materialize (\S\ref{sec:discussion}).

\begin{table}[h]
\centering
\small
\caption{Within-instance LP-quality alignment across datasets ($N{=}8$, $T{=}1.0$). $\rho$: median Spearman correlation between LP and F1 within each instance, over non-degenerate instances ($n_{\text{valid}}$: ${\geq}2$ distinct F1 values). Counts are from the primary seed (seed 42); text references use 3-seed mean degeneracy rates.}
\label{tab:within_rho}
\resizebox{\columnwidth}{!}{%
\begin{tabular}{@{}llccc@{}}
\toprule
Dataset & Model & Med.\ $\rho$ & $n_{\text{valid}}$ / $n_{\text{total}}$ & Valid\% \\
\midrule
SciERC & Qwen3-8B & 0.262 & 435 / 529 & 82\% \\
Few-NERD & Qwen3-8B & 0.667 & 28{,}055 / 32{,}565 & 86\% \\
CoNLL & Qwen3-8B & 0.506 & 1{,}124 / 2{,}756 & 41\% \\
CoNLL & LLaMA-3.1-8B & 0.512 & 298 / 2{,}756 & 11\% \\
\bottomrule
\end{tabular}}
\end{table}

\clearpage
\paragraph{Few-NERD LP compression threshold sensitivity.}\label{app:fewnerd_epsilon}

Tables~\ref{tab:fewnerd_epsilon} and \ref{tab:fewnerd_lp_bucket} extend the $\varepsilon$ analysis to Few-NERD ($n{=}32{,}565$, Qwen3-8B, $N{=}8$, seed${=}42$).

\begin{table}[ht!]
\centering
\small
\caption{$\varepsilon$ threshold vs.\ LP compression on Few-NERD ($N{=}8$, seed${=}42$). Compressed\%: instances with LP range ${<}\varepsilon$. $\rho$: Spearman correlation between LP and oracle F1. $\varepsilon{=}0.05$ (bold) is the elbow.}
\label{tab:fewnerd_epsilon}
\begin{tabular}{cccc}
\toprule
$\varepsilon$ & Compressed\% & $\rho$(non-comp.) & $\rho$(all) \\
\midrule
0.01 & 9.3 & .682 & .633 \\
0.03 & 19.9 & .695 & .633 \\
\textbf{0.05} & \textbf{34.5} & \textbf{.710} & \textbf{.633} \\
0.10 & 60.2 & .733 & .633 \\
0.15 & 75.9 & .723 & .633 \\
0.20 & 92.5 & .720 & .633 \\
\bottomrule
\end{tabular}
\end{table}

\begin{table}[ht!]
\centering
\small
\caption{LP-best selection advantage over MV by LP range bucket on Few-NERD ($N{=}8$, seed${=}42$). Instances with LP range ${\geq}0.05$ show $3$--$6{\times}$ larger advantage.}
\label{tab:fewnerd_lp_bucket}
\begin{tabular}{lccc}
\toprule
LP range & $n$ & MV F1 & LP$-$MV (pp) \\
\midrule
$[0, 0.01)$ & 3{,}018 & .818 & $+$0.8 \\
$[0.01, 0.03)$ & 3{,}457 & .812 & $+$1.5 \\
$[0.03, 0.05)$ & 4{,}744 & .786 & $+$2.3 \\
\midrule
$[0.05, 0.10)$ & 8{,}374 & .740 & $+$4.3 \\
$[0.10, 0.20)$ & 10{,}537 & .701 & $+$5.8 \\
$[0.20, \infty)$ & 2{,}435 & .771 & $+$3.6 \\
\bottomrule
\end{tabular}
\end{table}

$\varepsilon{=}0.05$ sits at the elbow: compressed instances increase from 19.9\% ($\varepsilon{=}0.03$) to 34.5\%, while non-compressed $\rho$ plateaus at $\varepsilon{\geq}0.10$ ($0.72$--$0.73$).
The LP-best advantage over MV shows a clear transition: compressed instances (LP range${<}0.05$) gain $+0.8$ to $+2.3$\,pp, while non-compressed instances gain $+3.6$ to $+5.8$\,pp (Table~\ref{tab:fewnerd_lp_bucket}).

\subsection{Token-Level Signal Comparison}\label{app:token_signal}

We compared mean log-probability (length-normalized, used throughout) with sequence-level log-probability (cumulative sum of token log-probs).
Instance-level Spearman $\rho$ between the two rankings averages 0.85, with rankings diverging for 10.6\% of instances where token count coefficient of variation exceeds 0.10.
On these disagreement instances, mean LP selection achieves F1${=}0.643$ vs.\ cumulative LP F1${=}0.219$ ($+42.4$\,pp).
Overall, mean LP outperforms cumulative LP by $+4.6$ F1 points (Wilcoxon $p{<}0.001$).
The divergence is driven by structured output length variation: sequences with more extracted entities are longer but not necessarily higher quality, and length normalization prevents this bias.

\section{Consistency-Stratified Selection Analysis}\label{app:vc_quartile}

Table~\ref{tab:vc_quartile} stratifies SciERC instances by voting confidence (VC) quartiles and reports selection $\Delta$ (gain over single-sample) for LP and SJ.

\begin{table}[t]
\centering
\small
\caption{Selection $\Delta$ (pp over single-sample) stratified by VC quartile on SciERC ($N{=}8$, Qwen3-8B). Oracle headroom collapses at Q4; LP rankings are unreliable at Q1--Q2. Neither quartile supports consistent selection gains.}
\label{tab:vc_quartile}
\begin{tabular}{lcccc}
\toprule
Quartile & Mean VC & LP $\Delta$ & SJ $\Delta$ & Oracle (pp) \\
\midrule
\multicolumn{5}{l}{\textit{NER}} \\
Q1 (low) & 0.30 & $+$1.55 & $+$0.54 & 16.85 \\
Q2 & 0.51 & $-$0.65 & $-$1.89 & 14.35 \\
Q3 & 0.68 & $+$1.97 & $+$1.84 & 16.83 \\
Q4 (high) & 0.92 & $-$0.41 & $-$0.04 & 4.78 \\
\midrule
\multicolumn{5}{l}{\textit{$^\dagger$RE}} \\
Q1 (low) & --- & $-$3.56 & --- & --- \\
Q4 (high) & --- & $+$4.01 & --- & --- \\
\bottomrule
\end{tabular}
\end{table}

The NER pattern shows two complementary failure modes rather than a monotonic relationship between consistency and selection gain:
\begin{itemize}
\item \textbf{Q4 ceiling effect.} High-consistency instances (mean VC${=}0.92$) have oracle headroom of only 4.78~pp. Samples have already converged, leaving little room for any signal to improve selection.
\item \textbf{Q1--Q2 noise.} Low-consistency instances retain large headroom (14--17~pp) but LP rankings are unreliable, and LP rankings produce inconsistent $\Delta$ signs across quartiles.
\end{itemize}

On $^\dagger$RE, the pattern differs: Q4 LP achieves $+4.01$~pp; LP may recover selection ability when high consistency coincides with structured output diversity. This is a single-dataset observation and should be interpreted cautiously.

On CoNLL NER, Q3--Q4 instances have VC${=}1.0$ (all samples identical), and $\Delta{=}0$ universally; this ceiling effect is even more extreme than SciERC Q4.

\section{Exploratory Extension to Relation Extraction}\label{app:re_extension}

We extend the NER evaluation to relation extraction (RE) on SciERC as an exploratory single-benchmark analysis.
RE produces sets of (head, relation, tail) triples with partial overlap; 40\% of RE instances receive zero F1, and entity span errors account for 55.5\% of RE false negatives, a prerequisite failure that consistency signals cannot detect (\S\ref{app:additional}).
All RE results use Qwen3-8B unless noted; single-seed ($\text{seed}{=}42$) for $N{=}8$ analyses, 3-seed for $N{=}16$.

\paragraph{Diagnostic framework.}
Applying the diagnostic framework to RE ($N{=}8$): instance degeneracy is $10.1$\%, comparable to SciERC NER ($13.4$\%), yet MV $\Delta{=}{-}5.90$\,pp---far worse than any NER dataset.
The driver is RE's $40$\% zero-F1 rate (vs.\ ${\leq}5.4$\% in NER): diverse but \textit{incorrect} samples produce noise that MV amplifies.
LP-best selection also fails ($\Delta{=}{-}0.94$\,pp), despite pooled LP-F1 $\rho{=}0.350$---above SciERC NER ($0.312$).
Oracle headroom is $+18.1$\,pp (greedy $0.404$, oracle $0.585$), the largest of any configuration, confirming that correct answers exist among samples but no signal recovers them.

\paragraph{Quality estimation.}
RE quality estimation is reported in Table~\ref{tab:multiseed} (3-seed means).
The signal ranking differs from NER: VC leads ($\rho{=}0.309{\pm}0.024$), but LP rises to \#2 ($\rho{=}0.297{\pm}0.022$), above FK, SJ, and EM, a reversal reflecting RE's higher sensitivity to per-sample scoring.
RE miscalibration is more severe than NER (ECE 0.109--0.339; Figure~\ref{fig:reliability}).

\paragraph{Selection.}
Table~\ref{tab:re_cond_selection} presents RE selection at $N{=}8$: no signal outperforms single-sample (0.404) on the full set ($n{=}397$) or the non-zero subset ($n{=}238$), with 18.1~pp of unrealized oracle headroom.
At $N{=}16$ (Qwen3-4B, Table~\ref{tab:re_n16_selection}), LP is the sole positive signal ($\Delta{=}{+}0.7$~pp, $p{=}0.058$), consistent with LP being the strongest RE selection signal across scales.

\paragraph{Scale ablation.}
At 4B scale (Qwen3-4B, $N{=}8$, $\text{seed}{=}42$), RE quality estimation shows FK $\rho{=}0.359$ and LP $\rho{=}0.351$, both above their 8B counterparts (FK 0.275, LP 0.266), while SJ declines ($0.267$ vs.\ $0.250$).
Selection remains below single-sample on RE ($\Delta$ up to $-3.4$~pp) despite 17.1~pp of oracle headroom.

\paragraph{Free-form decoding.}
RE LP selection consistently reverses direction under free-form decoding ($+1.8{\pm}0.1$~pp across 4 seeds vs.\ $-0.9$~pp under constrained decoding at 8B; \S\ref{app:freeform_ablation}); constrained decoding's fixed template likely obscures LP signal in RE (\S\ref{app:mbr_discussion}).

\subsection{RE Conditional Selection}\label{app:re_conditional_selection}

Table~\ref{tab:re_cond_selection} extends the selection analysis to the RE non-zero F1 subset.
Of $n{=}397$ RE instances, 238 (60\%) have non-zero F1 and 159 (40\%) have zero F1.
All five signals underperform single-sample on both subsets.
On the non-zero subset, selection performance is \textit{worse} than on the full set (e.g., VC: $-10.0$~pp on non-zero vs.\ $-4.5$~pp on full set), because single-sample is already strong on non-zero instances (F1${=}0.674$) and signals actively select inferior samples.
Oracle headroom exists (0.757 vs.\ 0.674) but no signal exploits it, ruling out the 40\% zero-F1 floor effect as the explanation for RE selection failure.

\begin{table}[t]
\centering
\small
\caption{RE selection F1 (macro, instance-averaged) on SciERC ($N{=}8$, $T{=}1.0$). Full set ($n{=}397$) and non-zero F1 subset ($n{=}238$). No signal outperforms greedy on either subset.}
\label{tab:re_cond_selection}
\begin{tabular}{lcc}
\toprule
Method & Full ($n{=}397$) & Non-zero ($n{=}238$) \\
\midrule
Oracle & 0.585 & 0.757 \\
\midrule
LP & 0.395 & 0.600 \\
SJ & 0.386 & 0.604 \\
EM & 0.372 & 0.583 \\
VC & 0.359 & 0.574 \\
\midrule
Greedy & 0.404 & 0.674 \\
\bottomrule
\end{tabular}
\end{table}

\subsection{RE $N{=}16$ Selection Analysis}\label{app:re_n16_selection}

Table~\ref{tab:re_n16_selection} extends the selection analysis to RE at $N{=}16$ with 3-seed statistics.
All four consistency signals underperform single-sample, with EM showing the largest deficit ($\Delta{=}{-}1.5$~pp, Fisher $p{>}0.99$).
LP is the only positive signal ($\Delta{=}{+}0.7$~pp, Fisher $p{=}0.058$), driven primarily by seed~456 ($\Delta{=}{+}1.5$~pp, $p{=}0.026$); the remaining seeds show smaller effects ($+0.1$ and $+0.6$~pp, $p{>}0.20$).
Even the most favorable seed captures only ${\sim}9\%$ of the 17.1~pp oracle headroom.
This cross-task consistency (LP as the only viable selection signal on both NER and RE) reinforces the persistent limitation of instance-level signals for candidate ranking across tested regimes.

\begin{table}[t]
\centering
\small
\caption{RE best-of-$N$ selection F1 (macro, instance-averaged) on SciERC (Qwen3-4B, $N{=}16$, $T{=}1.0$, 3-seed mean${\pm}\sigma$, $n{=}551$). Fisher combined $p$ vs.\ greedy. LP is the only signal with positive $\Delta$. $\sigma$ is population standard deviation over 3 seeds.}
\label{tab:re_n16_selection}
\begin{tabular}{lccc}
\toprule
Method & F1 & $\Delta$ (pp) & $p$(Fisher) \\
\midrule
Oracle & $.463{\pm}.006$ & $+$17.1 & --- \\
\midrule
LP & $\mathbf{.299}{\pm}.007$ & $+$0.7 & .058 \\
VC & $.284{\pm}.002$ & $-$0.8 & .973 \\
FK & $.283{\pm}.005$ & $-$0.9 & .985 \\
SJ & $.281{\pm}.005$ & $-$1.0 & .993 \\
EM & $.277{\pm}.007$ & $-$1.5 & .999 \\
\midrule
Greedy & $.291{\pm}.000$ & --- & --- \\
Random & $.262{\pm}.001$ & $-$2.9 & --- \\
\bottomrule
\end{tabular}
\end{table}

\section{Theoretical Results}\label{app:mbr_discussion}

\subsection{MBR Framework and Proof}\label{app:mbr_proof}

Entity-level construction (\S\ref{sec:prescriptive}) is a special case of minimum Bayes risk (MBR) decoding \cite{kumar2004minimum, eikema2020map}, extended from candidate selection ($\arg\max_{y \in \{y_1,\ldots,y_N\}}$) to construction over all subsets of $E_\cup{=}\bigcup_i y_i$.

\begin{proposition}[Entity construction as MBR for balanced set-overlap]\label{prop:construction_mbr}
Under entity-wise independence with utility $u(\hat{y}, y') = |\hat{y} \cap y'| - \alpha|\hat{y} \setminus y'|$ ($\alpha \geq 1$), the MBR-optimal constructed set is
\begin{equation}\label{eq:mbr_construction}
\hat{y}_{\textup{MBR}} \!=\! \bigl\{e \!{\in}\! E_\cup \!:\! \hat{p}(e) \!>\! \tfrac{\alpha}{\alpha{+}1}\bigr\},\; \hat{p}(e) \!=\! \tfrac{1}{N} \textstyle\sum\nolimits_{i} \mathbf{1}[e \!{\in}\! y_i]
\end{equation}
The special case $\alpha{=}1$ yields $\hat{p}(e) > \tfrac{1}{2}$, the majority-vote threshold.
\end{proposition}

\noindent\textbf{Remark.}
Set-level F1 is non-linear and admits no closed-form threshold rule \cite{dembczynski2011exact}. We use balanced linear utility ($\alpha{=}1$ throughout); balanced set-overlap and F1 induce similar optimal thresholds for typical NER output sizes ($|E_\cup| \leq 30$).

\paragraph{Proof.}
We show that under entity-wise independence, the MBR-optimal constructed set with utility $u(\hat{y}, y') = |\hat{y} \cap y'| - \alpha|\hat{y} \setminus y'|$ ($\alpha \geq 1$) includes entity $e$ iff $\hat{p}(e) > \alpha/(\alpha+1)$.
Consider constructing $\hat{y} \subseteq E_\cup$ to maximize expected utility against the empirical sample distribution.
Under entity-wise independence, each entity inclusion is an independent binary decision.
For entity $e$ with empirical frequency $\hat{p}(e) = \frac{1}{N} \sum_i \mathbf{1}[e \in y_i]$, including $e$ contributes $+1$ when $e \in y_i$ (probability $\hat{p}(e)$) and $-\alpha$ when $e \notin y_i$ (probability $1{-}\hat{p}(e)$).
The expected marginal contribution is $\hat{p}(e) - \alpha(1{-}\hat{p}(e)) = (1{+}\alpha)\hat{p}(e) - \alpha$, which is positive iff $\hat{p}(e) > \alpha/(\alpha+1)$.
The special case $\alpha{=}1$ yields the majority-vote threshold $\hat{p}(e) > 1/2$ used throughout.
With F1 utility, the optimal threshold depends on the full probability vector over $E_\cup$ \cite{dembczynski2011exact}; the constant $1/2$ threshold is a close approximation when $|E_\cup|$ is small (typically 5--15 in our experiments).
The independence assumption is approximate; we verify empirically that strict majority vote fails to improve over greedy (\S\ref{sec:prescriptive}; Table~\ref{tab:alpha_crit}), consistent with marginal-entity precision falling below $\alpha_{\textup{crit}}$.

\subsection{Connection to MBR Decoding}\label{app:mbr_connection}

Our selection strategies instantiate minimum Bayes risk (MBR) decoding \cite{eikema2022sampling, muller2021understanding}: each selects $\hat{y} = \arg\max_{y_i} \sum_{j \neq i} U(y_i, y_j)$ with utility $U$ defined by SJ, EM, or entity overlap.
Despite MBR's success in machine translation \cite{freitag2022high} and its extension to summarization, code generation, and captioning \cite{jinnai2024mbr}, all variants fail here: despite positive within-instance utility--F1 correlation, selection gains are limited by instance degeneracy, the core of the correlation-selection gap.
\citet{eikema2025structured} study structure-conditional MBR decoding and find that standard utility functions degrade on structured outputs, consistent with our finding that consensus-proximity selection fails in IE despite its success in machine translation.

\paragraph{Task-specific decoding interactions.}
The free-form ablation shows divergent behavior across tasks.
On RE, LP selection flips from $-0.68$~pp (constrained) to $+1.69$~pp (free-form; preliminary, single RE benchmark), a reversal that spans all 7 relation types and is attributable to selective variance reallocation ($p{<}0.0001$).
The win/loss ratio shifts from 36:59 (constrained) to 45:45 (free-form), with the gains concentrated on relation-dense instances (mean 3.3 gold relations).
This interaction does not generalize to NER, where selection remains negligible under both regimes.
This ablation uses a single seed (seed${=}42$); whether the RE reversal replicates across seeds and generalizes to other RE benchmarks warrants further investigation.
The non-tied subset ($n{=}210$) is partially reshuffled under free-form decoding: Jaccard similarity between the constrained and free-form non-tied sets is 0.472: 87 instances newly gain LP discrimination while others lose it, and the newly non-tied subset shows a $3{\times}$ increase in within-instance LP range for the newly non-tied subset.
Despite this reshuffling, aggregate NER selection gains remain negligible. Score compression is not a fixed property of specific instances but rather a pervasive regime that shifts across instances without producing net selection improvement.
Free-form decoding also slightly improves quality estimation (NER VC $\rho$: 0.379$\to$0.409), though this does not translate to selection gains.

\subsection{Degeneracy Decomposition}\label{app:bound_proof}

\begin{definition}[Instance-Level Degeneracy]
Given model $\mathcal{M}$, input $x$, and $N$ i.i.d.\ samples $\{y_1, \ldots, y_N\}$, instance $x$ is \emph{degenerate} if $y_1 = \cdots = y_N$. The degeneracy rate is $D = |\mathcal{D}|^{-1} \sum_{x} \mathbb{1}[y_1(x) = \cdots = y_N(x)]$.
\end{definition}

\begin{definition}[Non-Degenerate Oracle Gap]
$G_{\text{nd}} = \mathbb{E}[F_1^{\text{oracle}}(x) - F_1^{\text{greedy}}(x) \mid x \in \mathcal{D}_{\neg D}]$, where $\mathcal{D}_{\neg D}$ is the set of non-degenerate instances.
\end{definition}

\subsection{Proof of the Degeneracy Decomposition}

Partition $\mathcal{D}$ into degenerate ($\mathcal{D}_D$, fraction $D$) and non-degenerate ($\mathcal{D}_{\neg D}$, fraction $1{-}D$). By total expectation:
\begin{equation}\label{eq:bound_decomp}
\mathbb{E}[\Delta F_1(A)] = D \cdot \mathbb{E}[\Delta F_1 \mid \mathcal{D}_D] + (1{-}D) \cdot \mathbb{E}[\Delta F_1 \mid \mathcal{D}_{\neg D}].
\end{equation}
For degenerate instances, all $N$ samples are identical and equal the greedy output, so $\Delta F_1 = 0$. For non-degenerate instances, $\Delta F_1(A,x) \leq F_1^{\text{oracle}}(x) - F_1^{\text{greedy}}(x)$, giving $\mathbb{E}[\Delta F_1 \mid \mathcal{D}_{\neg D}] \leq G_{\text{nd}}$. Substituting: $\mathbb{E}[\Delta F_1(A)] \leq (1{-}D) \cdot G_{\text{nd}}$.

For signal-based methods, define alignment $\alpha = \mathbb{E}[\Delta F_1(A_s) \mid \mathcal{D}_{\neg D}] / G_{\text{nd}} \in [0,1]$. Direct substitution yields $\mathbb{E}[\Delta F_1(A_s)] = (1{-}D) \cdot G_{\text{nd}} \cdot \alpha$.

\begin{corollary}[Critical Degeneracy Threshold]\label{cor:theta_crit}
For SC gain to exceed significance threshold $\varepsilon$, the degeneracy rate must satisfy $D < 1 - \varepsilon / (G_{\text{nd}} \cdot \alpha) \stackrel{\text{def}}{=} \theta_{\text{crit}}$.
With $\varepsilon{=}0.3$\,pp, $G_{\text{nd}}{=}14.8$\,pp, $\alpha{=}0.068$: $\theta_{\text{crit}} = 1 - 0.3/(14.8 \times 0.068) = 0.70$.
\end{corollary}

\paragraph{Numerical verification.}
Table~\ref{tab:bound_verification} verifies the bound across 16 configurations. All positive-gain configurations have $D < \theta_{\text{crit}}$ and actual/bound ratios of 3--12\% (Figure~\ref{fig:bound}). The bound is loose because $\alpha$ (signal alignment) is the bottleneck: the best signal recovers only 6.8\% of the non-degenerate oracle gap on average.


\begin{table}[t]
\centering
\small
\caption{Bound verification across representative configurations ($N{=}8$). $G_{\text{nd}}$: non-degenerate oracle gap. $\alpha$: signal alignment. Bound: $(1{-}D) \cdot G_{\text{nd}} \cdot \alpha$. Ratio: actual$/$weak bound.}
\label{tab:bound_verification}
\begin{tabular}{@{}lccccr@{}}
\toprule
Config & $D$ & $G_{\text{nd}}$ & $\alpha$ & $\Delta$ & Ratio \\
\midrule
Q3-8B FT FN s123 & .131 & 14.8 & .121 & +1.56 & 12.1\% \\
Q3-8B FT FN s456 & .120 & 14.5 & .108 & +1.37 & 10.8\% \\
Q3-8B FT Sci s789 & .214 & 17.0 & .084 & +1.12 & 8.4\% \\
72B ZS Sci & .210 & 18.9 & .102 & +1.53 & 10.2\% \\
Q3-8B FT CoN s123 & .630 & 10.9 & .000 & $-$0.42 & --- \\
72B ZS CoN & .677 & 16.4 & .000 & $-$0.37 & --- \\
72B FS CoN & .700 & 17.0 & .047 & +0.24 & 4.7\% \\
\bottomrule
\end{tabular}
\end{table}

\subsection{Entity Selection Bound: Proof and Validation}\label{app:alpha_crit}

\begin{proposition}[Entity Selection Diagnostic]\label{prop:entity_selection_bound}
Let $\mathcal{M}^+ {=} \hat{y}_{\textup{MV}} \setminus \hat{y}_{\textup{gr}}$ be entities added by majority vote (precision $\pi_{\mathcal{M}^+}$) and $\delta$ the instance degeneracy rate.
\begin{enumerate}[label=\textup{(\alph*)},nosep,leftmargin=*]
\item On degenerate instances, $\hat{y}_{\textup{MV}} {=} \hat{y}_{\textup{gr}}$ and $\Delta F_1 {=} 0$.
\item $\mathbb{E}[\Delta F_1] = (1{-}\delta) \cdot \mathbb{E}[\Delta F_1 \mid \textup{non-deg.}]$.
\item $\Delta F_1 > 0$ only if $\pi_{\mathcal{M}^+} > \alpha_{\textup{crit}} \triangleq F_1^{\textup{gr}}/2$.
\end{enumerate}
\end{proposition}

\begin{proof}[Proof of Proposition~\ref{prop:entity_selection_bound}]
\textbf{Part (a).} On a degenerate instance, all $N$ samples produce the same entity set: $\textup{entities}(y_1) = \cdots = \textup{entities}(y_N)$. For every entity $e \in E_\cup$, the agreement count satisfies $k(e) = N$ (if $e$ appears) or $k(e) = 0$ (if not). Hence $k(e)/N \in \{0,1\}$, and the majority-vote threshold $k(e)/N \geq 1/2$ selects exactly the entities present in any single sample---identical to the greedy output. Thus $\hat{y}_{\textup{MV}} = \hat{y}_{\textup{gr}}$ and $\Delta F_1 = 0$.

\textbf{Part (b).} Partition instances into degenerate (fraction $\delta$) and non-degenerate (fraction $1{-}\delta$). By the law of total expectation:
\[
\mathbb{E}[\Delta F_1] = \delta \cdot \underbrace{\mathbb{E}[\Delta F_1 \mid \textup{deg.}]}_{= 0 \;\textup{by (a)}} + (1{-}\delta) \cdot \mathbb{E}[\Delta F_1 \mid \textup{non-deg.}].
\]

\textbf{Part (c).} Assume $\mathcal{M}^- = \hat{y}_{\textup{gr}} \setminus \hat{y}_{\textup{MV}} \approx \varnothing$ (greedy entities typically have agreement $k(e)/N \approx 1$). Let $m = |\mathcal{M}^+|$ and $\pi = |\mathcal{M}^+ \cap G|/|\mathcal{M}^+|$ denote the number and precision of entities added by MV. Define $S = |\hat{y}_{\textup{gr}}| + |G|$. Then:
\begin{align*}
\Delta F_1 &= \frac{2(\textup{TP}_{\textup{gr}} + m\pi)}{S + m} - \frac{2\,\textup{TP}_{\textup{gr}}}{S} \\
&= \frac{2m\bigl(\pi - \textup{TP}_{\textup{gr}}/S\bigr)}{S + m}
= \frac{2m\bigl(\pi - F_1^{\textup{gr}}/2\bigr)}{S + m},
\end{align*}
which is positive if and only if $\pi > F_1^{\textup{gr}}/2 \triangleq \alpha_{\textup{crit}}$.

When $\mathcal{M}^-$ is non-empty, the condition generalizes: the net precision across the full symmetric difference $\hat{y}_{\textup{MV}} \triangle \hat{y}_{\textup{gr}}$ must exceed a corresponding threshold. The addition-only case covers the dominant empirical regime, as greedy entities typically achieve near-unanimous agreement.
\end{proof}

\paragraph{Effective precision threshold.}
From parts (b) and (c), MV gain is attenuated by $(1{-}\delta)$ and requires marginal-entity precision $> \alpha_{\textup{crit}}$.
Equivalently, defining $\alpha_{\textup{crit,eff}} {=} \alpha_{\textup{crit}} / (1{-}\delta)$, MV gain requires marginal-entity precision $> \alpha_{\textup{crit,eff}}$; when $\alpha_{\textup{crit,eff}} > 1$, the required precision exceeds 1.0 and MV gain is structurally impossible regardless of entity quality.

\paragraph{Empirical validation.}
Table~\ref{tab:alpha_crit} reports MV performance under strict majority (count $> N/2$) using the MF4v2 pipeline (3-seed means).
MV fails to improve F1 on any dataset: Few-NERD $\Delta{=}{-}2.74$\,pp, SciERC $-0.74$\,pp, CoNLL $-0.04$\,pp.
We define the effective precision threshold $\alpha_{\textup{crit,eff}} {=} \alpha_{\textup{crit}} / (1{-}\delta)$; when $\alpha_{\textup{crit,eff}} > 1$, the required marginal-entity precision exceeds 1.0, making MV gain structurally impossible.
CoNLL ($\alpha_{\textup{crit,eff}}{=}1.148$) confirms this: $\Delta{=}{-}0.04{\pm}0.15$\,pp is indistinguishable from zero.
Few-NERD and SciERC ($\alpha_{\textup{crit,eff}} {<} 1$) show significantly negative $\Delta$, indicating that even where structural headroom exists, marginal-entity precision is insufficient.
SciERC's lower $\alpha_{\textup{crit}}$ ($0.334$, vs.\ $0.396$ for Few-NERD and $0.456$ for CoNLL) is consistent with it being the only dataset where a trained entity verifier (AUC${=}0.923$) achieves positive gain ($+1.28$\,pp).

\begin{table}[h]
\centering\small
\caption{MV performance under strict majority (count $> N/2$; Qwen3-8B FT, $N{=}8$, $T{=}1.0$, 3-seed mean${\pm}\sigma$). $\alpha_{\textup{crit}} = F_1^{\textup{gr}}/2$; $\alpha_{\textup{crit,eff}} = \alpha_{\textup{crit}} / (1{-}\delta)$. When $\alpha_{\textup{crit,eff}} > 1$, MV gain is structurally impossible.}
\label{tab:alpha_crit}
\resizebox{\columnwidth}{!}{%
\begin{tabular}{@{}lccccccccc@{}}
\toprule
Dataset & Seeds & Greedy & MV & Oracle & $\Delta_{\textup{SC}}$ (pp) & $\delta$ & $\alpha_{\textup{crit}}$ & $\alpha_{\textup{crit,eff}}$ & Prediction \\
\midrule
Few-NERD & 3 & .791 & .764 & .885 & $-2.74{\pm}0.54$ & 8.8\% & 0.396 & 0.434 & $\Delta < 0$ \\
CoNLL & 3 & .912 & .912 & .950 & $-0.04{\pm}0.15$ & 60.3\%$^*$ & 0.456 & 1.148 & $\Delta \approx 0$ \\
SciERC & 3 & .667 & .660 & .781 & $-0.74{\pm}0.46$ & 12.2\% & 0.334 & 0.380 & $\Delta < 0$ \\
\bottomrule
\end{tabular}}

\noindent{\footnotesize $^*$CoNLL $\delta{=}60.3$\% (full test set, 3-seed) vs.\ $61.7$\% in Table~\ref{tab:fewnerd} (gold-filtered); the difference reflects filtering. Tab.~\ref{tab:pretrained_summary} reports $57$--$60$\% (seed range, finetuned).}
\end{table}

\section{Free-Form Decoding Ablation}\label{app:freeform_ablation}

Table~\ref{tab:freeform_ablation} compares constrained and free-form decoding across four dimensions on SciERC ($N{=}8$, $T{=}1.0$, Qwen3-8B; seed${=}42$ shown, 4-seed RE statistics below).
Free-form decoding removes grammar constraints entirely; outputs are parsed post-hoc (parse rate 99.95\%, 4406/4408).
LP score distributions are nearly identical across regimes (Panel~A: within-instance std ratio $0.97{\times}$), and NER selection gains remain negligible under both regimes (Panel~B).
RE LP selection consistently reverses direction under free-form decoding across 4 seeds: $+1.84{\pm}0.12$~pp (seeds 42/123/456/789: $+1.69$, $+1.76$, $+2.00$, $+1.91$~pp) vs.\ $-0.68$~pp constrained.
The constrained reference ($-0.68$~pp) uses the ablation-matched baseline; the primary experiment (Table~\ref{tab:selection}) yields $-0.9$~pp on the same gold-filtered set.
Quality estimation improves modestly for most signals under free-form decoding (Panel~C), with the largest gain for FK ($+0.053$ NER, $+0.039$ RE).

\begin{table*}[t]
\centering
\caption{Comparison of constrained vs.\ free-form decoding on SciERC (Qwen3-8B, $N{=}8$, $T{=}1.0$, seed${=}42$; RE selection gap replicates across 4 seeds: $+1.84{\pm}0.12$~pp).}
\label{tab:freeform_ablation}

\vspace{2pt}
{\small\textsc{Panel A: Logprob Distribution}}
\vspace{2pt}

\begin{tabular}{lccc}
\toprule
Metric & Constrained & Free-form & $\Delta$ \\
\midrule
Mean LP std        & 0.0206 & 0.0199 & $-$0.0007 \\
Mean LP range      & 0.0623 & 0.0599 & $-$0.0024 \\
Tied fraction ($\varepsilon{=}0.05$) & 59.7\% & 57.7\% & $-$2.0\,pp \\
\bottomrule
\end{tabular}

\vspace{8pt}

{\small\textsc{Panel B: Selection F1}}
\vspace{2pt}

\begin{tabular}{llccc}
\toprule
Task & Method & Constrained & Free-form & $\Delta$ \\
\midrule
\multirow{3}{*}{NER}
  & Greedy        & 0.6182 & 0.6204 & $+$0.0022 \\
  & LP-selected   & 0.6241 & 0.6222 & $-$0.0019 \\
  & Selection gap & $+$0.0059 & $+$0.0018 & $-$0.0041 \\
\midrule
\multirow{3}{*}{RE}
  & Greedy        & 0.2913 & 0.2880 & $-$0.0033 \\
  & LP-selected   & 0.2845 & 0.3050 & \textbf{$+$0.0205} \\
  & Selection gap & $-$0.0068 & $+$0.0169 & \textbf{$+$0.0237} \\
\bottomrule
\end{tabular}

\vspace{8pt}

{\small\textsc{Panel C: QE Signal Quality (Spearman $\rho$)}}
\vspace{2pt}

\begin{tabular}{l ccc ccc}
\toprule
 & \multicolumn{3}{c}{NER ($n{=}529$)} & \multicolumn{3}{c}{RE ($n{=}397$)} \\
\cmidrule(lr){2-4} \cmidrule(lr){5-7}
Signal & Constr. & Free & $\Delta$ & Constr. & Free & $\Delta$ \\
\midrule
SJ & 0.360 & 0.377 & $+$0.018 & 0.250 & 0.273 & $+$0.023 \\
FK & 0.267 & 0.320 & $+$0.053 & 0.275 & 0.315 & $+$0.039 \\
LP & 0.205 & 0.213 & $+$0.008 & 0.266 & 0.304 & $+$0.038 \\
EM & 0.295 & 0.317 & $+$0.023 & 0.134 & 0.137 & $+$0.002 \\
VC & 0.379 & 0.409 & $+$0.030 & 0.350 & 0.323 & $-$0.026 \\
\bottomrule
\end{tabular}

\vspace{8pt}

{\small\textsc{Panel D: Parse Rate (Free-form)}}
\vspace{2pt}

\begin{tabular}{lc}
\toprule
Metric & Value \\
\midrule
Total samples          & 4408 \\
Parse rate             & 99.95\% \\
Instances 8/8 valid    & 549\,/\,551 (99.6\%) \\
\bottomrule
\end{tabular}

\end{table*}

\begin{table}[t]
\centering
\small
\caption{Cross-dataset free-form decoding summary (Qwen3-8B, $N{=}8$, $T{=}1.0$). LP std ratio = free-form\,/\,constrained within-instance LP std. Arrows: constrained $\to$ free-form.}
\label{tab:freeform_summary}
\begin{tabular}{@{}lccc@{}}
\toprule
 & SciERC & CoNLL & Few-NERD \\
\midrule
LP std ratio & $0.97{\times}$ & $0.86{\times}$ & $1.00{\times}$ \\
Degen.\ (\%) & $13.4 {\to} 12.1$ & $56.5 {\to} 58.8$ & $9.1 {\to} 9.1$ \\
MV $\Delta$ (pp) & ${-}3.3 {\to} {-}1.4$ & ${-}1.8 {\to} {-}2.1$ & ${-}9.6 {\to} {-}9.6$ \\
Greedy F1 & $.618 {\to} .620$ & $.907 {\to} .905$ & $.789 {\to} .789$ \\
\bottomrule
\end{tabular}
\end{table}

\paragraph{CoNLL NER free-form results.}
Extending the free-form ablation to CoNLL NER (Qwen3-8B, 4 seeds, $N{=}8$, $T{=}1.0$), LP range median is $0.058{\pm}0.007$ across seeds, higher than SciERC ($0.041$) but still compressed given CoNLL's 86.8\% tied fraction under constrained decoding.
Free-form decoding reduces the tied fraction modestly, yet selection remains below single-sample: LP-selected F1${=}0.896$ vs.\ single-sample $0.904$ ($\Delta{=}{-}0.77$~pp).
Quality estimation under free-form decoding yields SJ $\rho{=}0.461$, EM $\rho{=}0.471$, LP $\rho{=}0.326$, comparable to constrained results (Table~\ref{tab:quality_estimation}).
The failure to improve selection despite reduced compression is consistent with CoNLL's low oracle headroom (4.0~pp at $N{=}8$), which limits the selection space regardless of score distribution quality.

\paragraph{Few-NERD NER free-form results.}
On Few-NERD (Qwen3-8B, $N{=}8$, $T{=}1.0$, seed${=}42$), free-form decoding produces near-identical LP distributions (LP std ratio $1.00{\times}$) and unchanged degeneracy ($9.1$\%).
Greedy F1 is unchanged ($.789$); MV degradation is comparable under both regimes ($-9.6$~pp).
Table~\ref{tab:freeform_summary} summarizes all three datasets: LP std ratios range from $0.86{\times}$ to $1.00{\times}$, confirming that LP compression is model-intrinsic and not an artifact of constrained decoding.

\section{LP Selection Failure Stability}\label{app:lp_failure_stability}

We analyze RE LP selection failures across three random seeds (constrained decoding, $N{=}8$, $T{=}1.0$, Qwen3-8B) to assess whether failures are systematic or stochastic.

\paragraph{Pairwise overlap.}
LP selection failures (instances where LP-selected F1 $<$ single-sample F1) show moderate cross-seed stability: the average pairwise Jaccard similarity is 0.42 (range: 0.38--0.46); failures are partially systematic (attributable to instance complexity) but not fully deterministic.
The 3-way intersection contains 34 instances (27.9\% of any-seed failures), which constitutes a stable failure core.

\paragraph{Instance complexity co-occurs with failure.}
Table~\ref{tab:stable_failures} stratifies instances by failure stability.
Stable failures (all 3 seeds) involve more gold relations (2.12 vs.\ 1.62 for never-fail instances) and higher single-sample F1 (0.688 vs.\ 0.429); LP mis-selects on instances where the model already performs well, precisely the instances where selection would be most valuable.
Failure rate increases monotonically with gold relation count: 4.3\% at 1 relation, 20.7\% at 3 relations, 26.2\% at 6 relations.

\begin{table}[t]
\centering
\small
\caption{RE LP selection failure stability across 3 seeds (SciERC, $N{=}8$). Stable failures concentrate on relation-dense instances with high single-sample F1.}
\label{tab:stable_failures}
\resizebox{\columnwidth}{!}{
\begin{tabular}{lcccc}
\toprule
Group & $n$ & Gold rels & Greedy F1 & LP $\Delta$F1 \\
\midrule
Stable fail (3 seeds) & 34 & 2.12 & 0.688 & $-$0.499 \\
Fail 2+ seeds & 26 & 2.46 & 0.774 & $-$0.450 \\
Fail 1 seed & 62 & 2.32 & 0.692 & $-$0.239 \\
Never fail & 429 & 1.62 & 0.429 & $+$0.041 \\
\bottomrule
\end{tabular}
}
\end{table}

\paragraph{Flip instances.}
Only 5 RE instances with gold relations flip from LP failure (constrained) to LP success (free-form), with mean gold relation count 3.67 (well above the population mean of 1.86).
These instances show large constrained LP errors ($\Delta{=}{-}0.430$) that reverse under free-form decoding ($\Delta{=}{+}0.503$), but their rarity (5/397) confirms that the RE selection flip is attributable to aggregate variance reallocation rather than systematic instance-level correction.

\section{$N$-Scaling Detailed Analysis}\label{app:n_scaling}

Tables~\ref{tab:n-scaling}--\ref{tab:sigma-comparison} provide detailed $N$-scaling analysis on Qwen3-8B SciERC with 3-seed statistics.

\begin{table}[t]
\centering
\caption{Effect of sample pool size on signal quality (Qwen3-8B, SciERC, 3-seed mean Spearman $\rho$).}
\label{tab:n-scaling}
\small
\resizebox{\columnwidth}{!}{
\begin{tabular}{l ccc ccc}
\toprule
 & \multicolumn{3}{c}{\textbf{NER} $\rho$} & \multicolumn{3}{c}{\textbf{RE} $\rho$} \\
\cmidrule(lr){2-4} \cmidrule(lr){5-7}
Signal & $N\!=\!8$ & $N\!=\!16$ & $\Delta$ & $N\!=\!8$ & $N\!=\!16$ & $\Delta$ \\
\midrule
SJ & 0.3641 & 0.4105 & \textbf{+0.0464} & 0.2485 & 0.2758 & \textbf{+0.0273} \\
FK & 0.2536 & 0.2833 & \textbf{+0.0297} & 0.2903 & 0.2964 & \textbf{+0.0061} \\
VC & 0.3861 & 0.4092 & \textbf{+0.0231} & 0.3088 & 0.3738 & \textbf{+0.0650} \\
LP & 0.2167 & 0.2393 & \textbf{+0.0226} & 0.2970 & 0.3100 & \textbf{+0.0130} \\
EM & 0.3071 & 0.3395 & \textbf{+0.0324} & 0.1353 & 0.1506 & \textbf{+0.0153} \\
\midrule
Mean & 0.3055 & 0.3364 & \textbf{+0.0308} & 0.2560 & 0.2813 & \textbf{+0.0253} \\
\bottomrule
\end{tabular}
}
\end{table}

\begin{table}[t]
\centering
\caption{Cross-seed signal ranking stability (Kendall's $W$, 3 seeds).}
\label{tab:ranking-stability}
\small
\resizebox{\columnwidth}{!}{
\begin{tabular}{ll cc l}
\toprule
Model & Task & $N$ & Kendall's $W$ & Consensus ranking \\
\midrule
Qwen  & NER & 8  & 1.000 & VC $>$ SJ $>$ EM $>$ FK $>$ LP \\
Qwen  & NER & 16 & 0.956 & SJ $\approx$ VC $>$ EM $>$ FK $>$ LP \\
Qwen  & RE  & 8  & 0.822 & VC $>$ LP $>$ FK $>$ SJ $>$ EM \\
Qwen  & RE  & 16 & 0.956 & VC $>$ LP $>$ FK $>$ SJ $>$ EM \\
LLaMA & NER & 8  & 1.000 & VC $>$ SJ $>$ EM $>$ FK $>$ LP \\
\bottomrule
\end{tabular}
}
\end{table}

\begin{table}[t]
\centering
\caption{Cross-seed standard deviation ($\sigma$) of Spearman $\rho$ at $N\!=\!8$ vs.\ $N\!=\!16$ (Qwen3-8B, SciERC, 3 seeds). $\sigma$ is population standard deviation.}
\label{tab:sigma-comparison}
\small
\begin{tabular}{l cc cc}
\toprule
 & \multicolumn{2}{c}{\textbf{NER} $\sigma_\rho$} & \multicolumn{2}{c}{\textbf{RE} $\sigma_\rho$} \\
\cmidrule(lr){2-3} \cmidrule(lr){4-5}
Signal & $N\!=\!8$ & $N\!=\!16$ & $N\!=\!8$ & $N\!=\!16$ \\
\midrule
SJ & 0.0179 & 0.0102 & 0.0015 & 0.0067 \\
FK & 0.0157 & 0.0079 & 0.0121 & 0.0143 \\
VC & 0.0187 & 0.0260 & 0.0244 & 0.0158 \\
EM & 0.0076 & 0.0031 & 0.0157 & 0.0049 \\
LP & 0.0125 & 0.0017 & 0.0218 & 0.0152 \\
\midrule
Mean & 0.0145 & 0.0098 & 0.0151 & 0.0114 \\
\bottomrule
\end{tabular}
\end{table}

Signal rankings are consistent across our 3 seeds (Kendall's $W{=}1.000$ for NER at $N{=}8$, both models), dropping to $W{=}0.956$ at $N{=}16$ where SJ and VC converge (Table~\ref{tab:ranking-stability}); however, with only $n{=}3$ seeds statistical power is insufficient to confirm stability, and additional seeds may reveal ranking changes.
Doubling $N$ reduces cross-seed $\sigma$ by 32\% on NER and 24\% on RE (Table~\ref{tab:sigma-comparison}), with the largest reductions for SJ ($-43\%$ NER) and EM ($-59\%$ NER); additional samples stabilize structural signals more than LP.

\paragraph{$N$-scaling: selection headroom vs.\ selection gain.}
Table~\ref{tab:n_scaling_selection} traces the effect of sample count on degeneracy, oracle headroom, and per-signal selection gain.
Oracle headroom grows roughly linearly with $N$ ($+4.4$~pp at $N{=}2$ to $+20.0$~pp at $N{=}32$).
LP selection improves monotonically with $N$ ($+0.39$~pp at $N{=}8$ to $+0.95$~pp at $N{=}32$), but even at $N{=}32$, where degeneracy falls to 4\% and oracle headroom reaches ${\sim}20$~pp, the best signal captures less than 5\% of the available improvement.
Degeneracy drops monotonically (43\% at $N{=}2$ to 4\% at $N{=}32$), yet reduced degeneracy does not unlock selection, consistent with the LoRA-rank ablation finding (\S\ref{sec:robustness_main}).
VC selection degrades with increasing $N$ ($-2.66$~pp at $N{=}2$ to $-4.35$~pp at $N{=}32$), consistent with mode-seeking bias in voting-based selection as the candidate pool grows.

\paragraph{Few-NERD $N$-scaling: MV at $N{=}16$.}
On Few-NERD (Qwen3-8B, 3-seed), increasing from $N{=}8$ to $N{=}16$ reduces degeneracy ($8.8$\%$\to$$4.0{\pm}1.4$\%) and attenuates MV degradation ($-2.74$\,pp$\to$$-1.65{\pm}0.61$\,pp), but MV remains negative.
Additional samples reduce the MV penalty but do not reverse it, consistent with the correlated-error mechanism: more samples provide better frequency estimates but cannot overcome systematic shared mistakes.
SJ $\rho$ is used for quality estimation at $N{=}16$ (FK conditional $\rho$ shows sign instability across seeds: mean${=}0.051{\pm}0.062$).

\begin{table}[t]
\centering
\small
\caption{Effect of sample count on selection (SciERC NER, Qwen3-8B, $T{=}1.0$, gold-filtered, seed${=}42$; exp\_025).\protect\footnotemark\ ``Best $\Delta$'': highest selection gain across all 5 signals. Per-signal $\Delta$ in pp vs.~greedy.}
\label{tab:n_scaling_selection}
\resizebox{\columnwidth}{!}{
\begin{tabular}{ccccccccc}
\toprule
$N$ & Degen.\% & Oracle (pp) & LP & SJ & FK & EM & VC & Best $\Delta$ \\
\midrule
2 & 42.7\% & $+$4.4 & $-$0.89 & $-$2.72 & $-$2.72 & $-$2.72 & $-$2.66 & $-$0.89 \\
4 & 21.0\% & $+$10.2 & $+$0.03 & $-$0.31 & $-$0.40 & $-$0.74 & $-$2.16 & $+$0.03 \\
8 & 11.0\% & $+$14.5 & $+$0.39 & $+$0.20 & $+$0.15 & $-$0.77 & $-$2.87 & $+$0.39 \\
16 & 7.4\% & $+$17.7 & $+$0.78 & $+$0.24 & $-$0.02 & $-$0.06 & $-$3.95 & $+$0.78 \\
32 & 4.0\% & $+$20.0 & $+$0.95 & $+$0.35 & $+$0.65 & $+$0.22 & $-$4.35 & $+$0.95 \\
\bottomrule
\end{tabular}
}
\end{table}
\footnotetext{$N$-scaling analysis uses gold-filtered evaluation ($n{=}529$, greedy F1${=}0.640$, seed 42). Values differ from Table~\ref{tab:selection} (0.645, 2-seed) and Table~\ref{tab:fewnerd} (0.663, 3-seed) due to seed and aggregation differences; relative trends are consistent.}

\paragraph{Bootstrap significance at $N{=}32$.}
Table~\ref{tab:bootstrap_n32} reports bootstrap confidence intervals and $p$-values for selection at $N{=}32$ ($B{=}10{,}000$, $n{=}529$).
No signal achieves significant improvement over greedy at $\alpha{=}0.05$ (LP: $p{=}0.091$; FK: $p{=}0.101$), while VC is significantly \textit{worse} ($-4.36$~pp, $p{=}1.0$).
However, all signals except VC significantly outperform random selection ($p{<}0.001$, permutation test); the signals carry meaningful quality information yet fail to translate it into reliable selection gains.
The oracle gap ($+20.0$~pp, $p{<}0.001$) confirms substantial unrealized headroom.

\begin{table}[t]
\centering
\small
\caption{Bootstrap significance test at $N{=}32$ (SciERC NER, Qwen3-8B, $T{=}1.0$, $B{=}10{,}000$, $n{=}529$ gold-non-empty instances; 22 instances with no gold entities are excluded as they contribute zero selection differential). $\Delta$: selection gain vs.\ greedy. CI: 95\% bootstrap percentile interval. $p_{\text{greedy}}$: one-sided $p$-value vs.\ greedy; $p_{\text{rand}}$: permutation test vs.\ random selection.}
\label{tab:bootstrap_n32}
\resizebox{\columnwidth}{!}{
\begin{tabular}{lcccc}
\toprule
Signal & $\Delta$ (pp) & 95\% CI & $p_{\text{greedy}}$ & $p_{\text{rand}}$ \\
\midrule
LP & $+$0.95 & [$-$0.41, $+$2.34] & 0.091 & ${<}$0.001 \\
FK & $+$0.66 & [$-$0.36, $+$1.67] & 0.101 & ${<}$0.001 \\
SJ & $+$0.36 & [$-$0.66, $+$1.38] & 0.244 & ${<}$0.001 \\
EM & $+$0.22 & [$-$0.99, $+$1.44] & 0.355 & ${<}$0.001 \\
VC & $-$4.36 & [$-$6.26, $-$2.37] & 1.000 & 0.741 \\
\midrule
Oracle & $+$20.00 & [$+$17.98, $+$22.10] & ${<}$0.001 & --- \\
\bottomrule
\end{tabular}
}
\end{table}

\paragraph{Oracle ceiling scaling law.}
Oracle F1 follows $F_1(N) = a - b \cdot N^{-c}$, with $R^2 > 0.999$ across all model--dataset fits (Figure~\ref{fig:oracle_scaling}).
The exponent $c$ varies with dataset complexity: Few-NERD ($c{\approx}0.61$) scales comparably to reasoning-task estimates ($c{\approx}0.5$--$0.7$; \citealp{brown2024monkeys}), while SciERC converges more slowly ($c{\approx}0.35$), consistent with complex IE producing less marginal diversity per additional sample.
CoNLL's narrow headroom ($b{=}0.04$--$0.10$) makes $c$ unreliable (inter-model variance $2.5{\times}$); we report only $a$ and $b$ for CoNLL.
A single stochastic sample ($N{=}1$, $T{=}1.0$) underperforms greedy by 1--3\,pp because sampling noise outweighs diversity; $N{\geq}4$ is needed to match greedy.


\paragraph{Degeneracy and inverse scaling.}\label{app:degen_inverse}
Table~\ref{tab:degen_inverse} stratifies $N$-transitions by instance degeneracy status across all dataset--regime combinations.
Degenerate instances exhibit near-zero inverse scaling ($1.3$\% vs.\ $11.1$\% for non-degenerate; OR${=}0.109$, $p{<}10^{-300}$, Fisher's exact test over $195{,}675$ transitions).
This is expected: degenerate instances produce identical entity sets across all $N$ samples, so MV F1 is constant by construction---they cannot exhibit inverse scaling.
The effect holds across all settings (Table~\ref{tab:degen_inverse}), with the sole exception of SciERC pretrained ($p{=}0.083$), where low degeneracy (${\leq}10$\%) yields too few degenerate transitions for adequate power.
Pretrained models show a stronger protective effect (OR${=}0.066$ vs.\ $0.166$ for finetuned), consistent with higher overall inverse scaling rates creating more room for the degeneracy--stability association.

\begin{table}[t]
\centering
\small
\caption{Inverse scaling rate by degeneracy status. An $N{\to}N'$ transition is ``inverse'' if MV F1 at $N'$ is strictly lower than at $N$. OR: odds ratio (Fisher's exact test).}
\label{tab:degen_inverse}
\resizebox{\columnwidth}{!}{%
\begin{tabular}{@{}lrrrrrc@{}}
\toprule
Setting & \multicolumn{2}{c}{Degenerate} & \multicolumn{2}{c}{Non-degenerate} & OR & $p$ \\
\cmidrule(lr){2-3} \cmidrule(lr){4-5}
 & Trans. & Inv.\% & Trans. & Inv.\% & & \\
\midrule
SciERC FT & 485 & 0.8 & 2{,}160 & 9.6 & 0.078 & ${<}10^{-6}$ \\
CoNLL FT & 8{,}256 & 0.3 & 5{,}524 & 4.1 & 0.060 & ${<}10^{-6}$ \\
Few-NERD FT & 22{,}298 & 1.9 & 75{,}397 & 8.4 & 0.213 & ${<}10^{-6}$ \\
SciERC PT & 81 & 3.7 & 2{,}564 & 9.9 & 0.350 & 0.083 \\
CoNLL PT & 5{,}284 & 0.3 & 8{,}496 & 7.9 & 0.033 & ${<}10^{-6}$ \\
Few-NERD PT & 18{,}007 & 1.4 & 47{,}123 & 16.8 & 0.071 & ${<}10^{-6}$ \\
\midrule
\textbf{All} & \textbf{54{,}411} & \textbf{1.3} & \textbf{141{,}264} & \textbf{11.1} & \textbf{0.109} & $\mathbf{{<}10^{-300}}$ \\
\bottomrule
\end{tabular}}
\end{table}

\section{Cross-Scale and Cross-Family Comparison}\label{app:cross_scale}

\begin{table}[t]
\centering
\small
\caption{Cross-scale and cross-family comparison ($N{=}8$, $T{=}1.0$). Hdr.\ = oracle $-$ greedy (pp). Degen.: instance degeneracy (\%). F1 values are entity-type-averaged macro-F1, which differs from the instance-averaged F1 in Table~\ref{tab:fewnerd}. $^\star$Qwen2.5; unmarked: Qwen3. $^\dagger$Not computed. $^\ddagger$200-instance subset; 10.1\% parse failures excluded.}
\label{tab:scale_comparison}
\resizebox{\columnwidth}{!}{%
\begin{tabular}{@{}llccccccc@{}}
\toprule
Dataset & Regime & Greedy & Hdr. & Degen. & VC & SJ & FK & LP \\
\midrule
SciERC & 8B FT  & .667 & 11.4 & 12.2 & .326 & .324 & .177 & .159 \\
       & 72B ZS$^\star$ & .399 & 15.0 & 21.0 & .277 & .217 & .227 & .081 \\
\midrule
Few-NERD & 7B FT$^\star$ & .784 & 9.2 & 31.7 & .404 & .388 & .315 & .346 \\
         & 8B FT & .748 & 12.4 & 13.0 & .386 & .300 & .313 & .220 \\
         & 14B FT & .798 & 8.1 & ---$^\dagger$ & ---$^\dagger$ & ---$^\dagger$ & ---$^\dagger$ & ---$^\dagger$ \\
         & 72B ZS$^\star$ & .628 & ---$^\dagger$ & 46.2 & .416 & .345 & ---$^\dagger$ & .119 \\
         & GPT-5.5 ZS$^\ddagger$ & .638 & 8.2 & 50.0 & ---$^\dagger$ & .346 & ---$^\dagger$ & ---$^\dagger$ \\
\midrule
CoNLL  & 8B FT  & .908 & ${\approx}$4 & 61.7 & .424 & .436 & .400 & .302 \\
       & 72B ZS$^\star$ & .712 & 5.3 & 67.7 & .297 & .304 & .302 & $-$.002 \\
\bottomrule
\end{tabular}}
\end{table}

LP $\rho$ drops without fine-tuning (Table~\ref{tab:scale_comparison}) and LP selection fails on all 72B zero-shot datasets. Structural signals generalize better (Few-NERD VC $\rho{=}0.416$). Degeneracy increases at 72B (Few-NERD: $46.2$\% vs.\ $13.0$\%; CoNLL: $67.7$\% vs.\ $61.7$\%); the gap persists under this specific 72B zero-shot configuration, though scale, training regime, and quantization effects are confounded (see Limitations).
LP-weighted construction yields $+1.53$\,pp on SciERC 72B ZS ($p{=}0.008$, CI [$+0.25$, $+2.82$]), recovering $10.2$\% of headroom. Construction fails on Few-NERD ($-0.20$\,pp) and CoNLL ($-0.37$\,pp) where degeneracy exceeds $46$\%---consistent with FT-regime degeneracy gating.

\paragraph{Few-shot demonstrations improve QE but not construction.}
Adding 3-shot demonstrations to 72B improves greedy on SciERC and CoNLL ($+3.3$--$6.2$\,pp) but not Few-NERD ($-0.9$\,pp), and strengthens structural QE (CoNLL SJ $\rho$ $.304{\to}.442$, LP $-.002{\to}.106$). Yet construction is ineffective across all three datasets (SciERC $-0.08$\,pp, CoNLL $+0.03$\,pp, Few-NERD $+0.02$\,pp; all n.s.). On SciERC, this failure despite low degeneracy shows the bottleneck is regime-dependent (Appendix~\ref{app:fs_error}); on Few-NERD, high FS degeneracy ($45.2$\%) places it in the expected failure regime.

\paragraph{Degeneracy is model-family-dependent, not scale-dependent.}
Qwen2.5 exhibits systematically higher degeneracy than Qwen3 on Few-NERD (7B--72B: $31.7$--$46.2$\% vs.\ Qwen3-8B: $13.0$\%; Table~\ref{tab:scale_comparison}), confirming degeneracy is shaped by architecture rather than parameter count. Construction fails for all Qwen2.5 configurations on Few-NERD (both n.s.); 3$\to$5 epochs triples degeneracy within the same model (\S\ref{sec:fewnerd}).

\paragraph{Qwen3 8B vs.\ 4B scale ablation.}

\begin{table}[h]
\centering
\small
\caption{Scale ablation: Qwen3 8B vs.\ 4B QE $\rho$ ($N{=}8$, $T{=}1.0$, seed 42). RE in Appendix~\ref{app:re_extension}.}
\label{tab:scale_ablation_app}
\begin{tabular}{@{}ll cc r@{}}
\toprule
& & \textbf{8B} & \textbf{4B} & $\Delta$ \\
\midrule
CoNLL & SJ & .436 & .422 & $-$.014 \\
      & FK & .400 & .418 & $+$.018 \\
SciERC & SJ & .360 & .422 & $+$.062 \\
       & FK & .267 & .326 & $+$.059 \\
\bottomrule
\end{tabular}
\end{table}

Degeneracy scales with capability: 8B shows 12.2\% vs.\ 4B's 10.4\% on SciERC (Table~\ref{tab:degeneracy}).

\paragraph{14B scale extension.}
Qwen3-14B fine-tuned on Few-NERD achieves greedy F1${=}0.798$ ($+0.9$\,pp over 8B) with oracle headroom of 8.1\,pp ($+1.2$\,pp over 8B's 6.9\,pp) and $\rho$ advantage $+0.029$ over 8B.
The ceiling grows with model capability: larger models improve both greedy and oracle performance, but the aggregation gap widens rather than closes.

\paragraph{Pre-trained baseline (no fine-tuning).}
Table~\ref{tab:pretrained_summary} summarizes pretrained-vs-finetuned MV performance across three datasets (Qwen2.5-7B-Instruct pretrained, Qwen3-8B LoRA finetuned; $N{=}8$, $T{=}1.0$).
Across all three datasets, pretrained models benefit from MV aggregation ($+1.7$ to $+9.2$\,pp), while all finetuned models show negligible or negative MV gains.

\begin{table}[t]
\centering
\small
\caption{Pretrained vs.\ finetuned MV ($T{=}1.0$). Pretrained: Qwen2.5-7B-Instruct; finetuned: Qwen3-8B LoRA ($N{=}8$ unless noted). MV: strict majority ($>N/2$). Capture: $\text{MV}\,\Delta / \text{oracle headroom}$. CoNLL pretrained single-seed; others 3-seed mean${\pm}\sigma$ ($N{=}8$) or single-seed ($N{=}16{,}32$).}
\label{tab:pretrained_summary}
\resizebox{\columnwidth}{!}{%
\begin{tabular}{@{}llccccc@{}}
\toprule
Dataset & Setting & Greedy & MV $\Delta$ (pp) & Degen.\% & Capture & $\rho_{\textup{SJ}}$ \\
\midrule
CoNLL & Pretrained & .536 & $+$9.2 & 36 & 98\% & .467 \\
      & Finetuned & .912 & ${<}$0.1 & 57--60 & ${<}$2\% & .41--.45 \\
\addlinespace
Few-NERD & Pretrained ($N{=}8$) & .538 & $+$3.9${\pm}$0.2 & 11.7 & 36\% & .436 \\
         & Pretrained ($N{=}16$) & .537 & $+$4.1 & 7.6 & 29\% & .453 \\
         & Pretrained ($N{=}32$) & .539 & $+$3.9 & 5.4 & 19\% & .456 \\
         & Finetuned & .791 & $-$2.7${\pm}$0.5 & 8.8 & neg. & .31 \\
\addlinespace
SciERC & Pretrained & .244 & $+$1.7${\pm}$0.4 & 0.8 & --- & --- \\
       & Finetuned & .667 & ${<}$0.1 & 12.2 & --- & .36 \\
\bottomrule
\end{tabular}}
\end{table}

On CoNLL, MV captures $98$\% of oracle headroom for the pretrained model (greedy${=}0.536$, MV${=}0.628$, oracle${=}0.630$, degeneracy${=}36$\%) versus ${<}2$\% finetuned (greedy${=}0.912$, MV${=}0.912$, oracle${=}0.954$, degeneracy $57$--$60$\%).
Signal quality is comparable ($\rho_{\textup{SJ}}{=}0.47$ pretrained vs.\ $0.41$--$0.45$ finetuned).

On Few-NERD ($n{=}32{,}565$, 3-seed), finetuned degeneracy ($8.8$\%) is \textit{lower} than pretrained ($11.7$\%), yet finetuned MV is negative ($-2.74{\pm}0.54$\,pp).
Doubling the sample budget ($N{=}8{\to}16$) maintains the pretrained MV advantage ($+4.06$\,pp) while reducing degeneracy ($11.7{\to}7.6$\%) and expanding oracle headroom ($+2.9$\,pp); however, capture rate decreases ($36{\to}29$\%), indicating diminishing marginal returns from additional samples.
$N{=}32$ confirms the plateau (MV $\Delta{=}+3.9$\,pp), with capture rate declining further to $19$\%.
Fine-tuning thus decouples agreement from quality through dataset-dependent pathways: diversity elimination on CoNLL (degeneracy $36{\to}57$--$60$\%), but correlated error patterns on Few-NERD.
A baseline quality confound exists (pretrained greedy${=}0.538$ vs.\ finetuned${=}0.791$); per-type analysis (Appendix~\ref{app:pertype}) partially addresses this by showing MV degradation concentrates in low-F1 entity types ($\rho{=}0.90$, $p{=}0.002$).

On SciERC, MV gains $+1.67{\pm}0.40$\,pp despite very low greedy F1 ($0.244$) with near-zero degeneracy ($0.8$\%), contradicting a ``MV only works at high baseline'' explanation.
Structural signals outperform surface on CoNLL ($\rho_{\textup{SJ}}{=}0.467$, $\rho_{\textup{FK}}{=}0.435$) and LP AUROC${=}0.580$.
SciERC's low greedy F1 does establish a scope boundary for entity-level \textit{construction}: 46\% of samples produce F1${=}0$, so partial-correctness-based aggregation cannot improve over greedy.

\subsection{Per-Type MV Analysis}\label{app:pertype}

\begin{table}[t]
\centering
\small
\caption{Per-entity-type MV gain on Few-NERD (Qwen3-8B, $N{=}8$, $T{=}1.0$, 3-seed mean). Strict MV ($\theta{>}N/2$). Types sorted by finetuned greedy F1. 7/8 types flip MV direction between pretrained and finetuned.}
\label{tab:pertype_mv}
\begin{tabular}{@{}lrccc@{}}
\toprule
& & & \multicolumn{2}{c}{MV $\Delta$ (pp)} \\
\cmidrule(lr){4-5}
Entity type & $n$ & Greedy F1 & Pretr. & Finetuned \\
\midrule
person & ${\sim}$7{,}200 & .820 & $+$4.17 & $-$0.67 \\
location & ${\sim}$8{,}800 & .797 & $+$2.47 & $-$1.52 \\
art & ${\sim}$1{,}250 & .755 & $+$0.73 & $-$1.51 \\
organization & ${\sim}$6{,}700 & .725 & $+$3.15 & $-$2.65 \\
building & ${\sim}$1{,}540 & .696 & $+$1.98 & $-$3.90 \\
product & ${\sim}$2{,}260 & .657 & $+$1.78 & $-$3.38 \\
event & ${\sim}$1{,}315 & .645 & $+$2.34 & $-$3.17 \\
other & ${\sim}$3{,}485 & .639 & $-$0.45 & $-$3.02 \\
\midrule
\multicolumn{3}{@{}l}{$r$(greedy F1, MV\,$\Delta$)} & .86 & .90 \\
\multicolumn{3}{@{}l}{$p$-value} & .006 & .002 \\
\bottomrule
\end{tabular}
\end{table}
MV direction flips between pretrained and finetuned decoding for 7 of 8 entity types (Table~\ref{tab:pertype_mv}): types that benefit from MV under pretrained decoding ($+0.73$ to $+4.17$\,pp) are uniformly harmed under finetuning ($-0.67$ to $-3.90$\,pp), with only \textit{other} (lowest-F1) remaining negative in both regimes.
The correlation between greedy F1 and MV $\Delta$ is strong in both regimes ($r{=}0.86$, $p{=}0.006$ pretrained; $r{=}0.90$, $p{=}0.002$ finetuned), indicating entity-type difficulty predicts MV magnitude regardless of direction.
This rules out entity-type-specific factors as the explanation for MV failure: the same types succeed with pretrained decoding and fail with finetuned decoding, implicating the finetuning process itself.
Under a liberal MV threshold ($\theta{\geq}N/2$ instead of $\theta{>}N/2$), per-type correlation direction is unstable; all results above use strict MV, consistent with the main paper.

\subsection{Frontier Model Self-Consistency}\label{app:frontier_sc}

\begin{table}[h]
\centering
\small
\caption{Cross-family frontier SC on Few-NERD ($N{=}8$, $T{=}1.0$, 200 instances, zero-shot). Parse failures excluded from all metrics; $N_{\textit{eff}}$: mean effective samples after exclusion. Selection gap $= $ oracle $-$ MV. Text-based entity matching (Appendix~\ref{app:frontier_sc}, entity matching paragraph). Qwen3-8B FT shown for reference.}
\label{tab:frontier_sc_cross}
\resizebox{\columnwidth}{!}{%
\begin{tabular}{@{}lcccccc@{}}
\toprule
Model & Greedy & MV ($\Delta$) & Oracle & Degen.\% & Sel.\ gap & Parse\% \\
\midrule
DeepSeek & .399 & .424 ($+2.5$) & .714 & 5.5 & .290 & 46.9 \\
Gemini 2.5 Pro & .630 & .645 ($+1.5$) & .758 & 28.5 & .113 & 11.9 \\
GPT-5.5 & .638 & .652 ($+1.4$) & .720 & 50.0 & .068 & 10.1 \\
Claude Son.~4 & .607 & .623 ($+1.6$) & .654 & 74.0 & .032 & 6.4 \\
\midrule
Qwen3-8B FT & .789 & .792 ($+0.3$) & .858 & 13.0 & .066 & ${\approx}$0 \\
\bottomrule
\end{tabular}}
\end{table}

All four frontier families exhibit the correlation-selection gap (Table~\ref{tab:frontier_sc_cross}): MV recovers only a fraction of oracle headroom regardless of model capability.
Degeneracy and selection gap are negatively correlated across families: Claude Sonnet~4 (74\% degen, 0.032 gap) vs.\ DeepSeek (5.5\%, 0.290)---higher degeneracy eliminates the instances where aggregation could help.
Gemini 2.5 Pro achieves the highest oracle F1 (.758) and the highest QE correlation ($\rho_{\textup{SJ}}{=}0.350$, $\rho_{\textup{agr}}{=}0.389$) among frontier models, with balanced degeneracy (28.5\%) preserving a 0.113 selection gap---the largest gap among models with ${\geq}85$\% parse success.
DeepSeek's low degeneracy preserves the largest absolute gap (.290), but 47\% parse failures reduce effective $N$ to ${\approx}4.2$, limiting aggregation benefit.
GPT-5.5 occupies the middle ground: 50\% degeneracy with 6.8\,pp unrealized oracle headroom.

\paragraph{Parse failure handling.}
10.1\% of GPT-5.5 samples produced output that could not be parsed into valid entity sets.
These samples are excluded from all aggregation metrics (degeneracy, oracle, MV, selection); for instances where all $N$ samples fail parsing, the instance is excluded entirely.
The effective sample size per instance is therefore $N_{\textit{eff}} \leq 8$ (mean ${\approx}7.2$).

\paragraph{Entity matching methodology.}
Local experiments (fine-tuned models with constrained decoding) use offset-based matching: two entities match iff \texttt{(start, end, type)} are identical.
Frontier API experiments use text-based matching: \texttt{(text.lower(), type)}, because API outputs lack character-offset information.
Text-based matching is slightly more permissive (e.g., ignoring whitespace differences at span boundaries), producing ${\approx}2$\,pp higher MV $\Delta$ on controlled comparisons.
All within-pipeline comparisons (greedy vs.\ MV vs.\ oracle) use the same matching, so relative conclusions are unaffected; cross-pipeline absolute F1 values should not be directly compared.

\section{Temperature Ablation}\label{app:t_ablation}

Table~\ref{tab:t_ablation} reports the effect of sampling temperature on QE signals, oracle headroom, and output quality across four temperatures on SciERC NER.

\begin{table}[t]
\centering
\small
\caption{Temperature ablation on SciERC NER (Qwen3-8B, $N{=}8$, seed${=}42$). Oracle gap and QE $\rho$ across four temperatures. $\kappa$: Fleiss' $\kappa$ (inter-sample agreement). $T{=}0.5$ and $T{=}1.0$ use joint (NER+RE) subtask; $T{=}0.8$ and $T{=}1.2$ use NER-only.}
\label{tab:t_ablation}
\resizebox{\columnwidth}{!}{
\begin{tabular}{ccccccc}
\toprule
$T$ & Oracle gap (pp) & SJ $\rho$ & LP $\rho$ & $\kappa$ & Sample F1 \\
\midrule
0.5 & $+$9.1 & 0.297 & 0.218 & 0.577 & 0.636 \\
0.8 & $+$12.5 & 0.359 & 0.253 & 0.451 & 0.622 \\
1.0 & $+$13.0 & 0.360 & 0.205 & ${\approx}$0.39 & ${\approx}$0.612 \\
1.2 & $+$14.6 & 0.377 & 0.196 & 0.356 & 0.594 \\
\bottomrule
\end{tabular}
}
\end{table}

Oracle headroom increases monotonically with temperature ($+9.1$~pp at $T{=}0.5$ to $+14.6$~pp at $T{=}1.2$), confirming that higher temperature consistently expands the oracle ceiling under constrained structured decoding.
Structural QE signals improve with temperature: SJ $\rho$ rises from $0.297$ ($T{=}0.5$) to $0.377$ ($T{=}1.2$), consistent with more diverse outputs providing richer pairwise evidence.
LP $\rho$ peaks at $T{=}0.8$ ($0.253$) then degrades to $0.196$ at $T{=}1.2$.
This divergence confirms that LP compression is model-intrinsic: higher temperature increases between-instance LP variance (benefiting structural signals) but does not resolve within-instance compression, degrading LP's discriminative power.
Inter-sample agreement ($\kappa$) decreases monotonically ($0.577$ to $0.356$); higher temperature generates more diverse outputs. Mean sample F1 also decreases ($0.636$ to $0.594$), consistent with the quality--diversity trade-off.
The $T{=}0.7$ and $T{=}1.0$ comparison in Table~\ref{tab:robustness} is consistent with the trends reported here.

\paragraph{Temperature modulates degeneracy and LP spread.}
Table~\ref{tab:t_degeneracy} reports instance degeneracy (constant F1, gold-filtered) and median within-instance LP range across four temperatures on SciERC.
Both quantities are strictly monotonic in $T$: degeneracy falls from 41.4\% ($T{=}0.5$) to 14.9\% ($T{=}1.2$), while LP range increases from 0.015 to 0.058 nats.
Temperature thus provides a direct lever for practitioners to shift the degeneracy regime.
However, even at $T{=}1.2$ where degeneracy (14.9\%) approaches Few-NERD levels (${\sim}$13\%), within-instance $\rho(\text{LP}, \text{F1})$ remains at 0.225, far below Few-NERD's ${\sim}$0.5.
This confirms that temperature modulates degeneracy but cannot improve LP-quality alignment, which is largely associated with task structure and remains relatively stable across generation configurations (\S\ref{app:epoch_confound}). LP-quality alignment, not degeneracy, is the primary predictor of LP selection direction (\S\ref{sec:discussion}).

\begin{table}[h]
\centering
\small
\caption{Temperature sensitivity of degeneracy and LP spread (SciERC NER, Qwen3-8B, $N{=}8$). Degeneracy: constant-F1, gold-filtered. LP range: median within-instance range.}
\label{tab:t_degeneracy}
\resizebox{\columnwidth}{!}{
\begin{tabular}{ccccc}
\toprule
$T$ & Degen.\% & LP range (nats) & Fleiss' $\kappa$ & SJ mean \\
\midrule
0.5 & 41.4 & 0.015 & 0.572 & 0.830 \\
0.8 & 23.8 & 0.031 & 0.471 & 0.752 \\
1.0 & 17.8 & 0.042 & 0.405 & 0.700 \\
1.2 & 14.9 & 0.058 & 0.352 & 0.647 \\
\bottomrule
\end{tabular}
}
\end{table}

\section{Cross-Dataset OOD Analysis}\label{app:ood}

To test whether confidence signals generalize under domain shift, we evaluate models on out-of-distribution data: a SciERC-trained model applied to CoNLL instances (wider$\to$narrower schema) and vice versa (Table~\ref{tab:ood}).

\begin{table}[t]
\centering
\small
\caption{Cross-dataset OOD quality estimation (Spearman $\rho$, Qwen3-8B, $N{=}8$, $T{=}1.0$, seed${=}42$). SciERC$\to$CoNLL: SciERC-trained model on CoNLL instances ($n{=}2756$). CoNLL$\to$SciERC: CoNLL-trained model on SciERC instances ($n{=}529$, excluded from conclusions due to catastrophic model failure). In-domain $\rho$ from Table~\ref{tab:quality_estimation} (same model, same seed). \textsuperscript{***}$p{<}0.001$, \textsuperscript{*}$p{<}0.05$.}
\label{tab:ood}
\resizebox{\columnwidth}{!}{
\begin{tabular}{l ccc ccc}
\toprule
 & \multicolumn{3}{c}{\textbf{SciERC$\to$CoNLL}} & \multicolumn{3}{c}{\textbf{CoNLL$\to$SciERC}} \\
\cmidrule(lr){2-4} \cmidrule(lr){5-7}
Signal & $\rho$ & In-dom. & $\Delta$ & $\rho$ & In-dom. & $\Delta$ \\
\midrule
VC & .469\textsuperscript{***} & .379 & +.090 & +.105\textsuperscript{*} & .379 & $-$.274 \\
SJ & .378\textsuperscript{***} & .360 & +.018 & $-$.186\textsuperscript{***} & .360 & $-$.546 \\
EM & .300\textsuperscript{***} & .295 & +.005 & $-$.202\textsuperscript{***} & .295 & $-$.497 \\
FK & .285\textsuperscript{***} & .267 & +.018 & $-$.070 & .267 & $-$.337 \\
LP & .246\textsuperscript{***} & .205 & +.041 & $-$.152\textsuperscript{***} & .205 & $-$.357 \\
\bottomrule
\end{tabular}
}
\end{table}

\paragraph{Wider-to-narrower schema transfer.}
When the SciERC-trained model (6 NER types) is applied to CoNLL (4 types), all signals remain significant ($p{<}0.001$) and structural signals improve over their in-domain levels (same model, seed${=}42$, $N{=}8$; in-domain $\rho$ from Table~\ref{tab:quality_estimation}): VC reaches $\rho{=}0.469$ (in-domain $0.379$, $\Delta{=}{+}0.090$), SJ $\rho{=}0.378$ (in-domain $0.360$, $\Delta{=}{+}0.018$), and FK $\rho{=}0.285$ (in-domain $0.267$, $\Delta{=}{+}0.018$).
The wider training schema subsumes CoNLL's narrower schema, so the model produces reasonable outputs even out-of-domain (single-sample F1${=}0.522$), and structural signals benefit from the simpler target schema's reduced output ambiguity.

\paragraph{Narrower-to-wider schema transfer (excluded).}
The reverse direction is excluded from our conclusions due to catastrophic model failure, not signal failure.
The CoNLL-trained model recognizes only 4 news entity types (PER, ORG, LOC, MISC) and has no relation extraction capability; when applied to SciERC's 6 scientific entity types and 7 relation types, the model cannot produce meaningful structured outputs (single-sample NER F1${=}0.009$, RE F1${=}0.001$).
Under this regime, most signals flip to negative or near-zero correlations (Table~\ref{tab:ood}); only VC retains weak positive correlation ($\rho{=}0.105$, $p{=}0.016$).
These negative correlations reflect the model systematically generating confident but incorrect outputs, a capability failure analogous to WNUT-17 (\S\ref{sec:qe_results}) and not evidence against signal utility.

\paragraph{Asymmetry interpretation.}
The directional asymmetry is attributable to schema coverage: wider$\to$narrower preserves signal utility because the SciERC model's 6-type schema subsumes CoNLL's 4 types, while narrower$\to$wider fails because the CoNLL model lacks the entity and relation types required by SciERC.
This extends the capability boundary identified on WNUT-17 to the domain-shift setting: confidence signals require sufficient model capability as a prerequisite, and results from the CoNLL$\to$SciERC direction should not be interpreted as evidence of signal degradation under domain shift.

\section{Verbalized Confidence}\label{app:verbalized_confidence}

We evaluate verbalized confidence (VerbConf) as a second per-sample signal following \citet{lin2022teaching} and \citet{tian2023calibration}: after generating each structured output, the model is prompted to rate its confidence on a 1--10 scale.
Table~\ref{tab:verbconf} reports quality estimation and selection results.

\begin{table}[t]
\centering
\small
\caption{Verbalized confidence (VerbConf) vs.\ log-probability selection (Qwen3-8B, $N{=}8$, $T{=}1.0$, seed${=}42$). Parse fail: fraction of samples where the confidence score could not be extracted.}
\label{tab:verbconf}
\resizebox{\columnwidth}{!}{
\begin{tabular}{l cc cc}
\toprule
 & \multicolumn{2}{c}{\textbf{SciERC NER}} & \multicolumn{2}{c}{\textbf{CoNLL NER}} \\
\cmidrule(lr){2-3} \cmidrule(lr){4-5}
Metric & VerbConf & LP & VerbConf & LP \\
\midrule
$\rho$ & 0.297 & 0.205 & 0.337 & 0.316 \\
AUROC & 0.681 & 0.578 & 0.747 & 0.753 \\
Selection F1 & 0.620 & 0.647 & 0.889 & 0.904 \\
$\Delta$ vs.\ greedy & $-$2.45\,pp & $+$0.2\,pp & $-$1.93\,pp & $-$0.77\,pp \\
\midrule
Mean score & 8.58\,/\,10 & --- & 9.24\,/\,10 & --- \\
Parse fail rate & 0.26\% & --- & 0\% & --- \\
\bottomrule
\end{tabular}
}
\end{table}

\paragraph{Severe overconfidence.}
VerbConf exhibits extreme score compression: mean confidence is 8.58/10 on SciERC (single-sample F1${=}0.645$) and 9.24/10 on CoNLL (single-sample F1${=}0.904$), with most scores concentrated in the 8--10 range.
This compression is more severe than LP (which spans a 0.041-nat median range); score compression is a general property of model self-assessment, not an artifact of log-probability computation.

\paragraph{QE--selection decoupling.}
VerbConf achieves higher $\rho$ than LP on both datasets (SciERC: 0.297 vs.\ 0.205; CoNLL: 0.337 vs.\ 0.316) and competitive AUROC (SciERC: 0.681 vs.\ 0.578; CoNLL: 0.747 vs.\ 0.753), yet yields the worst selection among all signals ($\Delta{=}{-}2.45$~pp on SciERC, $-1.93$~pp on CoNLL).
This dissociation further demonstrates the correlation-selection gap: a signal can rank instances by expected quality without being able to discriminate candidates within an instance.
The AUROC--$\rho$ decoupling on CoNLL (LP leads AUROC at 0.753 while VerbConf leads $\rho$ at 0.337) reflects differing sensitivity to binary vs.\ ordinal discrimination.

\section{LLaMA-3.2-3B Scale Ablation}\label{app:llama3b}

To extend scale coverage below 4B, we fine-tune LLaMA-3.2-3B-Instruct on SciERC NER under identical training conditions (LoRA rank${=}32$, $\alpha{=}64$, 5 epochs) and evaluate quality estimation and selection (Table~\ref{tab:llama3b}).

\begin{table}[t]
\centering
\small
\caption{LLaMA-3.2-3B SciERC NER quality estimation and selection ($N{=}8$, $T{=}1.0$, seed${=}42$). Greedy F1${=}0.627$, Oracle F1${=}0.763$, headroom${=}13.63$~pp.}
\label{tab:llama3b}
\begin{tabular}{l cc cc}
\toprule
Signal & $\rho$ & AUROC & Sel.\ F1 & $\Delta$ vs.\ greedy \\
\midrule
VC & \textbf{.460} & .743 & 0.611 & $-$1.54\,pp \\
SJ & .375 & .738 & \textbf{0.633} & $+$0.61\,pp \\
EM & .334 & \textbf{.776} & --- & --- \\
FK & .268 & .671 & 0.630 & $+$0.28\,pp \\
LP & .266 & .633 & 0.628 & $+$0.12\,pp \\
\midrule
Greedy & --- & --- & 0.627 & --- \\
Oracle & --- & --- & 0.763 & --- \\
\bottomrule
\end{tabular}
\end{table}

The $\rho$ ranking at 3B (VC $>$ SJ $>$ EM $>$ FK $\approx$ LP) preserves the structural-over-LP ordering from 8B (Table~\ref{tab:quality_estimation}); ranking stability holds across two model families spanning 3B--8B.
AUROC rankings differ (EM $>$ VC $>$ SJ $>$ FK $>$ LP), consistent with the $\rho$--AUROC dissociation noted in the main text.
Despite 13.63~pp of oracle headroom, no signal achieves significant selection improvement (best: SJ, $+0.61$~pp, $p{=}0.175$); VC, the strongest QE signal, yields the worst selection ($-1.54$~pp), consistent with the correlation-selection gap observed at larger scales.

\section{Few-NERD Full Results}\label{app:fewnerd}

Table~\ref{tab:fewnerd_full} reports the complete Few-NERD INTRA coarse evaluation (Qwen3-8B, LoRA rank${=}32$, 3 epochs, $N{=}8$, $T{=}1.0$).
After gold-filtering (removing instances with no gold entities), the evaluation set contains $n{=}32{,}565$ instances across 8 coarse entity types; 11.9\% receive zero F1 under greedy decoding.
Instance degeneracy is 13.0\% (vs.\ 12.2\% on SciERC, 61.7\% on CoNLL).

\begin{table}[t]
\centering
\small
\caption{Few-NERD INTRA coarse: complete quality estimation and selection results (Qwen3-8B, $N{=}8$, $T{=}1.0$, gold-filtered $n{=}32{,}565$). $\rho$: Spearman correlation with instance F1. Sample AUROC: per-sample LP discrimination. Sel.\ $\Delta$: selection F1 minus greedy (0.748). CI: bootstrap 95\% confidence interval.}
\label{tab:fewnerd_full}
\resizebox{\columnwidth}{!}{
\begin{tabular}{lcccc}
\toprule
Signal & $\rho$ & Sample AUROC & Sel.\ $\Delta$ & 95\% CI \\
\midrule
LP & .220 & .578 & $\mathbf{+1.39}$\,pp & [+1.18, +1.53] \\
VC & .386 & --- & $-$1.79\,pp & --- \\
FK & .313 & --- & $-$2.26\,pp & --- \\
SJ & .300 & --- & $-$2.14\,pp & --- \\
EM & .287 & --- & $-$3.85\,pp & --- \\
\midrule
\multicolumn{3}{l}{Greedy: 0.748 \quad Oracle: 0.872} & \multicolumn{2}{l}{Headroom: 12.4\,pp} \\
\bottomrule
\end{tabular}
}
\end{table}

\paragraph{Key observations.}
(1)~All five signals achieve positive QE ($\rho$ 0.22--0.39); quality estimation works on Few-NERD. VC leads ($\rho{=}0.39$), followed by FK ($0.31$), SJ ($0.30$), EM ($0.29$), LP ($0.22$); the ranking matches SciERC where structural signals also outperform LP for quality estimation.
(2)~LP is the only signal achieving significant positive selection ($+1.39$~pp, $p{<}0.001$, CI excludes zero). The sample-level AUROC of 0.578 confirms LP provides per-sample discriminability beyond instance-level quality prediction.
(3)~Despite positive $\rho$, all structural signals degrade selection ($-1.79$ to $-3.85$~pp), further confirming that the correlation-selection gap is not an artifact of weak QE.
(4)~Low degeneracy (13.0\%) combined with strong LP-quality alignment ($\rho{\approx}0.5$--$0.7$) enables LP selection: candidates differ in F1 more often, and LP can rank them effectively. This contrasts with SciERC (12.2\% degeneracy but weak alignment $\rho{\approx}0.16$, LP $+0.61$~pp) and CoNLL (61.7\%, LP $-0.85$~pp).

\section{Per-Entity-Type LP Selection Analysis}\label{app:per_type_lp}

Table~\ref{tab:per_type_lp} stratifies LP selection by entity type on Few-NERD, sorted by degeneracy rate.
The degeneracy-gating mechanism holds at entity-type granularity: 7 of 8 types show positive LP selection gains, and the sole negative type (Art, $-0.17$~pp) has the highest degeneracy (19.8\%).
The correlation between per-type degeneracy and oracle headroom is strong and significant (Spearman $\rho{=}{-}0.881$, $p{=}0.004$; exact permutation $p{=}0.005$); degeneracy directly constrains the headroom available for selection.
Under a conservative $\text{df}{=}4$ adjustment for within-type dependence, the correlation remains significant ($p{=}0.033$).
Entity types show near-zero intra-class correlation (ICC $= -0.004$), supporting the independence assumption underlying per-type analysis.

\begin{table}[h]
\centering
\small
\caption{Per-entity-type LP selection on Few-NERD INTRA coarse (Qwen3-8B, $N{=}8$, $T{=}1.0$), sorted by degeneracy. LP $\Delta$: LP selection F1 minus greedy. 7/8 types show positive $\Delta$; Art (highest degeneracy) is the sole exception.}
\label{tab:per_type_lp}
\resizebox{\columnwidth}{!}{
\begin{tabular}{lrrrrrrr}
\toprule
Entity Type & $n$ & Degen.\% & Greedy & LP Sel & $\Delta$ (pp) & LP $\rho$ & Headroom \\
\midrule
Event & 1{,}306 & 6.0 & .655 & .662 & $+$0.65 & .457 & 16.1\,pp \\
Other & 3{,}433 & 6.6 & .637 & .685 & $\mathbf{+4.73}$ & .491 & 19.4\,pp \\
Organization & 6{,}647 & 6.7 & .720 & .732 & $+$1.21 & .509 & 13.3\,pp \\
Building & 1{,}587 & 7.2 & .702 & .705 & $+$0.23 & .478 & 14.9\,pp \\
Product & 2{,}276 & 8.2 & .656 & .677 & $+$2.01 & .483 & 15.8\,pp \\
Location & 8{,}023 & 11.4 & .793 & .798 & $+$0.53 & .520 & 10.3\,pp \\
Person & 8{,}016 & 13.4 & .822 & .835 & $+$1.25 & .544 & 9.2\,pp \\
Art & 1{,}277 & 19.8 & .756 & .754 & $-$0.17 & .451 & 9.2\,pp \\
\bottomrule
\end{tabular}
}
\end{table}

\begin{table}[h]
\centering
\small
\caption{CoNLL-2003 per-type LP selection (Qwen3-8B, $N{=}16$, $T{=}1.0$). All types exceed 20\% degeneracy and show negative LP selection, consistent with the dataset's ceiling-regime performance (F1 $>$ 90\%).}
\label{tab:per_type_lp_conll}
\begin{tabular}{lrrr}
\toprule
Entity Type & Degen.\% & LP $\Delta$ (pp) \\
\midrule
MISC & 23.2 & $-$2.19 \\
PER & 42.6 & $-$1.08 \\
LOC & 47.5 & $-$0.62 \\
ORG & 54.3 & $-$0.55 \\
\bottomrule
\end{tabular}
\end{table}

\begin{table}[h]
\centering
\small
\caption{SciERC per-type LP selection (Qwen3-8B, $N{=}8$, $T{=}1.0$). Within-dataset, degeneracy alone does not predict LP effectiveness (Spearman $\rho{=}0.14$, $p{=}0.79$); additional type-specific factors likely modulate selection at this granularity.}
\label{tab:per_type_lp_scierc}
\begin{tabular}{lrrr}
\toprule
Entity Type & Degen.\% & LP $\Delta$ (pp) \\
\midrule
Task & 10.1 & $+$2.58 \\
Generic & 10.0 & $+$2.07 \\
Material & 20.0 & $+$0.91 \\
Metric & 12.5 & $-$0.06 \\
Method & 16.3 & $-$0.03 \\
OtherScientificTerm & 6.2 & $-$0.57 \\
\bottomrule
\end{tabular}
\end{table}

The cross-dataset degeneracy gradient (61.7\% $\to$ 13.0\% $\to$ 12.2\%) is inversely associated with LP selection effectiveness at the macro level ($n{=}3$ datasets; exploratory).\footnote{Partial Spearman correlation controlling for entity-type frequency yields $\rho{=}{-}0.47$ (vs.\ raw $\rho{=}{-}0.38$); type frequency does not confound the degeneracy--LP relationship.}
At finer per-type granularity, selection depends on the interaction of degeneracy, oracle headroom, and type-specific LP-quality alignment.
A Kruskal-Wallis test finds no significant per-type heterogeneity in LP selection gains on SciERC ($p{=}0.796$), consistent with the small per-type sample sizes ($n{=}54$--$296$).
Low degeneracy alone does not ensure selection gains: CoNLL MISC (23\%, headroom 12.6~pp) and SciERC OtherScientificTerm (6.2\%, headroom 12.5~pp) both show negative LP selection despite available headroom; LP correlation strength with the type-specific quality function is an additional prerequisite.

\section{Calibration Baseline}\label{app:calibration_baseline}

Standard calibration methods are monotonic transforms of log-probability and therefore preserve within-instance sample rankings: Platt scaling \cite{platt1999probabilistic} fits a logistic regression $\sigma(a \cdot \text{LP} + b)$, and temperature scaling \cite{guo2017calibration} divides logits by a scalar $\tau$.
Since both preserve the argmax, selection F1 is identical to raw LP on all three datasets (Table~\ref{tab:calibration}).

To break monotonicity, we fit a multi-feature Platt model (MF-Platt) using LP, token count, and entity count as inputs, trained on a held-out calibration split (20\% of the test set, 5-fold cross-validation on the remainder).
MF-Platt can reorder candidates within an instance, yet yields no significant selection improvement: $\Delta{=}{-}1.1$~pp on SciERC ($p{=}0.72$), $+0.3$~pp on CoNLL ($p{=}0.31$), and $-0.6$~pp on Few-NERD ($p{=}0.83$).
Platt scaling substantially improves probability calibration (ECE drops from 0.26--0.39 to 0.01--0.08), but better-calibrated probabilities do not translate to better selection.
This parallels the supervised signal combination result (\S\ref{sec:robustness_main}): neither stronger aggregation nor additional features can compensate for weak per-sample discriminative signal.

\begin{table}[t]
\centering
\small
\caption{Calibration baseline results ($T{=}1.0$; single seed). SciERC uses $N{=}8$; CoNLL and Few-NERD use $N{=}16$ ($N$ does not affect greedy F1 or calibration conclusions). Platt and TempScale are monotonic transforms of LP that preserve the within-instance argmax (selection F1 $\equiv$ Raw LP, omitted). MF-Platt uses LP $+$ token count $+$ entity count; $p$: bootstrap test vs.\ Raw LP. Few-NERD evaluated on the first 5K test instances without gold-filtering (greedy .771; cf.\ gold-filtered full-set .748 in Table~\ref{tab:fewnerd}).}
\label{tab:calibration}
\resizebox{\columnwidth}{!}{
\begin{tabular}{l cccc cc}
\toprule
& & \multicolumn{2}{c}{Selection F1} & & \multicolumn{2}{c}{ECE} \\
\cmidrule(lr){3-4} \cmidrule(lr){6-7}
Dataset & Greedy & Raw LP & MF-Platt ($p$) & & Raw & Platt \\
\midrule
SciERC & .644 & .650 & .639 (.72) & & .260 & .048 \\
CoNLL & .909 & .898 & .901 (.31) & & .392 & .012 \\
Few-NERD & .771 & .785 & .778 (.83) & & .262 & .080 \\
\bottomrule
\end{tabular}
}
\end{table}

\section{Epoch Confound Analysis}\label{app:epoch_confound}

\paragraph{Training convergence modulates degeneracy.}
Table~\ref{tab:epoch_confound} compares the 3-epoch model (used throughout the main text) with a 5-epoch variant on Few-NERD (Qwen3-8B, $N{=}8$, $T{=}1.0$; gold-filtered $n{=}32{,}565$, 2-seed mean).
Extended training increases greedy F1 by $+0.7$~pp ($0.791 \to 0.798$) and triples instance degeneracy ($8.8$\% $\to$ $26.4$\%), reducing the fraction of instances with non-degenerate outputs.
The 5-epoch MV $\Delta$ ($-0.40$\,pp) appears smaller than 3-epoch ($-2.74$\,pp), but this is misleading: $26$\% of instances are fully degenerate (MV${\equiv}$greedy, $\Delta{=}0$), diluting the aggregate.
On non-degenerate instances, MV $\Delta{\approx}{-}0.55$\,pp, still negative---degeneracy masks the MV penalty rather than improving it.

LP selection shows even more dramatic degradation: 5-epoch LP $\Delta{=}{+}0.12{\pm}0.04$\,pp (2-seed mean), a 91\% reduction from the 3-epoch $+1.39$\,pp (both evaluated on the full test set under identical conditions; Table~\ref{tab:epoch_confound}). This confirms that training convergence not only increases degeneracy but also erodes LP's discriminative power within non-degenerate instances, consistent with the alignment degradation observed in SciERC (Table~\ref{tab:epoch_scierc}).

Evaluation loss increases after epoch 3, and the best checkpoint by validation loss occurs at approximately epoch 3.9.
Standard early stopping practice thus naturally selects the low-degeneracy regime.
This finding refines the narrative: degeneracy is not a fixed property of datasets but is modulated by the degree of training convergence. Over-trained models compress their output distributions, collapsing diversity even on datasets that support it under well-regularized training.

\begin{table}[t]
\centering
\small
\caption{Training convergence confound on Few-NERD (Qwen3-8B, $N{=}8$, $T{=}1.0$, gold-filtered $n{=}32{,}565$, 2-seed mean). ``Degen.'': entity-set exact-match degeneracy (all $N$ samples produce identical entity sets). ``MV $\Delta$'': strict majority vote gain over greedy (pp). ``Non-degen MV $\Delta$'': MV gain on non-degenerate instances only. LP $\Delta$ evaluated on the full test set ($n{=}37{,}648$, including gold-empty instances) for consistency with Table~\ref{tab:fewnerd}.}
\label{tab:epoch_confound}
\resizebox{\columnwidth}{!}{
\begin{tabular}{lcc}
\toprule
Metric & 3-epoch & 5-epoch \\
\midrule
Greedy F1 & 0.791 & 0.798 \\
MV $\Delta$ (pp) & $-$2.74 & $-$0.40 \\
Instance degen.\ (\%) & 8.8 & 26.4 \\
SJ mean & 0.685 & 0.802 \\
LP $\Delta$ (pp) & $+$1.39 & $+$0.12${\pm}$0.04 \\
Entity constr.\ $\Delta$ (pp) & $+$0.88 & $+$0.20 \\
\bottomrule
\end{tabular}
}
\end{table}

\paragraph{SciERC epoch curve.}
Table~\ref{tab:epoch_scierc} extends the analysis to SciERC with a 3-point curve (3, 5, 10 epochs; Qwen3-8B, $N{=}8$, $T{=}1.0$).
Degeneracy increases monotonically (6.4\% $\to$ 17.8\% $\to$ 47.8\%), and within-instance LP-quality correlation degrades in tandem (median $\rho$: 0.319 $\to$ 0.262 $\to$ 0.112; mean $\rho$: 0.273 $\to$ 0.179 $\to$ 0.066), confirming that overtraining erodes both diversity and LP alignment simultaneously.
The fraction of instances with positive LP-quality correlation also declines (76.4\% at 3 epochs $\to$ 64.1\% at 5 epochs $\to$ 54.0\% at 10 epochs).
LP selection does not yield significant gains at any training duration ($-1.24$, $+0.64$, $+1.02$~pp; all below the ${\approx}$1.5~pp MDE), despite oracle headroom ranging from 8.4 to 15.0~pp.
The 3-epoch model has higher within-instance alignment ($\rho{=}0.319$) than the 5-epoch model ($\rho{=}0.262$), yet LP selection is worse ($-1.24$ vs.\ $+0.64$~pp): at 6.4\% degeneracy, more non-degenerate instances enter the candidate pool, but LP ranking precision is too low to exploit them, so the broader pool hurts rather than helps. This illustrates that low degeneracy does not ensure selection gains; LP-quality alignment predicts selection direction.
At 10 epochs, the model retains competitive greedy F1 (0.649 vs.\ 0.650 at 3 epochs) yet conditions are unfavorable: 47.8\% degenerate instances and $\rho{=}0.112$.
The gradient confirms that degeneracy is convergence-driven rather than a fixed dataset property: extended training compresses output distributions, degrading both diversity and within-instance alignment. However, the epoch-to-epoch alignment variation ($\Delta\rho_{\text{median}}{=}0.057$) is small relative to the cross-dataset gap (SciERC ${\sim}$0.3 vs.\ Few-NERD ${\sim}$0.5), consistent with alignment being primarily dataset-intrinsic: training convergence modulates alignment at the margin but cannot shift a dataset into the strong-alignment regime.
Even at the most favorable epoch (3-epoch, $\rho{=}0.319$), SciERC alignment remains below the ${\sim}$0.5 threshold associated with successful LP selection on Few-NERD.

\begin{table}[t]
\centering
\small
\caption{SciERC NER epoch curve (Qwen3-8B, $N{=}8$, $T{=}1.0$). Training duration modulates both instance degeneracy and within-instance LP-quality alignment. ``Degen.'': constant-F1 instance degeneracy (gold-filtered). ``LP $\Delta$'': LP selection gain over greedy (gold-filtered; none reaches the ${\approx}$1.5~pp MDE). ``$\rho$ med'': median within-instance LP-quality correlation. Greedy F1 uses span-based evaluation consistent with Tables 1--3; $\rho$ and LP $\Delta$ use text-based matching consistent with Table~\ref{tab:within_rho}.}
\label{tab:epoch_scierc}
\resizebox{\columnwidth}{!}{
\begin{tabular}{lccc}
\toprule
Metric & 3-epoch & 5-epoch & 10-epoch \\
\midrule
Greedy F1 & 0.650 & 0.644 & 0.649 \\
Instance degen.\ (\%) & 6.4 & 17.8 & 47.8 \\
LP $\Delta$ (pp) & $-$1.24 & $+$0.64 & $+$1.02 \\
Within-inst.\ $\rho$ (median) & 0.319 & 0.262 & 0.112 \\
Within-inst.\ $\rho$ (mean) & 0.273 & 0.179 & 0.066 \\
\% positive $\rho$ & 76.4 & 64.1 & 54.0 \\
\bottomrule
\end{tabular}
}
\end{table}

\paragraph{Scope of epoch analysis.}
The Few-NERD epoch analysis covers 2 points (3 vs.\ 5 epochs); the SciERC curve provides a 3-point gradient showing monotonic degeneracy increase with simultaneous within-instance alignment degradation.
CoNLL epoch variation remains unverified; high-degeneracy datasets (CoNLL at 59--89\%) may exhibit a ceiling effect where additional training has limited marginal impact.
All Few-NERD results in the main text use the 3-epoch model (standard early stopping); the 5-epoch comparison is a controlled ablation.
All SciERC results use the 5-epoch model; the 3- and 10-epoch variants serve as controlled ablations.

\subsection{Degeneracy-Aware LP Selection}\label{app:dals}

The non-tied subset analysis (Table~\ref{tab:nontied_main}) identifies within-instance LP range (max minus min LP across $N$ samples) as a per-instance degeneracy proxy: instances with near-zero LP range carry no discriminative signal.
Degeneracy-aware LP selection (DALS) operationalizes this: when LP range exceeds $\tau{=}0.05$ (the threshold independently motivated by the non-tied analysis; Appendix~\ref{app:epsilon_robustness}), select the highest-LP sample; otherwise fall back to greedy.
On SciERC, DALS yields $+0.71{\pm}0.48$~pp (3-seed mean; individual seeds: $+0.19$ to $+1.14$~pp); on Few-NERD, $+0.30$~pp ($p{<}0.001$; gain remains positive across $\tau \in [0.01, 0.15]$); on CoNLL, where pure LP selection degrades F1 ($-0.85$~pp), DALS shows comparable degradation (Qwen: $-0.87$~pp; LLaMA: neutral).
DALS improves over pure LP selection when degeneracy is moderate, provides partial damage control when LP selection is harmful, and defaults to greedy when LP range signals no information.
While gains are modest and seed-sensitive, LP range serves as an operationally accessible degeneracy proxy, though the mixed cross-dataset results confirm that degeneracy gating alone cannot ensure positive selection outcomes.

\paragraph{LOO $\tau$ sensitivity.}
Leave-one-out (LOO) cross-validation across datasets shows that the optimal $\tau$ spans two orders of magnitude: $\tau^*{=}0.001$ on Few-NERD, $0.05$ on SciERC, and $0.10$ on CoNLL.
LOO regret (the gap between oracle and LOO-selected $\tau$) ranges from 0.51 to 1.21~pp.
The best single cross-dataset $\tau$ ($\tau{=}0.05$) matches the independently motivated non-tied threshold, but yields dataset-specific suboptimality: even with oracle $\tau$, CoNLL remains at $-0.31$~pp.
This sensitivity reinforces the diagnostic framing of DALS: per-dataset $\tau$ tuning is required for positive gains, and even optimally tuned DALS cannot overcome weak LP-quality alignment on high-degeneracy datasets.

\subsection{Conditional Decomposition of LP Selection}\label{app:conditional_decomp}

To test whether LP-quality alignment is associated with selection gains within datasets, we restrict LP selection to the favorable subset: non-degenerate instances whose within-instance LP range exceeds the dataset-specific median (SciERC: 0.040 nats, CoNLL: 0.062, Few-NERD: 0.067).
Table~\ref{tab:conditional_decomp} reports results across three datasets.

\begin{table}[h]
\centering
\small
\caption{LP selection on the favorable subset (non-degenerate instances with LP range above the dataset-specific median) per dataset (Qwen3-8B, $N{=}8$, $T{=}1.0$). CI: bootstrap 95\% confidence interval.}
\label{tab:conditional_decomp}
\resizebox{\columnwidth}{!}{
\begin{tabular}{lccc}
\toprule
Dataset & LP $\Delta$ (pp) & 95\% CI & Interpretation \\
\midrule
SciERC & $+$2.47 & [$+$0.30, $+$4.70] & significant positive \\
Few-NERD & $+$0.63 & [$-$0.02, $+$1.30] & marginal positive \\
CoNLL & $-$1.78 & [$-$2.83, $-$0.76] & significant negative \\
\bottomrule
\end{tabular}
}
\end{table}

On SciERC, restricting to the favorable subset recovers a significant LP gain, consistent with LP-quality alignment predicting selection within the non-degenerate pool.
On Few-NERD, the trend is directionally consistent but the confidence interval crosses zero; the full-dataset result ($+1.39$~pp, $p{<}0.001$) reflects the larger sample size rather than a stronger per-instance effect.
CoNLL remains negative even on the favorable subset, consistent with the ceiling-regime interpretation: at $>$90\% greedy F1, residual errors are idiosyncratic and LP cannot rank corrections.

\subsection{Power Analysis for Cross-Dataset Degeneracy--LP Correlation}\label{app:power_analysis}

Table~\ref{tab:power_analysis} reports post-hoc power for the $n{=}23$ cross-dataset correlation between degeneracy and LP selection gain ($\rho{=}0.052$, $p{=}0.81$; Figure~\ref{fig:degeneracy_scatter}).

\begin{table}[h]
\centering
\small
\caption{Post-hoc power analysis ($n{=}23$, two-tailed, $\alpha{=}0.05$). At this sample size, only large correlations ($\rho{\geq}0.5$) are detectable with adequate power.}
\label{tab:power_analysis}
\begin{tabular}{lc}
\toprule
True $\rho$ & Power at $n{=}23$ \\
\midrule
0.1 & 0.07 \\
0.2 & 0.13 \\
0.3 & 0.25 \\
0.4 & 0.44 \\
0.5 & 0.65 \\
0.6 & 0.84 \\
\bottomrule
\end{tabular}
\end{table}

\noindent 80\% power requires $n{\approx}85$ for $\rho{=}0.3$ and $n{\approx}47$ for $\rho{=}0.5$. The null result should be interpreted as absence of evidence for a large effect rather than evidence of absence.

\subsection{Entity-Token LP Decomposition}\label{app:entity_token_lp}

Full-sequence LP averages over all output tokens, diluting entity-relevant signal with formatting and schema tokens.
Entity-token LP restricts the average to tokens within entity text and type spans, yielding $4{\times}$ higher within-instance variance (median 0.000644 vs.\ 0.000161).
Table~\ref{tab:entity_token_lp} compares the two aggregation strategies on SciERC (Qwen3-8B, $N{=}8$, $T{=}1.0$).

\begin{table}[h]
\centering
\small
\caption{Full-sequence vs.\ entity-token LP on SciERC NER (Qwen3-8B, $N{=}8$, $T{=}1.0$). Entity-token LP has higher variance and stronger correlation, yet selection degrades.}
\label{tab:entity_token_lp}
\resizebox{\columnwidth}{!}{
\begin{tabular}{lcccc}
\toprule
Aggregation & Med.\ var. & $\rho$ & AUROC & $\Delta$ (pp) \\
\midrule
Full-sequence & 0.000161 & 0.297 & 0.642 & $+$0.35 \\
Entity text+type & 0.000644 & 0.440 & 0.719 & $-$0.73 \\
\bottomrule
\end{tabular}
}
\end{table}

Despite stronger correlation ($+$48\%) and higher AUROC, entity-token LP produces worse selection, reinforcing that compression operates at the representation level rather than token aggregation: the correlation-selection gap persists even with a more targeted signal.

\subsection{Entity-Level Majority Voting}\label{app:entity_merge}

Entity-level majority voting constructs a merged output by including all (text, type) entity tuples that appear in at least $k$ of $N$ samples (with $k = N/2$ as the default majority threshold).
Unlike best-of-$N$ selection (which picks one complete sample), merge combines evidence across samples at entity granularity.
We sweep $k \in \{2, 3, \ldots, 8\}$ on each dataset (Qwen3-8B, $N{=}8$, $T{=}1.0$) and report the best threshold per dataset.

\begin{table}[h]
\centering
\small
\caption{Entity-level majority voting results (Qwen3-8B, $N{=}8$, $T{=}1.0$). ``Greedy'': $T{=}0$ baseline. ``Best merge'': optimal $k$ per dataset from sweep $k \in \{2, \ldots, 8\}$. ``$\Delta$'': merge F1 minus greedy. ``Recovery'': $\Delta$ / (oracle $-$ greedy). Only SciERC shows positive recovery; even with oracle $k$, CoNLL and Few-NERD degrade.}
\label{tab:entity_merge}
\resizebox{\columnwidth}{!}{
\begin{tabular}{lcccc}
\toprule
Dataset & Greedy & Best merge ($k^*$) & $\Delta$ (pp) & Recovery (\%) \\
\midrule
SciERC & 0.644 & 0.660 ($k{=}3$) & $+$1.57 & $+$11.9 \\
CoNLL & 0.907 & 0.903 ($k{=}3$) & $-$0.38 & $-$9.0 \\
Few-NERD & 0.748 & 0.744 ($k{=}2$) & $-$0.43 & $-$3.6 \\
\bottomrule
\end{tabular}
}
\end{table}

On degenerate subsets (instances where all $N$ samples yield identical F1), merge is uniformly negative across all three datasets, confirming that entity-level aggregation cannot bypass the diversity bottleneck identified in the main text (\S\ref{sec:discussion}).

\subsection{Entity Construction Results}\label{app:entity_construction_table}

\begin{table}[h]
\centering
\small
\caption{Entity construction vs.\ greedy ($\Delta$ macro F1; $N{=}8$, $T{=}1.0$, Qwen3-8B, multi-seed mean${\pm}\sigma$). $\theta{=}2/N$. $p$: paired bootstrap vs.\ greedy. $^\dagger$Survives per-family Holm--Bonferroni correction (construction, $k{=}10$); global correction ($k{=}57$) yields $p_{\text{adj}}{=}0.055$ (Appendix~\ref{app:holm}). $^\ddagger$LP-weighted not computed.}
\label{tab:entity_construction}
\resizebox{\columnwidth}{!}{%
\begin{tabular}{@{}llcccc@{}}
\toprule
Dataset & Wt. & Greedy & Constr. & $\Delta$ & $p$ \\
\midrule
SciERC & LP & $.667{\pm}.001$ & $.670{\pm}.003$ & $+0.32$\,pp & n.s. \\
       & Unif. & $.667{\pm}.001$ & $.660{\pm}.004$ & $-0.74$\,pp & n.s. \\
\addlinespace
Few-NERD & LP & $.748$ & $\mathbf{.762{\pm}.005}$ & $+1.40{\pm}0.51$\,pp & ${<}.001^\dagger$ \\
         & Unif. & $.748$ & $\mathbf{.756{\pm}.006}$ & $+0.74{\pm}0.58$\,pp & .003 \\
\addlinespace
CoNLL$^\ddagger$ & Unif. & .908 & .904 & $-0.39$\,pp & n.s. \\
\bottomrule
\end{tabular}}
\end{table}

\subsection{Entity Construction Threshold Sensitivity}
\label{app:theta_sweep}

Table~\ref{tab:theta_sweep} reports entity construction F1 as a function of the confirmation threshold $\theta$ across three datasets (Qwen3-8B, $N{=}8$, $T{=}1.0$). The default $\theta{=}2/N{=}0.25$ achieves the best LP-weighted $\Delta$F1 on Few-NERD and is within 0.04\,pp of the optimum on SciERC. Performance is stable across $\theta \in [0.15, 0.35]$; below 0.15, spurious entities degrade precision, and above 0.40, aggressive filtering removes correct low-frequency entities.

\begin{table}[h]
\caption{Entity construction threshold sensitivity ($N{=}8$, Qwen3-8B, $T{=}1.0$). $\Delta$: F1 difference vs.\ greedy in percentage points. Best $\Delta$ per dataset--weighting pair in \textbf{bold}.}
\label{tab:theta_sweep}
\centering
\small
\resizebox{\columnwidth}{!}{
\begin{tabular}{ll rr rr}
\toprule
Dataset & $\theta$ & LP-w F1 & Unif.\ F1 & $\Delta$ LP-w & $\Delta$ Unif.\ \\
\midrule
\multirow{8}{*}{SciERC}
 & 0.10 & .620 & .620 & $-$2.37 & $-$2.37 \\
 & 0.15 & .659 & .659 & $+$1.54 & $+$1.54 \\
 & 0.20 & .659 & .659 & $+$1.54 & $+$1.54 \\
 & 0.25 & .661 & .659 & $+$1.67 & $+$1.54 \\
 & 0.30 & .661 & .661 & \textbf{$+$1.71} & \textbf{$+$1.71} \\
 & 0.35 & .661 & .661 & \textbf{$+$1.71} & \textbf{$+$1.71} \\
 & 0.40 & .653 & .654 & $+$0.94 & $+$0.97 \\
 & 0.50 & .649 & .654 & $+$0.46 & $+$0.97 \\
\midrule
\multirow{8}{*}{Few-NERD}
 & 0.10 & .677 & .676 & $-$7.12 & $-$7.15 \\
 & 0.15 & .751 & .751 & $+$0.32 & $+$0.32 \\
 & 0.20 & .751 & .751 & $+$0.32 & $+$0.32 \\
 & 0.25 & .757 & .751 & \textbf{$+$0.88} & $+$0.32 \\
 & 0.30 & .750 & .750 & $+$0.21 & $+$0.21 \\
 & 0.35 & .750 & .750 & $+$0.21 & $+$0.21 \\
 & 0.40 & .729 & .728 & $-$1.90 & $-$2.01 \\
 & 0.50 & .725 & .728 & $-$2.25 & $-$2.01 \\
\midrule
\multirow{8}{*}{CoNLL}
 & 0.10 & .844 & .842 & $-$6.44 & $-$6.56 \\
 & 0.15 & .896 & .896 & $-$1.19 & $-$1.19 \\
 & 0.20 & .896 & .896 & $-$1.19 & $-$1.19 \\
 & 0.25 & .901 & .896 & $-$0.71 & $-$1.19 \\
 & 0.30 & .904 & .904 & \textbf{$-$0.39} & \textbf{$-$0.39} \\
 & 0.35 & .904 & .904 & $-$0.41 & \textbf{$-$0.39} \\
 & 0.40 & .900 & .899 & $-$0.83 & $-$0.94 \\
 & 0.50 & .897 & .899 & $-$1.12 & $-$0.94 \\
\bottomrule
\end{tabular}
}
\end{table}

\subsection{Few-Shot Error Correlation Analysis}\label{app:fs_error}

We investigated whether few-shot demonstrations induce correlated entity-level errors that explain construction failure.
Pairwise error overlap (Jaccard similarity of per-sample false-positive and false-negative entity sets) shows no significant difference between FS and ZS regimes (SciERC: $0.454$ vs.\ $0.452$, $p{=}0.84$; CoNLL: $0.103$ vs.\ $0.104$, $p{=}0.64$; permutation tests, 10{,}000 resamples).
We additionally measure entity-level prediction entropy across $N{=}8$ samples per instance.
On SciERC, FS and ZS entropies are indistinguishable ($0.483$ vs.\ $0.485$, $p{=}0.81$).
On CoNLL, however, entity prediction entropy is significantly \emph{higher} under FS ($0.209$ vs.\ $0.184$, $p{<}10^{-9}$, $r{=}0.106$, small effect; Wilcoxon signed-rank test), contrary to the hypothesis that demonstrations concentrate predictions---they instead \emph{increase} output diversity on this dataset.
The mechanism by which few-shot demonstrations undermine construction remains an open question; entity-level error overlap is preserved while prediction entropy may even rise, suggesting that the \emph{nature} rather than the \emph{diversity} of errors shifts in ways that defeat LP-weighted voting.

\subsection{Construction Weighting Variants}\label{app:weighting_variants}

Table~\ref{tab:weighting_variants} compares LP-weighted, VC-weighted, and SJ-weighted entity construction on LLaMA-3.1-8B (3-seed mean, $N{=}8$, $T{=}1.0$, $\theta{=}2/N$).
All three weighting schemes beat greedy on SciERC, where degeneracy is low; LP and VC also improve Few-NERD, though gains are smaller than the Qwen results in Table~\ref{tab:entity_construction}.
The optimal weighting scheme is model-dependent (VC leads on LLaMA SciERC; LP leads on Qwen), but the core finding is robust: entity-level voting improves F1 in low-degeneracy regimes regardless of the specific weighting function.
SJ-weighted construction degrades on Few-NERD ($-0.20$\,pp), consistent with SJ over-weighting high-consensus samples, which correlates with the degeneracy direction rather than opposing it.

\begin{table}[h]
\centering
\small
\caption{Construction weighting variants (LLaMA-3.1-8B, $N{=}8$, $T{=}1.0$, $\theta{=}2/N$, 3-seed mean). $\Delta$: F1 difference vs.\ greedy in pp. Main-text construction results (Table~\ref{tab:entity_construction}) use Qwen3-8B.}
\label{tab:weighting_variants}
\begin{tabular}{lccc}
\toprule
Dataset & LP-wt $\Delta$ & VC-wt $\Delta$ & SJ-wt $\Delta$ \\
\midrule
SciERC & $+$0.91 & $+$1.05 & $+$0.72 \\
Few-NERD & $+$0.02 & $+$0.15 & $-$0.20 \\
\bottomrule
\end{tabular}
\end{table}

\subsection{Zero-Shot vs.\ Few-Shot Comparison}\label{app:zs_vs_fs}

Table~\ref{tab:72b_zs_vs_fs} presents the full zero-shot (ZS) vs.\ few-shot (FS) comparison on Qwen2.5-72B-Instruct-AWQ.
LP-weighted construction yields a significant gain only on SciERC under ZS ($+$1.53\,pp, $p{=}.008$); all five FS configurations show non-significant deltas ($|{\Delta}| \leq 0.08$\,pp).
Few-shot demonstrations raise the greedy baseline (e.g., SciERC: .399$\to$.461) but do not translate into construction improvements, suggesting that the dual-gating mechanism---where both signal quality and diversity are necessary---holds within the FT+ZS regime while FS introduces an independent failure mode whose origin remains open (\S\ref{app:fs_error}).

\begin{table*}[t]
\centering
\caption{Full zero-shot vs.\ few-shot comparison for Qwen2.5-72B-Instruct-AWQ ($N{=}8$). Only SciERC under zero-shot yields a significant construction gain.}
\label{tab:72b_zs_vs_fs}
\resizebox{\textwidth}{!}{%
\begin{tabular}{ll rrrrr rr r}
\toprule
Dataset & Regime & Degen\% & LP $\rho$ & SJ $\rho$ & FK $\rho$ & VC $\rho$ & Greedy F1 & Constr $\Delta$ & $p$ \\
\midrule
SciERC  & ZS & 21.0 & .081  & .217 & .227 & .277 & .399 & \textbf{+1.53} & .008 \\
SciERC  & FS & 16.6 & .060  & .293 & .301 & .258 & .461 & $-$0.08 & .544 \\
\midrule
CoNLL   & ZS & 67.7 & $-$.002 & .304 & .302 & .297 & .712 & $-$0.37 & --- \\
CoNLL   & FS & 70.0 & .106  & .442 & .451 & .426 & .745 & +0.03 & .417 \\
\midrule
Few-NERD & ZS & 46.2 & .119  & .345 & .321 & .416 & .628 & $-$0.20 & --- \\
Few-NERD & FS & 45.2 & .161  & .345 & .327 & .410 & .619 & +0.02 & .401 \\
\bottomrule
\end{tabular}%
}
\end{table*}

\section{Evaluation Protocol for Structured Self-Consistency}\label{app:eval_protocol}

We distill our experimental methodology into a reusable three-step protocol for assessing test-time compute scaling on set-valued structured prediction tasks.

\paragraph{Step 1: Measure instance degeneracy.}
Given $N$ stochastic samples per instance, compute the fraction of instances where all samples produce identical task-metric scores (\S\ref{sec:evaluation}).
Degeneracy ${>}30$\% indicates limited headroom for any selection or aggregation strategy.
Report degeneracy alongside all test-time scaling results to contextualize gains.

\paragraph{Step 2: Assess signal--quality alignment.}
For each candidate signal, compute (a)~cross-instance Spearman $\rho$ with the task metric (quality estimation strength) and (b)~within-instance $\rho$ between signal and task metric, restricted to non-degenerate instances (selection discriminability).
A large gap between (a) and (b) indicates the correlation-selection gap: reliable QE but poor selection prospects.

\paragraph{Step 3: Evaluate aggregation strategies.}
Compare three regimes against greedy ($T{=}0$) and oracle baselines:
\begin{itemize}
    \item \textit{Instance-level selection}: best-of-$N$ via per-sample signals (e.g., LP) and consensus-proximity (MBR utility).
    \item \textit{Entity-level construction}: majority-voted entity aggregation with threshold $\theta$ (e.g., $\theta{=}2/N$), optionally weighted by signal confidence.
    \item \textit{Selective prediction}: abstain on low-confidence instances; report coverage--F1 tradeoff.
\end{itemize}

\paragraph{Reporting standards.}
Report: (1)~degeneracy percentage, (2)~oracle headroom (oracle $-$ greedy), (3)~signal $\rho$ (cross-instance and within-instance), (4)~selection $\Delta$F1 with bootstrap $p$-values and effect sizes (Cohen's $d$), (5)~construction $\Delta$F1, (6)~selective prediction AUPRC and coverage--F1 curve.
Use paired bootstrap tests ($B{\geq}10{,}000$); report minimum detectable effect at 80\% power given the test set size.

\section{Cost-Benefit Analysis}\label{app:cost_benefit}

Table~\ref{tab:cost_benefit} summarizes the inference cost and F1 gain of each aggregation strategy relative to greedy decoding.
Entity construction achieves the only consistently positive result at zero marginal cost (reusing the $N$ samples already drawn for quality estimation); cross-model ensemble requires $2N{\times}$ inference for a smaller gain.

\begin{table}[h]
\centering
\small
\caption{Cost-benefit summary of aggregation strategies (Qwen3-8B, $N{=}8$, $T{=}1.0$). Inference cost is relative to greedy ($T{=}0$, single forward pass). $\Delta$F1 reports the primary result (dataset in parentheses).}
\label{tab:cost_benefit}
\begin{tabular}{@{}lcll@{}}
\toprule
Method & Cost & $\Delta$F1 & Verdict \\
\midrule
Greedy ($T{=}0$) & $1{\times}$ & baseline & default \\
SC + MV & $N{\times}$ & ${\approx}0$\,pp & no gain \\
SC + LP selection & $N{\times}$ & $+1.39$\,pp (FN) & dataset-dep. \\
SC + construction & $N{\times}$ & $+0.32$\textendash$1.40$\,pp & \textbf{positive} \\
Cross-model SC & $2N{\times}$ & $+0.42$\,pp (FN) & positive, $2{\times}$ cost \\
SC + DALS & $N{\times}$ & $+0.71$\,pp (SciERC) & seed-sensitive \\
SC + verifier & $N{\times}{+}\text{train}$ & $+0.30$\,pp & no transfer \\
\bottomrule
\end{tabular}
\end{table}

SC + LP selection is positive only on Few-NERD ($p{<}0.001$) and neutral to negative elsewhere (SciERC $+0.2$\,pp, $p{=}0.29$; CoNLL $-0.85$\,pp).
Entity construction is the only strategy that is positive on multiple datasets without additional model access, achieving zero marginal cost by reusing the $N$ samples already generated.
Cross-model ensemble breaks within-model degeneracy but doubles inference cost.

\section{Prior Work Differentiation Matrix}\label{app:differentiation}

Table~\ref{tab:differentiation} positions this work relative to four closely related studies.
All four investigate settings where self-consistency underperforms; we are the first to provide a multi-stage diagnostic with mechanistic probes, theory, and exhaustive method evaluation on structured generation.

\begin{table*}[h]
\centering
\small
\caption{Differentiation matrix: positioning relative to prior work on self-consistency limitations. ``Methods tested'' counts the number of distinct aggregation/selection categories evaluated (not hyperparameter variants).}
\label{tab:differentiation}
\resizebox{\textwidth}{!}{%
\begin{tabular}{@{}l p{2.8cm} p{2.5cm} p{2.5cm} p{2.8cm} p{3.2cm}@{}}
\toprule
Aspect & \citet{yun2025price} & \citet{chen2023twofailures} & \citet{aggarwal2023adaptive} & \citet{farquhar2024detecting} & \textbf{Ours} \\
\midrule
Root cause & Constrained decoding suppresses diversity & Reasoning failures (hypothetical, compositional) & Sample efficiency in allocation & Hallucination via semantic inconsistency & Capability ceiling (degeneracy $+$ LP compression $+$ verification failure) \\
\addlinespace
Tasks & Math, QA, structured gen. & Reasoning, QA & Math, code & QA, open-ended & NER ($+$SQL, RE) \\
\addlinespace
Methods tested & 1 (diversity analysis) & 1 (failure taxonomy) & 1 (adaptive $N$) & 1 (semantic entropy) & 9 categories (Table~\ref{tab:method_summary}) \\
\addlinespace
Theoretical contrib. & None & None & Allocation bounds & Entropy-based & MBR equivalence, degeneracy bound, scaling laws \\
\addlinespace
Mechanistic probe & No & No & No & No & Yes (ESJ, MRSC, LSC, hidden-state) \\
\bottomrule
\end{tabular}}
\end{table*}

\section{Extended Discussion}\label{app:objections}

\paragraph{Open questions.}
Three directions remain: (1)~whether learned verifiers (e.g., reward models) can break the introspection boundary and exploit hidden-state quality signals that prompt-based methods cannot access; (2)~whether diversity-promoting decoding strategies can preserve candidate variety without sacrificing output quality, mitigating the capability-diversity tension; (3)~whether the diagnostic framework transfers to other structured prediction tasks (e.g., coreference resolution, syntactic parsing) where output spaces similarly resist equivalence-class aggregation.

We address twelve anticipated objections, each referencing specific evidence in the paper.

\paragraph{O1: Results are NER-specific and may not generalize.}
We validate the diagnostic framework on Spider SQL parsing (GPT-5.5, $n{=}196$; \S\ref{sec:introduction}): degeneracy reaches 87.2\%, confirming the failure mode beyond NER. An RE extension on SciERC (Appendix~\ref{app:re_extension}) shows compound error propagation worsens all signals. The structural mismatch (set-valued outputs without equivalence classes) applies to any structured generation task.

\paragraph{O2: Cross-model gain of $+$0.42\,pp is marginal.}
$+0.42$\,pp is the only unconditionally positive result among 9 method categories tested (Table~\ref{tab:method_summary}). Its significance ($p{<}0.001$, CI $[0.27, 0.58]$; \S\ref{sec:cross_model}) and theoretical motivation (architecturally distinct models produce independent error distributions, breaking within-model degeneracy) distinguish it from the null results of other categories.

\paragraph{O3: Hyperparameters may not be optimal.}
We report full sweeps: construction threshold $\theta \in [0.05, 0.50]$ (Appendix~\ref{app:theta_sweep}), DALS threshold $\tau$ with LOO cross-validation (regret 0.51--1.21\,pp; Appendix~\ref{app:dals}), LP compression threshold $\varepsilon \in [0, 0.2]$ (Appendix~\ref{app:epsilon_robustness}), and temperature $T \in \{0.5, 0.7, 1.0, 1.2\}$ (Appendix~\ref{app:t_ablation}).

\paragraph{O4: Findings may be a fine-tuning artifact.}
Qwen2.5-72B zero-shot exhibits the same gap: LP $\rho$ drops without fine-tuning, structural signals generalize (VC $\rho{=}0.416$ on Few-NERD), and degeneracy is model-family-dependent rather than scale-dependent (Qwen2.5: 31.7--46.2\% vs.\ Qwen3: 13.0\%; Appendix~\ref{app:cross_scale}). The capability ceiling persists from 4B to 72B.

\paragraph{O5: Probe AUROC varies widely (0.63--0.88); why trust the result?}
The variation is complexity-dependent and interpretable: CoNLL (4 types, 91\% F1) at 0.882, Few-NERD (8 types, 75\% F1) at 0.636, SciERC (6 types, 65\% F1) at 0.630. LODO transfer drops to 0.56--0.67 (\S\ref{sec:mechanistic}), confirming dataset-specific encoding. The pattern---not the point estimate---is the finding.

\paragraph{O6: Theory is disconnected from experiments.}
Each formal result has explicit empirical validation: Proposition~\ref{prop:construction_mbr} (MBR equivalence) $\to$ Table~\ref{tab:entity_construction}; Proposition~\ref{prop:entity_selection_bound} (entity selection bound) $\to$ Table~\ref{tab:alpha_crit} across 3 datasets; Degeneracy Decomposition $\to$ Figure~\ref{fig:bound} across 16 configurations (\S\ref{sec:theory}); oracle scaling law $\to$ Figure~\ref{fig:oracle_scaling} ($R^2 > 0.999$; \S\ref{sec:scaling}).

\paragraph{O7: Why not use standard MBR decoding?}
Majority vote \textit{is} MBR decoding under set-overlap utility (Proposition~\ref{prop:construction_mbr}; \S\ref{sec:mbr_theory}). LP-weighted construction is the quality-aware MBR variant (\S\ref{sec:mbr_theory}). This follows directly from \citet{eikema2022sampling}; \citet{eikema2025structured} confirm that MBR utilities degrade on structured outputs, consistent with our LP compression findings.

\paragraph{O8: No qualitative error analysis.}
Appendix~\ref{app:failure_gallery} provides a failure mode gallery with 9 concrete instances: 3 degenerate (all samples identical), 3 LP selection failures (high-LP sample worse than low-LP), and 3 self-verification failures (ESJ misclassifies entity correctness). Each includes instance text, gold annotations, predictions, and signal values.

\paragraph{O9: Evaluation metric choice may bias results.}
We report per-instance F1 averaged across instances (macro F1) as the primary selection metric. Aggregate micro-F1 (pooled TP/FP/FN) confirms MV fails to improve over greedy: $\Delta$ ranges from $-2.74$ to ${\approx}0$\,pp (Table~\ref{tab:alpha_crit}; Appendix~\ref{app:alpha_crit}). The correlation-selection gap holds under both aggregations; the selection failure is not an artifact of the metric.

\paragraph{O10: Computational cost analysis is missing.}
Table~\ref{tab:cost_benefit} (Appendix~\ref{app:cost_benefit}) provides a cost-benefit breakdown. Entity construction achieves the best cost-efficiency: $+0.32$--$1.40$\,pp at $N{\times}$ cost with zero marginal overhead beyond sampling. Cross-model ensemble doubles inference cost for $+0.42$\,pp.

\paragraph{O11: Reproducibility concerns.}
All experiments use open-weight models (Qwen3-8B, LLaMA-3.1-8B, Qwen3-4B, Qwen2.5-7B, Qwen2.5-72B) and public benchmarks. Multi-seed stability ($\sigma{\leq}0.026$ for QE; Table~\ref{tab:multiseed}) and full hyperparameter sweeps are reported. Code, checkpoints, and per-seed results will be released upon acceptance.

\paragraph{O12: Overlap with \citet{yun2025price}.}
\S\ref{sec:yun_contrast} provides a three-way differentiation. Our LP compression (median 0.041 nats) corroborates Yun et al.'s diversity collapse, but we diverge in root cause (model-internal vs.\ format-induced), diagnostic depth (diagnostic framework with 9 method categories vs.\ diversity analysis), and scale ($4{\times}$: Qwen2.5-72B). Table~\ref{tab:differentiation} (Appendix~\ref{app:differentiation}) provides the full comparison matrix.

\section{Failure Mode Gallery}\label{app:failure_gallery}

We present 9 concrete instances illustrating the three primary failure modes: output degeneracy, LP confidence inversion, and consensus-based verification failure. All examples use Qwen3-8B on Few-NERD ($N{=}8$, $T{=}1.0$) from the MF4v2 pipeline.

\paragraph{Degenerate instances (3/3).}
All $N$ samples produce identical entity sets, eliminating any room for selection or construction improvement.

\begin{itemize}[nosep,leftmargin=*]
\item \textbf{Type confusion.} \textit{``His undergraduate studies at the University of Notre Dame were interrupted by Army service in World War II.''} Gold: \{University of Notre Dame / building\}. All 8 samples: \{University of Notre Dame / \textbf{organization}, World War II / event\}. F1${=}0.00$: correct span, wrong type, plus hallucinated entity.

\item \textbf{Type confusion (minimal).} \textit{``...the deployment of the `SPQ' template...''} Gold: \{SPQ / organization\}. All 8 samples: \{SPQ / \textbf{product}\}. F1${=}0.00$: entity text matches exactly but type disagrees across all samples.

\item \textbf{Complete miss.} \textit{``On stage, the band uses different kinds of instruments like synths, guitars, cello, glockenspiel and a bass played with a sickle.''} Gold: 6 entities (synths, guitars, cello, glockenspiel, bass, sickle / product). All 8 samples: $\emptyset$. F1${=}0.00$: total recall failure on all gold entities.
\end{itemize}

\paragraph{LP confidence failures (3/3).}
High log-probability does not guarantee correctness at the entity level. These entities appear in all 8 samples with near-zero negative log-probability (high confidence), yet are absent from the gold annotation.

\begin{itemize}[nosep,leftmargin=*]
\item \textit{``Cramer pitched for the Orange County Flyers of the independent Golden League in 2008.''} Entity: Cramer / person (mean LP${=}{-}0.046$, 8/8 samples). Gold: \{Orange County Flyers / org, independent Golden League / org\}. Plausible entity, not annotated.

\item \textit{``...according to UK Royal Air Force Wing Commander and FOD researcher Richard Friend.''} Entity: Richard Friend / person (mean LP${=}{-}0.050$, 8/8). Gold: \{UK Royal Air Force / org, FOD / org\}. Likely annotation gap: Richard Friend is a real person mentioned by name; gold omits it.

\item \textit{``...the Navy's forces within South Vietnam were operationally controlled by COMUSMACV.''} Entity: COMUSMACV / org (mean LP${=}{-}0.050$, 8/8). Gold: \{Seventh Fleet / org, South Vietnam / location\}. Model confidently extracts a valid entity missing from gold.
\end{itemize}

\noindent These examples highlight a fundamental limitation: LP signals the model's confidence, not correctness. When the model is correctly confident about entities that gold annotations omit, LP cannot serve as an error detector.

\paragraph{Consensus-based verification failures (3/3).}
Entity consensus frequency (fraction of $N$ samples containing a given entity) serves as the primary self-verification signal. It fails in both directions.

\begin{itemize}[nosep,leftmargin=*]
\item \textbf{Overconfident.} \textit{``Les Blank...was an American documentary filmmaker best known for his portraits of American traditional musicians.''} Entity: American / location (8/8 samples, consensus${=}1.0$). Gold: \{Les Blank / person\}. Maximum consensus for non-gold entity.

\item \textbf{Overconfident.} \textit{``Cramer pitched for the Orange County Flyers of the independent Golden League in 2008.''} Entity: Cramer / person (8/8 samples, consensus${=}1.0$). Gold: \{Orange County Flyers / org, independent Golden League / org\}. Same instance as LP failure above---simultaneous failure of both signals.

\item \textbf{Underconfident.} \textit{``They finished the season 14--19, 9--9 in C-USA play to finish in seventh place.''} Gold entity: C-USA play / event (consensus${=}0.0$, never predicted in any of 8 samples). Zero consensus for a correct entity.
\end{itemize}

\noindent These examples illustrate that the failure modes identified by the diagnostic framework are not statistical abstractions but manifest in individual instances. The LP and consensus failures frequently co-occur (e.g., ``Cramer'' instance), confirming that the signals share common blind spots rather than providing independent error detection.

\section{Practical Recommendations}\label{app:recommendations}

Measuring degeneracy and LP-F1 alignment at $N{=}8$ suffices: if ${>}50$\% degenerate, scaling provides no benefit; if degeneracy ${<}20$\% with $\rho{>}0.70$, construction may yield gains; otherwise prefer greedy.

\section{Use of AI Assistants}
\label{app:ai-assistants}

We used Claude (Anthropic) as an AI coding and writing assistant
throughout this work. Specifically:
\begin{itemize}
    \item \textbf{Code:} drafting experiment scripts, data-processing
    pipelines, and analysis utilities; debugging and refactoring.
    \item \textbf{Writing:} drafting and revising paper sections,
    polishing prose, and reformatting LaTeX.
    \item \textbf{Literature:} summarizing related papers to inform
    related-work section.
\end{itemize}
